\documentclass[10pt]{article}

\usepackage{tabu}
\usepackage{longtable}
% Cross-reference the parent report's equation labels. This requires
% rolling-contact-qp.aux to exist, so build rolling-contact-qp.tex first.
\usepackage{xr}
\externaldocument{rolling-contact-qp}
%%%%%%%%%%%%%%%%%%%%%%%%%%%%%%%%%%%%% FORMAT
%%%%%%%%%%%%%%%%%%%%%%%%%%% footer and header
\usepackage{lastpage}

\usepackage{fancyhdr}
\usepackage{xparse}
\pagestyle{fancy}
%%%%%%%%%%%%%%%%%%%%%%%%%%% Headers and footers

\fancyhead{} % clear all header fields
\fancyfoot{} % clear all footer fields


\rfoot{Page \thepage\ of \pageref{LastPage}}
% Define the footer separation line
\renewcommand{\footrulewidth}{0.4pt}
\addtolength{\headwidth}{4.8cm}
\usepackage[headheight=80pt,tmargin=85pt,headsep=25pt, footskip=25pt]{geometry}

\geometry{a4paper, margin=2cm}

\usepackage[numbers]{natbib}

%%%%%%%%%%%%%%%%%%%%%%%%%%% Modify the description list
\usepackage{nameref}

\makeatletter
\let\orgdescriptionlabel\descriptionlabel
\renewcommand*{\descriptionlabel}[1]{%
  \let\orglabel\label
  \let\label\@gobble
  \phantomsection
  \edef\@currentlabel{#1\unskip}%
  \let\label\orglabel
  \orgdescriptionlabel{#1}%
}
\makeatother
%%%%%%%%%%%%%%%%%%%%%%%%%%% Font colour

\usepackage[utf8]{inputenc}
\usepackage[english]{babel}
\usepackage[dvipsnames]{xcolor}

\usepackage{listings}
\usepackage{color}
% %%%%%%%%%%%%%%%%%%%%%%%%% Tikz packages:
\usepackage{tikz}
\usepackage{tikz-layers}
\usepackage{array}
\newcolumntype{C}[1]{>{\centering\arraybackslash}p{#1}}
%
% the positioning-plus library (is on TeX.sx),
% the node-families library (is on TeX.sx),
% the fit library (loaded by positioning-plus),
% the backgrounds library to draw stuff behind other stuff,
% the calc library for some funny coordinate calculations and the
% shapes.geometric library for the ellipse shape

\usetikzlibrary{shapes.geometric,automata,positioning,arrows,shadows,backgrounds,calc}

\definecolor{frameColor}{RGB}{128,128,128}
\definecolor{initialStateColor}{RGB}{10,100,10}
\definecolor{impactStateColor}{RGB}{0,255,0}
\definecolor{admittanceStateColor}{RGB}{0,0,255}

\definecolor{resetStateColor}{RGB}{10,100,10}

\definecolor{stateColor}{RGB}{51,204,204}

\definecolor{controllerColor}{RGB}{10,10,255}
\definecolor{plannerColor}{RGB}{10,150,150}
\definecolor{estimatorColor}{RGB}{10,200,10}
\definecolor{robotColor}{RGB}{255,0,0}
\definecolor{modelColor}{RGB}{180,100,10}

%%% Enumeration
\usepackage{enumitem}

%%%%%%%%%%%%%%%%%%%%%%%%%%% Table
\usepackage{multirow}


%%%%%%%%%%%%%%%%%%%%%%%%%%% Adjust font size
\usepackage{setspace}

%%%%%%%%%%%%%%%%%%%%%%%%%%%
%%%%%%%%%%%%%%%%%%%%%%%%%%%%%%%%%%%%%% MATH:
%%%%%%%%%%%%%%%%%%%%%%%%% AMS packages:
\usepackage{mathtools}
\usepackage{amsmath, amsthm, amsfonts}
\usepackage{amssymb}
\usepackage{graphicx}
\usepackage[font=normal,labelfont=bf]{caption}
\usepackage{subcaption}
\usepackage{url}
\usepackage[colorlinks=true,hyperindex=true,bookmarksnumbered]{hyperref} % Turn on "bookmarks" and "hyperindex" to generate outlines.

%%%%%%%%%%%%%%%%%%%%%%%%% Answer and reply
\usepackage{thmtools}
\declaretheoremstyle[headfont=\bfseries,
bodyfont=\normalfont,
numbered=no]{normalhead}
\declaretheorem[style=normalhead]{Reply}

\theoremstyle{plain}
\newtheorem*{Reply*}{Reply}

\newtheorem*{AuthorReply*}{AuthorReply}

\declaretheoremstyle[
  headfont=\color{blue}\normalfont\bfseries,
  bodyfont=\color{red}\normalfont\itshape,
  numbered=no
]{quoteColor}

\declaretheorem[style=quoteColor]{Quote}
\newtheorem*{Quote*}{Quote}

\declaretheoremstyle[
  headfont=\color{black}\normalfont\bfseries,
  bodyfont=\color{blue}\normalfont,
  numbered=no
]{reviewColor}
\declaretheorem[style=reviewColor]{review}
\newtheorem*{review*}{review}

\theoremstyle{definition}

\theoremstyle{remark}
% -----------------------------------------------------------------


%%%%%%%%%%%%%%%%%%%%%%%%%%%%%%%%%%%%%%%%%%%%%%%%%%%%%%%%%%%%%%%%%%%%%%%%%%%%%%%%
%% 1. PACKAGES AND ENVIRONMENT SETUP
%%%%%%%%%%%%%%%%%%%%%%%%%%%%%%%%%%%%%%%%%%%%%%%%%%%%%%%%%%%%%%%%%%%%%%%%%%%%%%%%

\usepackage{xparse} % For enhanced command definitions, coupled with transposeRow
\usepackage{xstring} % For string parsing, e.g., \IfStrEq{#1}{dot}{}{}
\usepackage{ifthen} % For semantic selectors used by frame and notation macros.
\usepackage{mathtools} % Core display math and \DeclareMathOperator for legacy direct inputs.
\usepackage{amsfonts,amssymb}
\usepackage{enumitem} % Structured enumerate labels used throughout the reports.

\ExplSyntaxOn
\NewDocumentCommand{\yuquanIfBlankTF}{mmm}{%
  \tl_if_blank:nTF {#1} {#2} {#3}%
}
\ExplSyntaxOff

\usepackage{mdframed}
\usepackage{natbib} % For \citet command

\newmdenv[
shadow=true,
shadowsize=6pt,
linewidth=0pt,
backgroundcolor=gray!20,
shadowcolor=black!40,
roundcorner=10pt
]{shadowed}

\usepackage{makecell}

%%%%%%%%%%%%%%%%%%%%%%%%%%% For  the cross mark
\usepackage[utf8]{inputenc}
\usepackage[english]{babel}

\usepackage{color}
% %%%%%%%%%%%%%%%%%%%%%%%%% Tikz packages:
\usepackage{tikz}
\usepackage{tikz-layers}
\usepackage{array}
%
% the positioning-plus library (is on TeX.sx),
% the node-families library (is on TeX.sx),
% the fit library (loaded by positioning-plus),
% the backgrounds library to draw stuff behind other stuff,
% the calc library for some funny coordinate calculations and the
% shapes.geometric library for the ellipse shape

\usetikzlibrary{shapes.geometric,automata,positioning,arrows,shadows,backgrounds,calc,fit}

\usepackage{setspace}




%%%%%%%%%%%%%%%%%%%%%%%%%% MATLAB setup

\usepackage{listings}
\lstset{
  language=Matlab,                		% choose the language of the code
  %	basicstyle=10pt,       				% the size of the fonts that are used for the code
  numbers=left,                  			% where to put the line-numbers
  numberstyle=\footnotesize,      		% the size of the fonts that are used for the line-numbers
  stepnumber=1,                   			% the step between two line-numbers. If it's 1 each line will be numbered
  numbersep=5pt,                  		% how far the line-numbers are from the code
  %	backgroundcolor=\color{white},  	% choose the background color. You must add \usepackage{color}
  showspaces=false,               		% show spaces adding particular underscores
  showstringspaces=false,         		% underline spaces within strings
  showtabs=false,                 			% show tabs within strings adding particular underscores
  %	frame=single,	                			% adds a frame around the code
  %	tabsize=2,                				% sets default tabsize to 2 spaces
  %	captionpos=b,                   			% sets the caption-position to bottom
  breaklines=true,                			% sets automatic line breaking
  breakatwhitespace=false,        		% sets if automatic breaks should only happen at whitespace
  escapeinside={\%*}{*)}          		% if you want to add a comment within your code
}


%%%%%%%%%%%%%%%%%%%%%%%%%% Self-defined commands:


\usepackage{xspace}
\lstset{frame=tb,
  language=C++,
  backgroundcolor=\color{black!5}, % set backgroundcolor
  basicstyle=\footnotesize,% basic font setting
  aboveskip=3mm,
  belowskip=3mm,
  showstringspaces=false,
  columns=flexible,
  numbers=none,
  numberstyle=\tiny\color{gray},
  keywordstyle=\color{blue},
  commentstyle=\color{dkgreen},
  stringstyle=\color{mauve},
  breaklines=true,
  breakatwhitespace=true,
  tabsize=3
}


\usepackage{soul}

\usepackage{pifont}% http://ctan.org/pkg/pifont

\usepackage{stackengine}
\newcommand\barbelow[1]{\stackunder[1.2pt]{$#1$}{\rule{.8ex}{.075ex}}}

% Algorithms.  Use the standard algorithmicx packages when available.  A
% small compatibility implementation keeps the reports buildable on minimal
% TeX installations; it supports the constructs used in this repository.
\IfFileExists{algorithm.sty}{%
  \usepackage{algorithm}% http://ctan.org/pkg/algorithms
  \usepackage{algorithmicx}
  \usepackage{algpseudocode}% http://ctan.org/pkg/algorithmicx
}{%
  \usepackage{float}
  \floatstyle{ruled}
  \newfloat{algorithm}{tbp}{loa}
  \floatname{algorithm}{Algorithm}
  \newlength{\algorithmicindent}
  \setlength{\algorithmicindent}{1.5em}
  \NewDocumentEnvironment{algorithmic}{O{}}{%
    \begin{list}{}{%
      \setlength{\leftmargin}{1.5em}%
      \setlength{\itemsep}{0.15em}%
      \setlength{\parsep}{0pt}%
    }%
  }{\end{list}}
  \newcommand{\State}{\item}
  \newcommand{\Procedure}[2]{\item\textbf{procedure} ##1\yuquanIfBlankTF{##2}{}{\ (##2)}}
  \newcommand{\EndProcedure}{\item\textbf{end procedure}}
  \newcommand{\For}[1]{\item\textbf{for} ##1 \textbf{do}}
  \newcommand{\EndFor}{\item\textbf{end for}}
  \newcommand{\If}[1]{\item\textbf{if} ##1 \textbf{then}}
  \newcommand{\Else}[1]{\item\textbf{else}\yuquanIfBlankTF{##1}{}{ ##1}}
  \newcommand{\EndIf}{\item\textbf{end if}}
  \newcommand{\While}[1]{\item\textbf{while} ##1 \textbf{do}}
  \newcommand{\EndWhile}{\item\textbf{end while}}
  \newcommand{\Comment}[1]{\hfill$\triangleright$~##1}
}

%%% predecessor:
\usepackage{dsfont}

%%%%%%%%%%%%%%%%%%%%%%%%%%%%%%%%%%%%%%%%%%%%%%%%%%%%%%%%%%%%%%%%%%%%%%%%%%%%%%%%
%% 2. COLOR DEFINITIONS
%%%%%%%%%%%%%%%%%%%%%%%%%%%%%%%%%%%%%%%%%%%%%%%%%%%%%%%%%%%%%%%%%%%%%%%%%%%%%%%%

\definecolor{frameColor}{RGB}{128,128,128}
\definecolor{initialStateColor}{RGB}{10,100,10}
\definecolor{impactStateColor}{RGB}{0,255,0}
\definecolor{admittanceStateColor}{RGB}{0,0,255}

\definecolor{resetStateColor}{RGB}{10,100,10}

\definecolor{stateColor}{RGB}{51,204,204}

\definecolor{controllerColor}{RGB}{10,10,255}
\definecolor{plannerColor}{RGB}{10,150,150}
\definecolor{estimatorColor}{RGB}{10,200,10}
\definecolor{robotColor}{RGB}{255,0,0}
\definecolor{modelColor}{RGB}{180,100,10}

%%% Enumeration

%%%%%%%%%%%%%%%%%%%%%%%%%%% Adjust font size
\definecolor{dkgreen}{rgb}{0,0.6,0}
\definecolor{gray}{rgb}{0.5,0.5,0.5}
\definecolor{mauve}{rgb}{0.58,0,0.82}

\definecolor{svaColor}{RGB}{7, 160, 2}
\definecolor{uncertainColor}{RGB}{255, 0, 0}

\definecolor{spatialVelColor}{RGB}{189, 38, 96}
\definecolor{geometricColor}{RGB}{66, 126, 245}
\definecolor{bodyVelColor}{RGB}{13, 191, 191}

\definecolor{velColor}{RGB}{50, 205, 50}

\definecolor{textColor}{RGB}{0, 11, 178}


\colorlet{svaColorVariable}{svaColor}
\colorlet{geometricColorVariable}{geometricColor}

% %%%%%%%%%%%%%%%%%%%%%%%%% comment
\definecolor{csrStateColor}{RGB}{255,255,153}
\definecolor{scrStateColor}{RGB}{153,204,0}
\definecolor{crStateColor}{RGB}{153,204,155}

%%%%%%%%%%%%%%%%%%%%%%%%%% DDP commands:

% Open-loop gain for computing the optimal control action

%%%%%%%%%%%%%%%%%%%%%%%%%%%%%%%%%%%%%%%%%%%%%%%%%%%%%%%%%%%%%%%%%%%%%%%%%%%%%%%%
%% 3. FORMATTING AND PRESENTATION
%%%%%%%%%%%%%%%%%%%%%%%%%%%%%%%%%%%%%%%%%%%%%%%%%%%%%%%%%%%%%%%%%%%%%%%%%%%%%%%%

\newcommand{\shadowRemark}[1]{
\begin{shadowed}
\begin{remark}
{#1}
\end{remark}
\end{shadowed}
}

\newcommand{\shadowProblem}[2]{
\begin{shadowed}
\infoBlock{Problem:} \hlBlock{#1}
{#2}
        \end{shadowed}
}

%%%%%%%%%%%%%%%%%%%%%%%%%%% Splitting lines within a table cell
\newcommand{\crossmark}{
  \ding{55}
}


%%%%%%%%%%%%%%%%%%%%%%%%%%% Font colour
\newcommand{\mcrtc}{mc$\_$rtc\xspace}
\newcommand{\sota}{state-of-the-art\xspace}

\newcommand{\poincare}{Poincar\'e\xspace}

%%%%%%%%%%%%%%%%%%%%%%%%%% Units:
\newcommand{\ith}{$i$-th\xspace}
%%%%%%%%%%%%%%%%%%%%%%%%%% Units:
\newcommand{\highlight}[1]{%
  \colorbox{gray!20}{$\displaystyle#1$}}

\newcommand{\highlightblue}[1]{%
  \colorbox{blue!20}{$\displaystyle#1$}}
\newcommand{\highlightgreen}[1]{%
  \colorbox{green!20}{$\displaystyle#1$}}


%%%%%%%%%%%%%%%%%%%%%%%%%% Bounds on tasks
\newcommand{\threshold}[1]{\Theta_{#1}}

%%%%%%%%%%%%%%%%%%%%%%%%%% Control Lyapunov function

%%%%%%%%%%%%%%%%%%%%%%%%%% Control Barrier function


%%%%%%%%%%%%%%%%%%%%%%%%%% Optimal Control
% Estimated quantity
\NewDocumentCommand{\cpp}{v}{%
\texttt{\textcolor{blue}{#1}}%
}

%%%%%%%%%%%%%%%%%%%%%%%%%%%%% New reference command:
\NewDocumentCommand{\secRef}{m}{%
Sec.~\ref{#1}}

\NewDocumentCommand{\remarkRef}{m}{%
Remark.~\ref{#1}}

\NewDocumentCommand{\tableRef}{m}{%
Table.~\ref{#1}}

\NewDocumentCommand{\algRef}{m}{%
Algorithm.~\ref{#1}}

\NewDocumentCommand{\lemmaRef}{m}{%
Lemma.~\ref{#1}}

\NewDocumentCommand{\theoremRef}{m}{%
Theorem.~\ref{#1}}

\NewDocumentCommand{\figRef}{m}{%
  Fig.~\ref{#1}}

\NewDocumentCommand{\appRef}{m}{%
  Appendix~\ref{#1}}

\NewDocumentCommand{\lineRef}{m}{%
line.~\ref{#1}}

\NewDocumentCommand\orderedTwoS{mm}%
  {$<$#1,#2$>$}
\NewDocumentCommand\orderedThreeS{mmm}%
  {$<$#1,#2,#3$>$}

\newif\ifBlockComment
\BlockCommentfalse %\BlockCommenttrue
% %%%%%%%%%%%%%%%%%%%%%%%%%


\newcommand{\comment}[1]{{\bf \colorbox{teal}{\textcolor{white}{#1}}}}
\newcommand{\fix}[1]{{\bf \colorbox{purple}{\textcolor{white}{#1}}}}
% %%%%%%%%%%%%%%%%%%%%%%%%% Highlight with colors
\sethlcolor{teal}
\newcommand{\hlBlock}[1]{{\sethlcolor{teal}\color{white}\hl{#1}}}
\newcommand{\errorBlock}[1]{{\sethlcolor{violet}\color{white}\hl{#1}}}
% Question block
\newcommand{\qtBlock}[1]{{\sethlcolor{violet}\color{white}\hl{#1}}}
\newcommand{\infoBlock}[1]{{\sethlcolor{olive}\color{white}\hl{#1}}}
\newcommand{\bashBlock}[1]{{\sethlcolor{black}\color{white}\hl{#1}}}

\newcommand{\tBlock}[1]{{\bf \colorbox{teal}{\textcolor{white}{#1}}}}
\newcommand{\oBlock}[1]{{\bf \colorbox{olive}{\textcolor{white}{#1}}}}

\newcommand{\bBlock}[1]{{\sethlcolor{brown}\color{white}\hl{#1}}}


\newcommand{\fixedwidth}[1]{
{\parbox{\dimexpr\linewidth-20\fboxsep}{#1}}
}
\newcommand{\wrongBlock}[1]{{\color{brown}\ding{55}}~~{\bf \colorbox{brown}{\textcolor{white}{#1}}}}
\newcommand{\correctBlock}[1]{
{{\color{olive}\ding{51}}~{\bf \colorbox{teal}{\textcolor{white}{
	\fixedwidth{#1}
	}}}}
}
% Delay the lightweight fallback until every package has been loaded.  This
% lets reports opt into todonotes after inputting this legacy configuration
% without triggering a duplicate \todo definition.
\AtBeginDocument{%
  \providecommand{\todo}[1]{To do:~
    \textcolor{bodyVelColor}{#1}%
  }%
}


% The critical frequency
\newcommand{\cmark}{\ding{51}}%
\newcommand{\xmark}{\ding{55}}%


% The motion constraint Jacobian

%%%%%%%%%%%%%%%%%%%%%%%%%%%%%%%%%%%%%%%%%%%%%%%%%%%%%%%%%%%%%%%%%%%%%%%%%%%%%%%%
%% 4. UNITS
%%%%%%%%%%%%%%%%%%%%%%%%%%%%%%%%%%%%%%%%%%%%%%%%%%%%%%%%%%%%%%%%%%%%%%%%%%%%%%%%

\newcommand{\unitHz}{~Hz\xspace}
\newcommand{\unitMs}{~ms\xspace}
\newcommand{\unitTime}{~s\xspace}
\newcommand{\unitDegree}{~deg.\xspace}
\newcommand{\unitMass}{~kg\xspace}
\newcommand{\unitMeter}{~m\xspace}
\newcommand{\unitHertz}{~N/m$^{3/2}$\xspace}
\newcommand{\unitGPa}{~GPa\xspace}


\newcommand{\unitPosTS}{~m\xspace}
\newcommand{\unitPosJS}{~rad\xspace}
\newcommand{\unitVelTS}{~m/s\xspace}
\newcommand{\unitVelJS}{~rad/s\xspace}
\newcommand{\unitAccTS}{~m/s$^2$\xspace}
\newcommand{\unitAccVel}{~rad/s$^2$\xspace}
\newcommand{\unitForce}{~N\xspace}

%%%%%%%%%%%%%%%%%%%%%%%%%%%
\newcommand{\unitImpulse}{~N$\cdot$s\xspace}

%%%%%%%%%%%%%%%%%%%%%%%%%% LQR commands:

%%%%%%%%%%%%%%%%%%%%%%%%%%%%%%%%%%%%%%%%%%%%%%%%%%%%%%%%%%%%%%%%%%%%%%%%%%%%%%%%
%% 5. BASIC MATH NOTATION
%%%%%%%%%%%%%%%%%%%%%%%%%%%%%%%%%%%%%%%%%%%%%%%%%%%%%%%%%%%%%%%%%%%%%%%%%%%%%%%%

\newcommand{\setDef}[2]{
  \defeq \{#1 \mid #2 \}
}

\newcommand{\quickEq}[2]{
  \begin{equation}
    \label{#1}
    {#2}
  \end{equation}
}

\newcommand{\quickBM}[1]{
    \begin{bmatrix}
    {#1}
    \end{bmatrix}
}

\newcommand{\quickAd}[1]{
  \begin{aligned}
    {#1}
  \end{aligned}
}

\newcommand{\de}[1]{
  \frac{d}{dt}(#1)
}
%%%%%%%%%%%%%%%%%%%%%%%%%% Collision observer


\newcommand{\setSymbol}{\mathcal{X}}

\newcommand{\feasibleSymbol}{\mathcal{F}}
\newcommand{\viableSymbol}{\nu}

\newcommand{\feasibleSet}[1]{\feasibleSymbol_{#1}}
\newcommand{\viableSet}[1]{\viableSymbol_{#1}}

\newcommand{\set}[1]{\setSymbol_{#1}}
% Interval
\newcommand{\interval}[2]{[\ #1,~#2 ]\ }



% %%%%%%%%%%%%%%%%%%%%%%%%% COLOUR:

\newcommand{\quadraticPN}[4]{
#1\pTwo{#4}  #2{#4}   #3 = 0
}
%%%%%%%%%%%%%%%%%%%%%%%%% Momenta
\newcommand{\weightMatrix}[1]{\mathbb{#1}}
\newcommand{\weightmatrixH}{\mathbf{H}}
\newcommand{\weightmatrixW}{\mathbf{W}}
\newcommand{\weightmatrixQ}{\mathbf{Q}}


\newcommand{\matrixA}{\mathbf{A}}

\IfFileExists{accents.sty}{%
  \usepackage{accents}
  \newcommand{\ubar}[1]{\underaccent{\bar}{##1}}
}{%
  % A conventional underline is preferable to making the entire package fail
  % when the optional accents package is not installed.
  \newcommand{\ubar}[1]{\underline{##1}}
}
\newcommand{\boundSymbol}{l}

\newcommand{\applyBoundSuperscript}[1]{%
    \IfStrEqCase{#1}{%
        {min}{\text{min}}%
        {max}{\text{max}}%
    }[] % "default" string
}

\newcommand{\applyBoundSubscript}[1]{%
        #1
}


\NewDocumentCommand{\bound}{O{}O{}}{%
  \ensuremath{%
    \boundSymbol
    \yuquanIfBlankTF{#2}{}{^{\applyBoundSubscript{#2}}}%
    \yuquanIfBlankTF{#1}{}{_{\applyBoundSuperscript{#1}}}%
  }
}


\newcommand{\lowerBound}[1]{{#1}_{\text{min}}}
\newcommand{\upperBound}[1]{{#1}_{\text{max}}}
%%%%%%%%%%%%%%%%%%%%%%%%%%%%%%%%%%%%%%%%%%%%%%%%%%%%%%%%%%%%%%%%% Bound

\newcommand{\stateBound}[1]{
        \lowerBound{#1} \leq #1 \leq \upperBound{#1}
}
\newcommand{\stateAbsBound}[1]{
        | #1 |  \leq \upperBound{#1}
}
\newcommand{\stateBoundConstraint}[2]{
        \lowerBound{#1} \leq #2 \leq \upperBound{#1}
}
\newcommand{\doubleBound}[1]{\ubar{\bar{#1}}}


\newcommand{\defeq}{\vcentcolon=}
\newcommand{\eqdef}{=\vcentcolon}


% Operators
% New operators must defined as such to have them typeset
% correctly. As an example we define the Jacobian:
% -----------------------------------------------------------------
\DeclareMathOperator{\Jac}{Jac}


% Shortcuts.
% One can define new commands to shorten frequently used
% constructions. As an example, this defines the R and Z used
% for the real and integer numbers.
% -----------------------------------------------------------------
\def\RR{\mathbb{R}}
\def\ZZ{\mathbb{Z}}

% Integrate over time:  #1 starting time, #2 finishing time, #3 quantity
\newcommand{\timeIntegral}[3]{
\int^{#2}_{#1} #3 dt
}

\newcommand{\RRv}[1]{\mathbb{R}^{#1}}
\newcommand{\RRm}[2]{\mathbb{R}^{#1 \times #2}}

%%%%%%%%%%%%%%%%%%%%%%%%%%% Impact-handling

\DeclareMathOperator*{\argmin}{arg\,min} % thin space, limits underneath in displays
\DeclareMathOperator*{\lexmin}{lex\,min} % Lexicographic-order minimization

%%%%%%%%%%%%%%%%%%%%%%%%%% Probabilistic robotics


%%%%%%%%%%%%%%%%%%%%%%%%%% Collision avoidance notations
% A link
\newcommand{\fpd}[2]{{#1}_{#2}}
% Second-order partial derivative #1 quantity, #2 variable, #3 second variable
\newcommand{\spd}[3]{{#1}_{#2#3}}



% The argument denotes the initial state
\newcommand{\Vee}[1]{\left(#1\right)^\vee}


%%%%%%%%% Represent the derivative operator
\newcommand{\derivative}[1]{\frac{d}{dt}(#1)}
%%% #1:function, #2:var1, #3:var2
\newcommand{\pdv}[3]{
  \frac{\partial^2 #1}{\partial #2 \partial #3}
}
%%% #1:function, #2:var1
\newcommand{\pd}[2]{
  \frac{\partial #1}{\partial #2 }
}

\newcommand{\up}[1]{\skewMatrix{#1}}

\newcommand{\liebra}[2]{
  [#1,#2]
}

% #1 Notation of the twist coordinate
\newcommand{\liebraResult}[2]{
  \Vee{\liebraDefDef{#1}{#2}}
}
%%%% #1:linear velocity #2: angular velocity #3:linear velocity #4: angular velocity
\newcommand{\seLiebraResult}[4]{
(\cross{#2}{#3} - \cross{#4}{#1}, \skewMatrix{#2}#4 - \skewMatrix{#4}#2 )
}



\newcommand{\liebraDef}[2]{
  \skewMatrixTwo{\cross{#1}{#2}}
}

\newcommand{\liebraDefDef}[2]{
  #1#2 - #2#1
}

%%%%%%%%%%% ----------------------------------------

%% atan2
% #1: y #2: x
\newcommand{\twoNorm}[1]{
  {{\left\lVert#1\right\rVert}^2}
}

\ExplSyntaxOn
\NewDocumentCommand{\applyNormSymbols}{mm}{%
  \str_case:nnF {#1}
    {
      {two}{\left\lVert#2\right\rVert^2_2}% Squared Euclidean norm.
      {square}{\left\lVert#2\right\rVert^2}% Squared unspecified norm.
      {one}{\left\lVert#2\right\rVert_1}% One norm.
      {abs}{\left\lvert#2\right\rvert}% Scalar absolute value.
      {inf}{\left\lVert#2\right\rVert_{\infty}}% Infinity norm.
    }
    {\left\lVert#2\right\rVert}% Default norm.
}
\ExplSyntaxOff

\NewDocumentCommand{\norm}{O{}m}{%
  \applyNormSymbols{#1}{#2}%
}

\newcommand{\abs}[1]{\left\vert#1\right\vert}
\newcommand{\cross}[2]{ #1 \times #2 }
\newcommand{\crossDef}[1]{ % Define the 3-dimensional cross operator for a vector
\begin{bmatrix}
0 & -{#1}_3 & {#1}_2 \\
{#1}_3 & 0 &  -{#1}_1 \\
-{#1}_2 & {#1}_1 & 0
\end{bmatrix}
}
\newcommand{\crossDefTwo}[1]{ % Define the 3-dimensional cross operator for a vector
\begin{bmatrix}
0 & -\z{#1} & \y{#1} \\
\z{#1} & 0 &  -\x{#1} \\
-\y{#1} & \x{#1} & 0
\end{bmatrix}
}



\newcommand{\dualCross}[2]{ #1 \times^* #2 }
\newcommand{\dualCrossDef}[2]{ -#1 \times^\top #2 }
\newcommand{\innerP}[2]{
  \transpose{#1}#2
}

% Computes the outer product
\newcommand{\outerP}[2]{
  #1\transpose{#2}
}



% Calculates the reciprocal product between two twists
% #1:v_1,  #2:w_1,  #3:v_2, #4:w_2.
\newcommand{\agg}[3]{
\sum^{#2}_{#1=1}{#3}
}
%%% #1:index, #2:first index, #3:last index, #3 component
\newcommand{\aggTwo}[4]{
\sum^{#3}_{#1=#2}{#4}
}

%%% #1:index, #2:last index, #3 component
\newcommand{\union}[3]{
\bigcup^{#2}_{#1=1}{#3}
}


%%%%%%%%%%% #1:row index, #2:column index, #3 index (between 1 and dof)
\newcommand{\columnTwo}[2]{
  \begin{bmatrix}
    #1&  #2
\end{bmatrix}
}

\newcommand{\vectorTwo}[2]{
  \begin{bmatrix}
    #1\\
    #2
\end{bmatrix}
}

\newcommand{\vectorThree}[3]{
  \begin{bmatrix}
    #1\\
    #2\\
    #3
\end{bmatrix}
}

\newcommand{\vectorFour}[4]{
  \begin{bmatrix}
    #1\\
    #2\\
    #3\\
    #4
\end{bmatrix}
}
\newcommand{\vectorTwoRow}[2]{
  [#1^\top, #2^\top]^\top
}

\newcommand{\vectorThreeRow}[3]{
  [{#1}^\top, {#2}^\top, {#3}^\top]^\top
}
\newcommand{\matrixTwo}[4]{
  \begin{bmatrix}
    #1 & #2\\
    #3 & #4
\end{bmatrix}
}

\newcommand{\matrixThree}[9]{
  \begin{bmatrix}
    #1 & #2 & #3\\
    #4 & #5 & #6\\
    #7 & #8 & #9
\end{bmatrix}
}

\newcommand{\order}[1]{\mathcal{O}(#1)}


\newcommand{\SE}[1]{SE{(#1)}}
\newcommand{\se}[1]{se{(#1)}}

\newcommand{\SO}[1]{SO{(#1)}}
\renewcommand{\so}[1]{so{(#1)}}
\newcommand{\soPowerTwoDef}[1]{
  \projectionDef{#1} - \twoNorm{#1} \identityMatrix
}
\newcommand{\soPowerThreeDef}[1]{
  - \twoNorm{#1} \skewMatrix{#1}
}


\newcommand{\mapDef}[2]{
  : #1 \rightarrow #2
}


% transform a twist coordinate from coordinate #1 to coordinate #2
\newcommand{\quadraticObj}[2]{\transpose{#1}#2#1}

\newcommand{\complexNum}[2]{
  (#1 ~ #2*i)
}

%% #1 upper bound, #2 middle #3 lower bound
\newcommand{\saturation}[3]{
\left\{
  \begin{array}{lr}
    #1 \\
    #2 \\
    #3
  \end{array}
\right.
}
% #1 velocity #2 mass
\newcommand{\quadratic}[2]{\frac{1}{2} #1^\top #2 #1}
\newcommand{\squareTwo}[1]{ #1^\top #1}
\newcommand{\squared}[1]{  {#1}^2 }


\ExplSyntaxOn

 % Define a sequence variable to store processed items
 \seq_new:N \trSeq

% Define the command to process and display the list
\NewDocumentCommand{\transposeRow}{m}
 {

  % Clear the sequence at the start
  \seq_clear:N \trSeq

  % Process each item in the list and store in the sequence
  \clist_map_inline:nn { #1 }
   {
    \seq_put_left:Nn \trSeq { \transpose{##1} }
   }

  % Output all items in the sequence, separated by spaces (or any other delimiter)
  \seq_use:Nn \trSeq {}
 }
\ExplSyntaxOff


\ExplSyntaxOn

%% % Define a sequence variable to store processed items
\seq_new:N \imSeq

% Define the command to process and display the list
\NewDocumentCommand{\innerMany}{m}
 {
  % Clear the sequence at the start
  \seq_clear:N \imSeq

  % Process each item in the list and store in the sequence
  \clist_map_inline:nn { #1 }
   {
    \seq_put_left:Nn \imSeq { \transpose{##1} }
    \seq_put_right:Nn \imSeq { {##1} }
   }

  % Output all items in the sequence, separated by spaces (or any other delimiter)
  \seq_use:Nn \imSeq {}
 }
\ExplSyntaxOff


% # 1 vel #2 mass # 3 partial d of vel
\newcommand{\quadraticd}[3]{ #3^\top #2 #1 +  #1^\top #2 #3}
\newcommand{\pTwo}[1]{#1^{2}}


%%%%%%%%%%%%%%%%%%%%%%%%%% Cross product
\newcommand{\paral}[2]{#1 \parallel #2}
\newcommand{\perpen}[2]{#1 \perp #2}

\newcommand{\bs}{\boldsymbol}
\newcommand{\backfill}{\hfill\(\blacksquare\)}

\newcommand{\skewMatrix}[1]{\widehat{#1}}

\newcommand{\skewMatrixTwo}[1]{({#1})^{\widehat{~}}}


\newcommand{\vectriple}[3]{#1\times(#2 \times #3)}
\newcommand{\vectripleexpan}[3]{
({#1}\bullet{#3}){#2} - ({#1}\bullet{#2}){#3}
}
\newcommand{\vectripleexpannegate}[3]{
-({#1}\bullet{#3}){#2} + ({#1}\bullet{#2}){#3}
}
\newcommand{\scalartriple}[3]{{#1}\bullet({#2} \times {#3})}
\newcommand{\vcd}[3]{({#1}\bullet{#3}){#2} - ({#1}\bullet{#2}) {#3}}  % Vector cross product definition

\newcommand{\inner}[2]{#1\bullet#2} % Inner product
\newcommand{\innerDef}[2]{\transpose{{#1}}#2} % Inner product
\newcommand{\normal}[1]{\bs{n}_{{#1}}}

\newcommand{\nc}[1]{\frac{#1}{| #1| }}% normal calc

\newcommand{\transpose}[2][default]{
  \IfStrEq{#1}{inv}
  {
    % Inverse transpose:
    {#2}^{-\top}
  }
  {
    % Default:
    {#2}^\top
  }
}

\newcommand{\transInverse}[1]{#1^{-\top}}
\newcommand{\pseudoInverse}[1]{{#1}^{+}}
% When  #row >  #column
\newcommand{\pseudoInverseColumnDef}[1]{\inverse{( \transpose{#1} #1 )}\transpose{#1}}
% When  #column >  #row
\newcommand{\pseudoInverseRowDef}[1]{\transpose{#1}\inverse{( #1 \transpose{#1} )}}

\newcommand{\weightedPseudoInverse}[1]{#1^{\dagger}}
% % #1 Jacobian #2 The weight, e.g., the inertia matrix.
\newcommand{\weightedPseudoInverseRowDef}[2]{\inverse{#2}\transpose{#1}\inverse{( #1 \inverse{#2} \transpose{#1} )}}
\newcommand{\weightedPseudoInverseColumnDef}[2]{\inverse{( \transpose{#1} \inverse{#2} #1 )}\transpose{#1}\inverse{#2}}

\newcommand{\inverse}[2][default]{
\IfStrEq{#1}{trans}
{
{#2}^{-\top}
}{
{#2}^{-1}
}
}


%%%%%%%%%%%%%%%%%%%%%%%%%% Resultant torque
% The argument are: force, torque, application point of the input wrench, application point of the resultant wrench.
\newcommand{\resultantTorque}[4]{#2 + (#3 - #4)\times #1}
\newcommand{\resultantTorqueTwo}[4]{#2 + \pointer{#4}{#3}\times #1}

\newcommand{\torqueFrame}[2]{{\torque_{#1#2}}}

\newcommand{\frameSymbol}{\mathcal{F}}


\newcommand{\applyFrameSubscript}[1]{%
    \ifthenelse{\equal{#1}{inertial}}{\inertialFrame}{%
    \ifthenelse{\equal{#1}{cmm}}{\cmmFrame}{%
    \ifthenelse{\equal{#1}{com}}{\com}{%
    \ifthenelse{\equal{#1}{fb}}{\fbFrame}{%
    \ifthenelse{\equal{#1}{torso}}{\torsoFrame}{%
    \ifthenelse{\equal{#1}{ee}}{\text{e}}{%
    \ifthenelse{\equal{#1}{cp}}{\cpFrame}{%
    \text{$#1$} % Default case, just print the argument as text
    }}}}}}}%
}

\NewDocumentCommand{\cframe}{O{}}{{\frameSymbol_{\applyFrameSubscript{#1}}}}



\newcommand{\origin}{\point[][][\inertialFrame]}


% The Gravito-inertial wrench
\newcommand{\giwrench}{\bs{W}^{gi}}
% The contact wrench
\newcommand{\cwrench}{\bs{W}^c}
\newcommand{\numEE}{m}

\newcommand{\forceJacobian}[2]{\jacobian_{#1 #2}}

\newcommand{\configurationFrame}[2]{g_{#1 #2}}



\newcommand{\spatialJacobian}[2]{
  \textcolor{spatialVelColor}{
    \spatial{\jacobian}_{#1 #2}
  }
}

\newcommand{\spatialJacobiand}[2]{
  \textcolor{spatialVelColor}{
    \dot{\spatial{\jacobian}}_{#1 #2}
  }
}

\newcommand{\forceJacobianDef}[2]{
\begin{bmatrix}
  \rotation{#1}{#2} \\
  \pointer{#1}{#2} \times \rotation{#1}{#2}
\end{bmatrix}}

% Transform from #1 to #2
\newcommand{\svaForceJacobianDef}[2]{
  \textcolor{svaColorVariable}{
    \vectorTwo{
      \translationSkew{#2}{#1}\rotationInv{#1}{#2}
    }{
      {\rotationInv{#1}{#2}}
    }
  }
}
% Transform from #1 to #2
\newcommand{\svaForceJacobianDefTwo}[2]{
  \textcolor{svaColorVariable}{
    \vectorTwo{
      -\rotation{#2}{#1}\translationSkew{#1}{#2}
    }{
      {\rotation{#2}{#1}}
    }
  }
}


\newcommand{\mass}{\text{m}}

\newcommand{\applyNumSymbol}[1]{%
       \IfStrEq{#1}{upper}{\text{N}}{%
         \text{n}
    }%
}

\newcommand{\applyNumSuperscript}[1]{%
    \IfStrEqCase{#1}{%
    {float}{
    \text{fb}
    }%
    {actuated}{
    \text{at}
    }%
    {ieq}{
    \text{ieq}
    }%
    {eq}{
    \text{eq}
    }%
    }[] % "default" string
}
\newcommand{\applyNumSubscript}[1]{%
    \IfStrEqCase{#1}{%
    {obj}{
    \text{obj}
    }%
    {cons}{
    \text{cons}
    }%
    {joints}{
    \text{j}
    }%
    {contact}{
    \text{sc}
    }%
    {impact}{
    \impulse
    }%
    }[] % "default" string
}

\NewDocumentCommand{\num}{O{default}O{default}O{default}}{%
    \ensuremath{%
        % #3 Shall we use the lower case symbol n?
        \applyNumSymbol{#3}_{\applyNumSubscript{#1}}% Check the first optional argument for obj, cons, joints, contact, impact,
        % Check the second optional argument for actuated, floating-base
        ^{\applyNumSuperscript{#2}}
    }
}


\newcommand{\numContacts}{{\text{m}_{\text{sc}}}} % The number of sustained contacts

\newcommand{\numImpacts}{{\text{m}_{\text{im}}}}


\newcommand{\numFree}{\numEE - \numContacts -\numImpacts }
\newcommand{\contactPoint}{\text{c}}

% End-effector
\newcommand{\ee}{\text{e}}

%%%%%%%%%%%%%%%%%%%%%%%%%%%%%%%% wrench, force, and torque
\newcommand{\applyWrenchPrescript}[1]{%
#1
}

\newcommand{\applyWrenchSuperscript}[1]{%
    \IfStrEqCase{#1}{%
        {ref}{\text{ref}}%
        {des}{\text{des}}%
        {mea}{\circ}%
        {min}{\text{min}}%
        {max}{\text{max}}%
    }[#1] % "default" string
}

\newcommand{\applyRelativeFrameSubscript}[2]{%
    {#1
    \ifx\hfuzz#1\hfuzz
         % #1 is empty
    \else % #1 is not empty
        \ifx\hfuzz#2\hfuzz
        % #2 is empty
        \else
          ,
          \fi
    \fi
    #2}
}

\newcommand{\wrenchSymbol}{{\textcolor{textColor}{\text{\bf w}}}\xspace}

\NewDocumentCommand{\test}{O{}O{}}{%
    {\ensuremath{
        {\wrenchSymbol_{#1}}_{#2}
    }
    }
}


\NewDocumentCommand{\wrench}{O{}O{}O{}O{}}{%
    {\ensuremath{
        % #1 reference frame #2 application point frame
        ~^{\applyWrenchPrescript{#4}}
        \wrenchSymbol_{\applyRelativeFrameSubscript{#1}{#2}}
        % #3 ref, mea, or des
        ^{\applyWrenchSuperscript{#3}}
    }
    }
}

\newcommand{\applyForceScalarSuperscript}[1]{%
    \IfStrEqCase{#1}{%
        {x}{\text{x}}%
        {y}{\text{y}}%
        {z}{\text{z}}%
    }[] % "default" string
}


\newcommand{\forceSymbol}{{{\bf f}}\xspace}
\NewDocumentCommand{\force}{O{}O{}O{}O{}}{%
{    \ensuremath{
      %% ^{\applyWrenchPrescript{#4}}
        % #1 reference frame #2 application point frame
       \bs{f}_{\applyRelativeFrameSubscript{#1}{#2}}
        ^{
          \applyForceScalarSuperscript{#4}% #4 x,y, or z
          \ifx\hfuzz#3\hfuzz
          % #3 is empty
          \else % #3 is not empty
          \ifx\hfuzz#4\hfuzz
          % #4 is empty
          \else
          ,
          \fi
          \fi
          %% \ifx\hfuzz#3\hfuzz\else
          %% \ifx\hfuzz#4\hfuzz\else
          %%   ,\applyWrenchSuperscript{#3}%
          %% \fi
          %% \fi
          \applyWrenchSuperscript{#3}
          }
        }
}
}

\NewDocumentCommand{\torque}{O{}O{}O{}O{}}{%
{    \ensuremath{
        % #1 reference frame #2 application point frame
       \bs{\tau}_{\applyRelativeFrameSubscript{#1}{#2}}
        ^{
          \applyForceScalarSuperscript{#4}% #4 x,y, or z
          \ifx\hfuzz#3\hfuzz\else
          \ifx\hfuzz#4\hfuzz\else
            ,\applyWrenchSuperscript{#3}%
          \fi
          \fi
          \applyWrenchSuperscript{#3}
          }
        }
}
}

\newcommand{\applyDurationSubscript}[2]{%
    \IfStrEqCase{#1}{%
    {stance}{
      {#2}_{\text{st.}}
    }%
    {flight}{
      {#2}_{\text{flt.}}
    }%
    {nominal}{
      {#2}_{\text{nom.}}
    }%
    {horizon}{
      {#2}_{\text{hor.}}
    }%
    }[
    #2
    ] % "default" string
}


%%%%%%%%%%%%%%%%%%%%%%%%%%%%%%%% ZMP, COM, and DCM

\newcommand{\applySuperscript}[1]{%
    \IfStrEqCase{#1}{%
    {ref}{
    \text{ref}
    }%
    {des}{
    \text{des}
    }%
    {mea}{
    \circ
    }%
    {est}{
    \wedge
    }%
    {min}{
    \text{min}
    }%
    {max}{
    \text{max}
    }%
    }[#1] % "default" string
}

\newcommand{\applyOverheadSymbols}[2]{%
    \IfStrEqCase{#1}{%
    {dddot}{
    \dddot{#2}
    }%
    {ddot}{
    \ddot{#2}
    }%
    {dot}{
    \dot{#2}
    }%
    }[
    #2
    ] % "default" string
}

\newcommand{\applyPostscript}[1]{%
    \IfStrEqCase{#1}{%
    {time}{
    (t)
    }%
    {current}{
    (k)
    }%
    {next}{
    (k+1)
    }%
    {previous}{
    (k-1)
    }%
    }[#1] % "default" string
}

\newcommand{\applySubscript}[1]{%
    \IfStrEqCase{#1}{%
    {plane}{
    \text{x,y}
    }%
    {x}{
    \text{x}
    }%
    {y}{
    \text{y}
    }%
    {z}{
    \text{z}
    }%
    }[#1] % "default" string
}




\newcommand{\zmpSymbol}{\text{\bf z}}
\NewDocumentCommand{\zmp}{O{}O{}O{}O{}}{%
{    \ensuremath{%
        % Check the first optional argument for dot or double dot
        %% \applyOverheadSymbols{#1}\text{\bf z}
        \applyOverheadSymbols{#1}{\zmpSymbol}
        % Check the second optional argument for 'ref'
        ^{\applySuperscript{#2}}%
        _{\applySubscript{#3}}
        \applyPostscript{#4}
    }
    }
}



\newcommand{\comSymbol}{{\textcolor{textColor}{\text{\bf c}}}\xspace}
\NewDocumentCommand{\com}{O{}O{}O{}O{}}{%
{    \ensuremath{%
        % Check the first optional argument for dot or double dot
        \applyOverheadSymbols{#1}{\comSymbol}
        % Check the second optional argument for 'ref'
        ^{\applySuperscript{#2}}
        _{\applySubscript{#3}}
        \applyPostscript{#4}
    }
    }
}

\newcommand{\dcmSymbol}{\textcolor{textColor}{\bs{\xi}}}
\NewDocumentCommand{\dcm}{O{}O{}O{}O{}}{%
{    \ensuremath{%
        % Check the first optional argument for dot or double dot
        \applyOverheadSymbols{#1}{\dcmSymbol}
        % Check the second optional argument for 'ref'
        ^{\applySuperscript{#2}}%
        _{\applySubscript{#3}}
        \applyPostscript{#4}
    }
    }
}

\newcommand{\pointSymbol}{\boldsymbol{q}}
\NewDocumentCommand{\point}{O{}O{}O{}O{}}{%
  {\ensuremath{%
        %% % Check the first optional argument for dot or double dot
        %% \IfStrEq{#1}{ddot}{\ddot{\boldsymbol{q}}}{%
        %%     \IfStrEq{#1}{dot}{\dot{\boldsymbol{q}}}{\boldsymbol{q}}%
        %% }%
        \applyOverheadSymbols{#1}{\pointSymbol}
        % Check the second optional argument for 'ref'
        ^{\applySuperscript{#2}}%
        % The mandatory subscript
        \ifx&#3&%
        % #3 is empty
        \else
        _{#3}%
        \fi
        % Post script
        \applyPostscript{#4}
    }
  }
}


\newcommand{\applySigmoidSymbols}[2]{%
    \IfStrEqCase{#1}{%
    {def}{
      \frac{\exp{#2}}{ 1 + \exp{#2}}
    }%
    %
    }[
      \sigmoidSymbol(#2)
    ] % "default" string
}


\newcommand{\sigmoidSymbol}{\sigma}
\NewDocumentCommand{\sigmoid}{O{default}O{default}}{%
  {\ensuremath{%
    \applySigmoidSymbols{#2}{#1} % #1, argument #2, def or not
    }
  }
}


\newcommand{\durationSymbol}{\text{T}}
\NewDocumentCommand{\duration}{O{default}}{%
  {\ensuremath{%
    \applyDurationSubscript{#1}{\durationSymbol}
    }
  }
}

\NewDocumentCommand{\pointFrame}{O{default}O{default}O{default}O{default}}{%
{    \ensuremath{%
       \point[#3][#4]_{\applyRelativeFrameSubscript{#1}{#2}} % #1 Reference frame, #2 Notation
    }
}
}


\newcommand{\forceScalar}{f}
\newcommand{\forceFrame}[2]{{\force^{#1}_{#2}}}


% The arguments are: (1) frame (2) Application point.
\newcommand{\impulseFrame}[2]{{\impulse^{#1}_{#2}}}
\newcommand{\applyWrenchFrameSuperscript}[1]{%
    \IfStrEq{#1}{ref}{, \text{ref}}{%
        \IfStrEq{#1}{des}{, \text{des}}{%
            \IfStrEq{#1}{mea}{, \circ}{}%
        }%
    }%
}
\NewDocumentCommand{\wrenchFrame}{mmO{default}}{%
  % #1 = representation frame
  % #2 = application point
  % #3 = optional argument, specifies ref, des, or mea
   \wrench^{{#1}\applyWrenchFrameSuperscript{#3}}_{#2}
}

% The foundamental subspaces
% #1 The matrix
\newcommand{\nullSpace}[1]{{N(#1)}}
% #1 The range (column) space
\newcommand{\colSpace}[1]{C(#1)}
\newcommand{\csProjector}[1]{{P}_{#1}}
% #1 Jacobian with: #col > #row
\newcommand{\csProjectorDef}[1]{\pseudoInverse{#1}#1}
% #1 Twist of a joint with: #col == 1
\newcommand{\jcsProjectorDef}[1]{#1\transpose{#1}}

\newcommand{\nsProjector}[1]{{N}_{#1}}
% #1 Jacobian with: #col > #row
\newcommand{\nsProjectorDef}[1]{(\identityMatrix - \csProjector{#1})}
\newcommand{\nsProjectorDefDef}[1]{(\identityMatrix - \csProjectorDef{#1})}
% #1 Twist of a joint with: #col == 1
\newcommand{\jnsProjectorDef}[1]{(\identityMatrix - \jcsProjectorDef{#1})}


% #1 The row  space
\newcommand{\rowSpace}[1]{C(\transpose{#1})}
% #1 The null space of the rows
\newcommand{\rowNullSpace}[1]{N(\transpose{#1})}
\newcommand{\rnsProjector}[1]{{N}_{\transpose{#1}}}
% #1 Jacobian with: #col > #row
\newcommand{\rnsProjectorDef}[1]{\transpose{(\identityMatrix - \csProjectorDef{#1})}}
% #1 Twist of a joint with: #col == 1
\newcommand{\jrnsProjectorDef}[1]{\transpose{(\identityMatrix - \jcsProjectorDef{#1})}}

\newcommand{\rsProjector}[1]{{P}_{\transpose{#1}}}
\newcommand{\rsProjectorDef}[1]{\transpose{(\csProjector{#1})}}
% #1 Twist of a joint with: #col == 1
\newcommand{\jrsProjectorDef}[1]{\transpose{(\jcsProjectorDef{#1})}}

% Joint row null space projector (to obtain the constraint wrench P*LagrangeMultiplier)
\newcommand{\jcwProjector}{T}

% The joint torque null-space projector of a redundant robot.
% #1 The Jacobian: v = J \dot{q}
\newcommand{\jtnsProjector}[1]{N_{\transpose{#1}}}
\newcommand{\jtnsProjectorDef}[1]{\transpose{(\identityMatrix - \gJacobianInv{#1} #1)}}

% The generalized Jacobain inverse (inertia-weighted pseudoinverse):
% #1 The Jacobian: v = J \dot{q},
% \gJacobianInv{J} = J^{-1}, when the manipulator is non-redundant.
\newcommand{\gJacobianInv}[1]{\overline{#1}}
\newcommand{\gJacobianInvDef}[1]{\inverse{\inertiaMatrix} \transpose{#1}\osInertia }

% The arguments are: (1) Equivalent frame name (2) Application point of the force
\newcommand{\forceGraspMatrix}[2]{
\begin{bmatrix}
    \rotation{\bs{#1}}{#2}^\top \\
    \pointer{#1}{#2} \times \rotation{#1}{#2}^\top
  \end{bmatrix}
}
\newcommand{\graspMatrix}[2]{
  {G_{#1#2}}
}

\newcommand{\graspMatrixDef}[2]{
\begin{bmatrix}
  \rotation{#1}{#2} & 0\\
  \pointer{#1}{#2} \times \rotation{#1}{#2} & \rotation{#1}{#2}
\end{bmatrix}
}


\newcommand{\pendulumC}{w}
\newcommand{\const}{\text{C}}
\renewcommand{\exp}[1]{e^{#1}}

\newcommand{\pendulumCDef}{\sqrt{\frac{\gSymbol}{\com[][][z]}}}


% Hadamard product
\newcommand{\hProduct}[2]{{#1}\bigodot {#2}}


%%%%%%%%%%%%%%%%%%%%%%%%%%%%%%%%%%%%%%%%%%%%%%%%%%%%%%%%%%%%%%%%%%%%%%%%%%%%%%%%
%% 6. COORDINATE FRAMES
%%%%%%%%%%%%%%%%%%%%%%%%%%%%%%%%%%%%%%%%%%%%%%%%%%%%%%%%%%%%%%%%%%%%%%%%%%%%%%%%

\newcommand{\dFrame}{\text{des}}

% The reference frame
\newcommand{\rFrame}{\text{ref}}

% Operational space coriolios and centrifugal forces

\newcommand{\coFrame}{\mathcal{C}}



\newcommand{\hybridFrame}{H}
\newcommand{\cmmFrame}{\text{G}}
\newcommand{\fbFrame}{\text{B}}
\newcommand{\torsoFrame}{\text{t}}
\newcommand{\inertialFrame}{\text{O}}
\newcommand{\anchor}{\text{A}} % The anchor frame of a legged robot
% The inertialFrame
\newcommand{\iFrame}{\text{O}}

\newcommand{\bodyFrame}{\text{b}}
\newcommand{\spatialFrame}{\text{s}}


\newcommand{\nz}{\text{\bf n}}
\newcommand{\nzDef}{
\begin{bmatrix}
0\\
0\\
1
\end{bmatrix}
}

\newcommand{\linkFrame}{L}
\newcommand{\jointFrame}{J}

% Gravity force

%%%%%%%%%%%%%%%%%%%%%%%%%%%%%%%%%%%%%%%%%%%%%%%%%%%%%%%%%%%%%%%%%%%%%%%%%%%%%%%%
%% 7. RIGID BODY GEOMETRY
%%%%%%%%%%%%%%%%%%%%%%%%%%%%%%%%%%%%%%%%%%%%%%%%%%%%%%%%%%%%%%%%%%%%%%%%%%%%%%%%

\newcommand{\epMapDef}[1]{
\identityMatrix + #1  + \frac{#1^2}{2!}  + \frac{#1^3}{3!} + \ldots
}
%%% #1:unit skew matrix #2 realnumber
\newcommand{\unitSkewExpDef}[2]{
\identityMatrix + \skewMatrix{#1}sin(#2) + \skewMatrix{#1}^2(1-cos(#2))
}

%%%%%%%%%%% -------------------------------------- Jacobian:
%%%% #1 reference frame #2 body frame #3 joint angle
\newcommand{\transformConfig}[3]{
  \textcolor{geometricColorVariable}{
    \transform{#1}{#2}(#3)
  }}
\newcommand{\transformConfigDef}[3]{
  \textcolor{geometricColorVariable}{
    \epMap{\twist{#1}{#2}#3} \transform{#1}{#2}(0)
  }}

\newcommand{\transformdConfigDef}[3]{
  \textcolor{geometricColorVariable}{
    \twist{#1}{#2} \dot{#3}\epMap{\twist{#1}{#2}#3} \transform{#1}{#2}(0)
  }}
\newcommand{\transformdConfigDefTwo}[3]{
  \textcolor{geometricColorVariable}{
    \epMap{\twist{#1}{#2}#3}\twist{#1}{#2} \dot{#3} \transform{#1}{#2}(0)
  }}

\newcommand{\transformddConfigDef}[3]{
  \textcolor{geometricColorVariable}{
    % \twist \ddot{#3}\epMap{\twist#3} \transform{#1}{#2}(0)
    % +\twist \dot{#3}\twist \dot{#3}\epMap{\twist#3} \transform{#1}{#2}(0)
    (\twist{#1}{#2} \ddot{#3}%\epMap{\twist#3} \transform{#1}{#2}(0)
    +\twist{#1}{#2} \twist{#1}{#2} \dot{#3}^2)\epMap{\twist{#1}{#2}#3} \transform{#1}{#2}(0)
  }}
\newcommand{\transformddConfigDefTwo}[3]{
  \textcolor{geometricColorVariable}{
    % \twist \ddot{#3}\epMap{\twist#3} \transform{#1}{#2}(0)
    % +\twist \dot{#3}\twist \dot{#3}\epMap{\twist#3} \transform{#1}{#2}(0)
    \epMap{\twist{#1}{#2}#3} (\twist{#1}{#2} \ddot{#3}%\epMap{\twist#3} \transform{#1}{#2}(0)
    +\twist{#1}{#2} \twist{#1}{#2} \dot{#3}^2)\transform{#1}{#2}(0)
  }}

% #1:reference frame, #:2 end-effector frame
\newcommand{\poeDef}[2]{
  \textcolor{geometricColorVariable}{
    \twistEp{#1}{1}{1}\twistEp{#1}{2}{2}\ldots\twistEp{#1}{#2}{#2}\transformConfig{#1}{#2}{\zeroVector}
  }}
\newcommand{\invPoeDef}[2]{
  \textcolor{geometricColorVariable}{
    \transformInvConfig{#1}{#2}{\zeroVector}\invTwistEp{#1}{#2}{#2}\ldots\invTwistEp{#1}{2}{2}\invTwistEp{#1}{1}{1}
  }}

%%%% #1 reference frame #2 body frame #3 joint angle
\newcommand{\transformInvConfig}[3]{
  \textcolor{geometricColorVariable}{
    \transformInv{#1}{#2}(#3)
  }}
\newcommand{\transformdConfig}[3]{
  \textcolor{geometricColorVariable}{
    \transformd{#1}{#2}(#3)
  }}
\newcommand{\transformddConfig}[3]{
 \textcolor{geometricColorVariable}{
   \transformdd{#1}{#2}(#3)
 }}
\newcommand{\transformInvConfigDef}[3]{
  \textcolor{geometricColorVariable}{
  \transformInv{#1}{#2}(0) \epMap{-\twist{#1}{#2}#3}
  }}



%%% Lie bracket
\newcommand{\epMap}[1]{e^{#1}}

\newcommand{\rotation}[2]{R_{#1 #2}}




\newcommand{\rotationInv}[2]{R^{\top}_{#1 #2}}
\newcommand{\rotationInvD}[2]{\dot{R}^{\top}_{#1 #2}}
\newcommand{\rotationInvDD}[2]{\ddot{R}^{\top}_{#1 #2}}
\newcommand{\rotationd}[2]{\dot{R}_{#1 #2}}
\newcommand{\rotationdd}[2]{\ddot{R}_{#1 #2}}

% #1 reference frame (point) #2 point
\newcommand{\translationDef}[2]{
  \textcolor{geometricColorVariable}{
    \pointer{#1}{#2} =  #2 - #1
}}


\newcommand{\translation}[2]{
  \vc{p}_{#1 #2}
}

\newcommand{\translationInv}[2]{
   \textcolor{geometricColorVariable}{
     \translation{#2}{#1}
     }
}
\newcommand{\translationInvSkew}[2]{
   \textcolor{geometricColorVariable}{
     \translationSkew{#2}{#1}
     }
}
\newcommand{\translationInvDef}[2]{
   \textcolor{geometricColorVariable}{
     -\rotationInv{#1}{#2}\translation{#1}{#2}
     }
}


% Complexity of an algorithm
\newcommand{\translationSkew}[2]{\skewMatrix{\vc{p}}_{#1 #2}}
\newcommand{\translationdSkew}[2]{\skewMatrix{\dot{\vc{p}}}_{#1 #2}}
\newcommand{\translationd}[2]{\dot{\vc{p}}_{#1 #2}}
\newcommand{\translationdd}[2]{{\ddot{\vc{p}}_{#1 #2}}}

\newcommand{\identityMatrix}{\mathds{1}} %matrix % alternativ \mathds{1} or {\mathbb{1}}
\newcommand{\idMatrix}[1]{{\identityMatrix_{#1}}} %matrix % alternativ \mathds{1} or {\mathbb{1}}

\newcommand{\zeroVector}{\mathbf{0}} %matrix % alternativ \mathds{0} or {\mathbb{0}} or \mathcal{O}
\newcommand{\zeroMatrix}{0} %matrix % alternativ \mathds{0} or {\mathbb{0}} or \mathcal{O}

\newcommand{\transform}[2]{
  \textcolor{geometricColorVariable}{
    g_{#1#2}
  }}

\newcommand{\transformd}[2]{
 \textcolor{geometricColorVariable}{
   \dot{g}_{#1#2}
 }}
\newcommand{\transformdDef}[2]{
 \textcolor{geometricColorVariable}{
   \matrixTwo{
     \rotationd{#1}{#2}
   }{
     \translationd{#1}{#2}
   }{
     \zeroVector
   }{
     0
   }
 }}
\newcommand{\transformdd}[2]{
 \textcolor{geometricColorVariable}{
   \ddot{g}_{#1#2}
 }}
\newcommand{\transformddDef}[2]{
 \textcolor{geometricColorVariable}{
   \matrixTwo{
     \rotationdd{#1}{#2}
   }{
     \translationdd{#1}{#2}
   }{
     \zeroVector
   }{
     0
   }
 }}
\newcommand{\transformInv}[2]{
  \textcolor{geometricColorVariable}{
    g^{-1}_{#1#2}
  }}
\newcommand{\transformInvD}[2]{
  \textcolor{geometricColorVariable}{
    % \dot{g}^{-1}_{#1#2}
    \derivative{g^{-1}_{#1#2}}
  }}
\newcommand{\transformInvDConfig}[3]{
  \textcolor{geometricColorVariable}{
    % \dot{g}^{-1}_{#1#2}
    \derivative{g^{-1}_{#1#2}(#3)}
  }}
\newcommand{\transformInvDDef}[2]{
  \textcolor{geometricColorVariable}{
    \matrixTwo{
      \rotationInvD{#1}{#2}
    }{
      -\rotationInvD{#1}{#2} \translation{#1}{#2}  - \rotationInv{#1}{#2} \translationd{#1}{#2}
    }{
      \zeroVector
    }{
      0
    }
  }}

\newcommand{\transformDef}[2]{
  \textcolor{geometricColorVariable}{
    \matrixTwo{
      \rotation{#1}{#2}
    }{
      \translation{#1}{#2}
    }{
      \zeroVector
    }{
      1
    }
  }
}

\newcommand{\transformInvDef}[2]{
  \textcolor{geometricColorVariable}{
    \matrixTwo{
      \rotationInv{#1}{#2}
    }{
      -\rotationInv{#1}{#2} \translation{#1}{#2}
    }{
      \zeroVector
    }{
      1
    }
  }
}
\newcommand{\adg}[2]{
  \textcolor{geometricColorVariable}{
    Ad_{g_{#1#2}}
  }
}

\newcommand{\adgDef}[2]{
  \textcolor{geometricColorVariable}{
  \begin{bmatrix}
    \rotation{#1}{#2} & \translationSkew{#1}{#2}\rotation{#1}{#2} \\
    \zeroMatrix & \rotation{#1}{#2}
  \end{bmatrix}
}}

\newcommand{\adgInv}[2]{
  \textcolor{geometricColorVariable}{
  Ad^{-1}_{g_{#1#2}}
}}
\newcommand{\adgTwoInv}[2]{
  \textcolor{geometricColorVariable}{
  Ad_{g^{-1}_{#1#2}}
}}

\newcommand{\adgTrans}[2]{
  \textcolor{geometricColorVariable}{
  Ad^{\top}_{g_{#1#2}}
}}
\newcommand{\adgTransDef}[2]{
  \textcolor{geometricColorVariable}{
  \begin{bmatrix}
    \rotationInv{#1}{#2} & \zeroMatrix \\
     -\rotationInv{#1}{#2}\translationSkew{#1}{#2} & \rotationInv{#1}{#2}
\end{bmatrix}
}}


\newcommand{\adgInvDefTwo}[2]{
  \textcolor{geometricColorVariable}{
  \begin{bmatrix}
    \rotationInv{#1}{#2} & -(\rotationInv{#1}{#2}\translation{#1}{#2})^{\skewMatrix{}}\rotationInv{#1}{#2} \\
    \zeroMatrix & \rotationInv{#1}{#2}
\end{bmatrix}
}}
\newcommand{\adgInvDef}[2]{
  \textcolor{geometricColorVariable}{
  \begin{bmatrix}
    \rotationInv{#1}{#2} & -\rotationInv{#1}{#2}\translationSkew{#1}{#2} \\
    \zeroMatrix & \rotationInv{#1}{#2}
\end{bmatrix}
}}

% From frame #1 to frame #2
\newcommand{\adgInvTrans}[2]{
  \textcolor{geometricColorVariable}{
  Ad^{\top}_{g^{-1}_{#1#2}}
}}

\newcommand{\adgInvTransDef}[2]{
  \textcolor{geometricColorVariable}{
  \begin{bmatrix}
    \rotation{#1}{#2} & \zeroMatrix \\
    \translationSkew{#1}{#2}\rotation{#1}{#2}  & \rotation{#1}{#2}
  \end{bmatrix}
  }
}



\newcommand{\adgTransform}[3]{
  \textcolor{geometricColorVariable}{
  \transform{#1}{#2} \skewMatrix{#3} \transformInv{#1}{#2}
}}


%%%%%%%%%%%%%%%%%%%%%%%%%%%%%%%%%%%%%%%%%%%%%%%%%%%%%%%%%%%%%%%%%%%%%%%%%%%%%%%%
%% 8. SCREW THEORY
%%%%%%%%%%%%%%%%%%%%%%%%%%%%%%%%%%%%%%%%%%%%%%%%%%%%%%%%%%%%%%%%%%%%%%%%%%%%%%%%

\newcommand{\atanTwo}[2]{
atan2({#1}, {#2})
}


%%%%%%%%%%% ----------------------------------------Scew


%%%%%%%%%%%% Transform defined by a screw
%%% #1:axis, #2:theta, #3:point, #4:pitch(finite)
\newcommand{\screwTransform}[4]{
\matrixTwo{
  \epMap{\skewMatrix{#1}#2}
}{
  (\identityMatrix - \epMap{\skewMatrix{#1}#2})#3
  + #4 #1 #2
}{
\zeroVector
}{
1
}
}

%%%%%%%%%%%% Wrench defined by a screw
%%% #1:magnitude, #2:point, #3:direction, #4:pitch(finite)
%%% #1: reference frame #2: screw name
\newcommand{\geoWrenchDef}[2]{
  \screwM_{#2} \vectorTwo{
    \angularVel_{#2}
  }{
    -\cross{\angularVel_{#2}}{\translation{\origin{#1}}{#2}} + \pitch_{#2}\angularVel_{#2}
  }
}

%%%%%% When the pitch is infinite:
%%% #1: screw name
\newcommand{\geoWrenchDefTwo}[1]{
  \screwM_{#1} \vectorTwo{
    \zeroVector
  }{
    \angularVel_{#1}
  }
}

\newcommand{\distance}{d}
%%% #1:screw 1, #2:screw 2
\newcommand{\screwAngle}[2]{
  \alpha_{#1#2}
}
\newcommand{\screwAngleDef}[2]{
atan2(\inner{(\cross{\angularVel_{#1}}{\angularVel_{#2}})}{\unitVec_{#1#2}}, \inner{\angularVel_{#1}}{\angularVel_{#2}} ).
}
\newcommand{\unitVec}{\bs{n}}

% Calculates the reciprocal product
\newcommand{\recip}[2]{
  #1 \odot{}   #2
}

\newcommand{\recipScrewDef}[2]{
  \screwM_{#1}\screwM_{#2}\left(
    (\pitch_{#1} + \pitch_{#2})\cos\screwAngle{#1}{#2} - \distance\sin\screwAngle{1}{2}
  \right)
}

%%%%%%%%%%%% Pitch:
\newcommand{\pitch}{
  h
}



% Twist2Screw:  #1 linear vel #2 angular vel
\newcommand{\pitchTwistDef}[2]{
  \frac{\innerP{#2}{#1}}{\twoNorm{#2}}
}

% Wrench2Screw: #1 force #2 moment
\newcommand{\pitchWrenchDef}[2]{
  \frac{\innerP{#1}{#2}}{\twoNorm{#1}}
}


\newcommand{\axis}{
  l
}

%%% #1:point, #2:angularVel direction
\newcommand{\screwAxisDef}[2]{
  \{ #1 + \lambda#2 : \lambda \in \RRv{} \}
}

% When angularVel is Not zero: #1:linearVel, #2:angularVel, #3:lambda (scalar)
\newcommand{\axisTwistDefThree}[3]{
\frac{\cross{#2}{#1}}{\twoNorm{#2}} + \axisTwistDefTwo{#1}{#3}
}

% When angularVel is zero: #1:linearVel, #2:lambda (scalar)
\newcommand{\axisTwistDefTwo}[2]{
   #1 #2
}

% When force is NOT zero: #1:force, #2:torque, #3:lambda (scalar)
\newcommand{\axisWrenchDefThree}[3]{
\frac{\cross{#1}{#2}}{\twoNorm{#1}} + \axisWrenchDefTwo{#2}{#3}
}

% When force is zero: #1 torque #2: lambda (scalar)
\newcommand{\axisWrenchDefTwo}[2]{
   #1 #2
}

%%%%%%%%%%%% Magnitude:
\newcommand{\screwM}{
  M
}

% Wrench2Screw: #1: force,  if force is nonzero. Otherwise #1 is the torque
% Twist2Screw: #1: angularVel,  if angularVel is nonzero. Otherwise #1 is the linearVel
\newcommand{\screwMDef}[1]{
  \norm{#1}
}
%%%%%%%%%%%% Twist:

% #1 axis #2 theta
\newcommand{\twistE}[2]{
  e^{\skewMatrix{#1}#2}
}


% Computes the iner product
\newcommand{\recipVelDef}[4]{
  \innerP{#1}{#4} + \innerP{#3}{#2}
}

% Calculates the reciprocal product between two wrenches
% #1:force,  #2:torque,  #3:force, #4:torque.
\newcommand{\recipWrenchDef}[4]{
  \innerP{#1}{#4} + \innerP{#3}{#2}
}

% Screw
\newcommand{\screw}{S}
%%% #1 screw notation
\newcommand{\screwDef}[1]{
  (\screwM_{#1}, \pitch_{#1}, \point[][][#1], \angularVel_{#1})
}


% Detail the wrench along a screw
\newcommand{\wrenchAS}[2]{S}

% Detail the twist about a screw
%%% #1:axis, #2:\point, #3 pitch
\newcommand{\linearTwistAS}[3]{
-\cross{#1}{#2} + #3#1
}

\newcommand{\linearTwistASTwo}[3]{
  \cross{#2}{#1} + #3#1
}

% Detail the twist about a screw, when the pitch is zero:
%%% #1:axis, #2:\point,
\newcommand{\linearTwistASRevolute}[2]{
-\cross{#1}{#2}
}
\newcommand{\linearTwistASRevoluteTwo}[2]{
{#2}\cross{#1}
}



%%%%%%%%%%% Composite rigid body (CRB) algorithm:
\newcommand{\twistDef}[2]{
  \matrixTwo{
    #2
  }{
    #1
  }{\zeroVector
  }{0
  }
}
%%%%%%%%%%%%%%% The exponential map of a twist:
%% #1 linear part #2 angularPart   #3 theta
\newcommand{\twistEpDef}[3]{
  \matrixTwo{
    \epMap{\skewMatrix{#2} #3}
  }{
    (\identityMatrix - \epMap{\skewMatrix{#2} #3})
    (\cross{#2}{#1})
    +
    \projectionDef{#2}#1#3
  }{
    \zeroVector
  }{
    1
  }
}
%% #1: reference frame, #2 twist notation #3 theta notation
\newcommand{\twistEp}[3]{
\epMap{\twist{#1}{#2} \theta_{#3}}
}
%% #1: reference frame, #2 twist notation #3 theta notation
\newcommand{\invTwistEp}[3]{
\epMap{-\twist{#1}{#2} \theta_{#3}}
}


%%%%%%%%%%%%%%%%%%%%%% Dynamics
%%% #1:index, #2:last index, #3 component
\newcommand{\twistFromScrew}[1]{
  \screwM_{#1}\vectorTwo{
    \linearTwistAS{\angularVel_{#1}}{\point[][][#1]}{\pitch_{#1}}
  }{
    \angularVel_{#1}
  }
}

\newcommand{\twistFromScrewRevolute}[1]{
  \screwM_{#1}\vectorTwo{
    \linearTwistASRevolute{\angularVel_{#1}}{\point[][][#1]}
  }{
    \angularVel_{#1}
  }
}


%%%%%%%%%%%%%%%%%%%%%%%%%%%%%%%%%%%%%%%%%%%%%%%%%%%%%%%%%%%%%%%%%%%%%%%%%%%%%%%%
%% 9. TWIST COORDINATES AND VELOCITIES
%%%%%%%%%%%%%%%%%%%%%%%%%%%%%%%%%%%%%%%%%%%%%%%%%%%%%%%%%%%%%%%%%%%%%%%%%%%%%%%%

\newcommand{\twistVel}{\bs{V}}


\newcommand{\twoCases}[4]{
\begin{dcases}
  {#1} & {#2} \\
  {#3} & {#4}
\end{dcases}
}
\newcommand{\crossSixD}[1]{
  \begin{bmatrix}
    \cross{\angularVel_{#1}}{} & \cross{\linearVel_{#1}}{}\\
    0 &\cross{\angularVel_{#1}}{}
  \end{bmatrix}
}

\newcommand{\dualCrossSixD}[1]{
  \begin{bmatrix}
    \cross{\angularVel_{#1}}{} & 0 \\
    \cross{\linearVel_{#1}}{} &\cross{\angularVel_{#1}}{}
  \end{bmatrix}
}


\newcommand{\twist}[2]{
  \textcolor{spatialVelColor}{
    \skewMatrix{\xi}_{#1#2}
  }
}

%%% #1: infinitesimal translation generator #2:so{3}
\newcommand{\twistC}[2]{
  \textcolor{spatialVelColor}{
    \xi_{#1#2}
  }
}

% #1:reference frame, #2:notation of the twist
\newcommand{\twistCd}[2]{
  \textcolor{spatialVelColor}{
    \dot{\xi}_{#1#2}
  }
}

% #1:reference frame, #2:notation of the twist
\newcommand{\spatialTwist}[2]{
  \textcolor{spatialVelColor}{
    \hat{\xi}'_{#1#2}
  }}

% #1:reference frame, #2:notation of the twist
\newcommand{\spatialTwistC}[2]{
  \textcolor{spatialVelColor}{
    \xi'_{#1#2}
  }}

% #1:reference frame, #2:notation of the twist
\newcommand{\spatialTwistCd}[2]{
  \textcolor{spatialVelColor}{
    \dot{\xi}'_{#1#2}
  }}

% #1:reference frame, #2:notation of the twist
\newcommand{\bodyTwistC}[2]{
  \textcolor{bodyVelColor}{
    \xi^{\dagger}_{#1#2}
  }}
% #1:reference frame, #2:index of the twist
\newcommand{\bodyTwistCd}[2]{
  \textcolor{bodyVelColor}{
    \dot{\xi}^{\dagger}_{#1#2}
  }}

\newcommand{\bodyTwist}[2]{
  \textcolor{bodyVelColor}{
    \hat{\xi}^{\dagger}_{#1#2}
  }}

% #1:reference frame, #2:index  of the twist, #3:index of the starting link
%%% t_0 means that starting time.
\newcommand{\spatialTwistCDef}[3]{
  \textcolor{spatialVelColor}{
    Ad_{\epMap{\twist{#1}{#3(t_0)} \theta_{#3}}\ldots \epMap{\twist{#1}{#2-1(t_0)} \theta_{#2-1}} }\twistC{#1}{#2(t_0)}
    }
}

% #1:reference frame, #2:index of the twist, #3:index of the last link
\newcommand{\bodyTwistCDef}[3]{
  \textcolor{bodyVelColor}{
    Ad^{-1}_{\epMap{\twist{#1}{#2(t_0)} \theta_{#2}}\ldots \epMap{\twist{s}{#3(t_0)} \theta_{#3}}\transformConfig{s}{#3}{0} }\twistC{#1}{#2(t_0)}
    }
}


%%% The twist coordinate
%%% #1: notation of the joint
\newcommand{\twistCRevolute}[1]{
\vectorTwo{
  \linearTwistASRevolute{\angularVel_{#1}}{\point[][][#1]}
  }{
    \angularVel_{#1}
  }
}

%%% #1: notation of the joint
\newcommand{\twistCPrismatic}[1]{
\vectorTwo{
  \linearVel_{#1}
  }{
    \zeroVector
  }
}

%%% #1:linearVel, #2: angularVel
\newcommand{\twistCDef}[2]{(#1, #2)}

%%% #1:notation of screw
\newcommand{\vc}[1]{\bs{#1}}


\newcommand{\caseTwo}[4]{
  \begin{cases}
    #1 & #2 \\
    #3 & #4
  \end{cases}
}
\newcommand{\caseThree}[6]{
  \begin{cases}
    #1 & #2 \\
    #3 & #4 \\
    #5 & #6 \\
  \end{cases}
}

\newcommand{\translationVel}{
  \bs{v}
}

\newcommand{\linearVel}{
\translationVel
}
\newcommand{\linearVelScalar}{
  v
}

\newcommand{\angularVel}{
  \bs{w}
}

\newcommand{\angularVelScalar}{
  w
}



%%% #1 reference frame #2 body frame
\newcommand{\spatialVel}[2]{
  \textcolor{spatialVelColor}{
    \vc{V}^{\text{s}}_{#1#2}
  }}

%%% #1 reference frame #2 body frame
\newcommand{\spatialVelDef}[2]{
  \textcolor{spatialVelColor}{
\vectorTwo{
  \spatialTVDef{#1}{#2}
  }{
    % (\rotationd{#1}{#2}\rotationInv{#1}{#2})^{\vee}
    \spatialRVDef{#1}{#2}
  }
}}
\newcommand{\spatialVelTwistC}[2]{
  \textcolor{spatialVelColor}{
\vectorTwo{
  \spatialTV{#1}{#2}
  }{
    % (\rotationd{#1}{#2}\rotationInv{#1}{#2})^{\vee}
    \spatialRV{#1}{#2}
  }
}}
\newcommand{\bodyVelTwistC}[2]{
  \textcolor{bodyVelColor}{
\vectorTwo{
  \bodyTV{#1}{#2}
  }{
    % (\rotationd{#1}{#2}\rotationInv{#1}{#2})^{\vee}
    \bodyRV{#1}{#2}
  }
}}

\newcommand{\spatialVelTwistDefTwo}[2]{
  \textcolor{spatialVelColor}{
    \transformdConfig{#1}{#2}{\theta}\transformInvConfig{#1}{#2}{\theta}
  }}

\newcommand{\bodyVelTwistDefTwo}[2]{
  \textcolor{spatialVelColor}{
    \transformInvConfig{#1}{#2}{\theta}\transformdConfig{#1}{#2}{\theta}
}}


% Defines the reference frame and body frame of the angular velocity
%%% #1 reference frame #2 body frame
\newcommand{\hybridVel}[2]{
  \textcolor{spatialVelColor}{
    \vc{V}^{h}_{#1#2}
  }}


\newcommand{\hybridVelDef}[2]{
  \textcolor{spatialVelColor}{
    \vectorTwo{
      \translationd{#1}{#2}
    }{
      \spatialRV{#1}{#2}
    }
  }}
\newcommand{\hybridVelDefDef}[2]{
  \textcolor{spatialVelColor}{
    \vectorTwo{
      \translationd{#1}{#2}
    }{
      \spatialRVDef{#1}{#2}
    }
  }}
%%%%%%%%%%%%%%%% Convert hybrid vel to body vel
%%% #1 Reference frame  #2 body frame
\newcommand{\hTwoB}[2]{
  \begin{bmatrix}
    \rotationInv{#1}{#2} & \zeroMatrix \\
    \zeroMatrix & \rotationInv{#1}{#2}
  \end{bmatrix}
}

%%% #1 Reference frame  #2 body frame
\newcommand{\bTwoH}[2]{
  \begin{bmatrix}
    \rotation{#1}{#2} & \zeroMatrix \\
    \zeroMatrix & \rotation{#1}{#2}
  \end{bmatrix}
}

%%%%%%%%%%%%%%%% Convert hybrid vel to body vel
\newcommand{\spatialTV}[2]{
  \textcolor{spatialVelColor}{
  \translationVel^{\spatialFrame}_{#1#2}
}}
\newcommand{\spatialTVDef}[2]{
  \textcolor{spatialVelColor}{
    -\rotationd{#1}{#2}\rotationInv{#1}{#2}\translation{#1}{#2} + \translationd{#1}{#2}
}}
\newcommand{\spatialTA}[2]{
  \textcolor{spatialVelColor}{
    \dot{\translationVel}^{\spatialFrame}_{#1#2}
  }}
\newcommand{\spatialTADef}[2]{
  \textcolor{spatialVelColor}{
    \translationdd{#1}{#2} - \spatialRATwist{#1}{#2}\translation{#1}{#2} - \spatialRVTwist{#1}{#2}\translationd{#1}{#2}
  }}
\newcommand{\spatialTADefTwo}[2]{
  \textcolor{spatialVelColor}{
    \translationdd{#1}{#2} + \translationSkew{#1}{#2}\spatialRA{#1}{#2} +\translationdSkew{#1}{#2} \spatialRV{#1}{#2}
  }}
\newcommand{\spatialTADefDef}[2]{
  \textcolor{spatialVelColor}{
\translationdd{#1}{#2} -\rotationdd{#1}{#2}\rotationInv{#1}{#2} \translation{#1}{#2}
-\rotationd{#1}{#2} \rotationInvD{#1}{#2} \translation{#1}{#2} - \rotationd{#1}{#2}\rotationInv{#1}{#2}\translationd{#1}{#2}
  }}
\newcommand{\spatialRADef}[2]{
  \textcolor{spatialVelColor}{
    \rotationdd{#1}{#2}\rotationInv{#1}{#2} + \rotationd{#1}{#2}\rotationInvD{#1}{#2}
  }}
\newcommand{\bodyRADef}[2]{
  \textcolor{bodyVelColor}{
    \rotationInv{#1}{#2}\rotationdd{#1}{#2} + \rotationInvD{#1}{#2}\rotationd{#1}{#2}
  }}
\newcommand{\spatialAcc}[2]{
  \textcolor{spatialVelColor}{
    \dot{\vc{V}}^{\spatialFrame}_{#1#2}
  }
}
\newcommand{\spatialAccTwist}[2]{
  \textcolor{spatialVelColor}{
    \skewMatrix{\dot{\vc{V}}}^{\spatialFrame}_{#1#2}
  }
}
\newcommand{\bodyAcc}[2]{
  \textcolor{bodyVelColor}{
    \dot{\vc{V}}^{\bodyFrame}_{#1#2}
  }
}
\newcommand{\bodyAccTwist}[2]{
  \textcolor{bodyVelColor}{
    \skewMatrix{\dot{\vc{V}}}^{\bodyFrame}_{#1#2}
  }
}
\newcommand{\bodyAccTwistDef}[2]{
  \textcolor{bodyVelColor}{
    \matrixTwo{
      \bodyRATwist{#1}{#2}
    }{
      \bodyTA{#1}{#2}
    }{
      \zeroVector
    }{
      0
    }
  }
}
\newcommand{\spatialAccTwistDef}[2]{
  \textcolor{spatialVelColor}{
    \matrixTwo{
      \spatialRATwist{#1}{#2}
    }{
      \spatialTA{#1}{#2}
    }{
      \zeroVector
    }{
      0
    }
  }
}
\newcommand{\spatialAccDef}[2]{
  \textcolor{spatialVelColor}{
    \matrixTwo{
      \spatialRATwist{#1}{#2}
    }{
      \translationdd{#1}{#2} - \spatialRATwist{#1}{#2}\translation{#1}{#2} - \spatialRVTwist{#1}{#2}\translationd{#1}{#2}
    }{
      \zeroVector
    }{
      0
    }
  }}

\newcommand{\spatialAccTwistC}[2]{
  \textcolor{spatialVelColor}{
    \vectorTwo{
      \translationdd{#1}{#2} - \spatialRATwist{#1}{#2}\translation{#1}{#2} - \spatialRVTwist{#1}{#2}\translationd{#1}{#2}
    }{
      \spatialRA{#1}{#2}
    }
  }}

\newcommand{\spatialAccDefDef}[2]{
  \textcolor{spatialVelColor}{
    \matrixTwo{
      \rotationdd{#1}{#2}\rotationInv{#1}{#2} + \rotationd{#1}{#2}\rotationInvD{#1}{#2}
    }{
      \translationdd{#1}{#2} -\rotationdd{#1}{#2}\rotationInv{#1}{#2} \translation{#1}{#2}
      -\rotationd{#1}{#2} \rotationInvD{#1}{#2} \translation{#1}{#2} - \rotationd{#1}{#2}\rotationInv{#1}{#2}\translationd{#1}{#2}
    }{
      \zeroVector
    }{
      0
    }
  }}
\newcommand{\bodyAccDefDef}[2]{
  \textcolor{bodyVelColor}{
    \matrixTwo{
      \rotationInv{#1}{#2}\rotationdd{#1}{#2} + \rotationInvD{#1}{#2}\rotationd{#1}{#2}
    }{
      \rotationInvD{#1}{#2}\translationd{#1}{#2} + \rotationInv{#1}{#2}\translationdd{#1}{#2}
    }{
      \zeroVector
    }{
      0
    }
  }}
\newcommand{\bodyAccDef}[2]{
  \textcolor{bodyVelColor}{
    \matrixTwo{
      \bodyRATwist{#1}{#2}
    }{
      % \bodyTV{#1}{#2} + \rotationInv{#1}{#2}\translationdd{#1}{#2}
      \bodyTA{#1}{#2}
    }{
      \zeroVector
    }{
      0
    }
  }}


\newcommand{\spatialRA}[2]{
  \textcolor{spatialVelColor}{
    \dot{\angularVel}^{\spatialFrame}_{#1#2}
  }}
\newcommand{\spatialRATwist}[2]{
  \textcolor{spatialVelColor}{
    \skewMatrix{\dot{\angularVel}}^{\spatialFrame}_{#1#2}
  }}
\newcommand{\bodyRATwist}[2]{
  \textcolor{bodyVelColor}{
    \skewMatrix{\dot{\angularVel}}^{\bodyFrame}_{#1#2}
  }}

\newcommand{\spatialRV}[2]{
  \textcolor{spatialVelColor}{
    \angularVel^{\spatialFrame}_{#1#2}
  }}

\newcommand{\spatialRVDef}[2]{
  \textcolor{spatialVelColor}{
    (\rotationd{#1}{#2}\rotationInv{#1}{#2})^{\vee}
}}
\newcommand{\spatialRVDefTwo}[2]{
  \textcolor{spatialVelColor}{
    \rotationd{#1}{#2}\rotationInv{#1}{#2}
}}

\newcommand{\spatialRVd}[2]{
  \textcolor{spatialVelColor}{
    \dot{\angularVel}^{\spatialFrame}_{#1#2}
}}

\newcommand{\spatialTVd}[2]{
  \textcolor{spatialVelColor}{
    \dot{\translationVel}^{\spatialFrame}_{#1#2}
}}

\newcommand{\spatialVelTwist}[2]{
  \textcolor{spatialVelColor}{
    \skewMatrix{\vc{V}}^{\spatialFrame}_{#1#2}
  }}

\newcommand{\spatialVelTwistDef}[2]{
  \textcolor{spatialVelColor}{
    \matrixTwo
    {
      \spatialRVTwist{#1}{#2}
    }{
      \spatialTV{#1}{#2}
    }{
      \zeroVector
    }{
      0
    }
}}
\newcommand{\spatialTVTwist}[2]{
  \textcolor{spatialVelColor}{
    \skewMatrix{\vc{v}}^{\spatialFrame}_{#1#2}
}}
\newcommand{\spatialRVTwist}[2]{
  \textcolor{spatialVelColor}{
    \skewMatrix{\vc{w}}^{\spatialFrame}_{#1#2}
}}
\newcommand{\spatialRVdTwist}[2]{
  \textcolor{spatialVelColor}{
    \skewMatrix{\dot{\vc{w}}}^{\spatialFrame}_{#1#2}
  }}

%%% #1 Representation frame, #2 Base frame, and #3 Body frame
\newcommand{\vel}[3]{
  \textcolor{velColor}{
    \vc{V}^{#1}_{#2#3}
  }}
\newcommand{\velJacobian}[3]{
  \textcolor{velColor}{
    \jacobian^{#1}_{#2#3}
  }}
\newcommand{\tV}[3]{
  \textcolor{velColor}{
    \vc{v}^{#1}_{#2#3}
  }}

\newcommand{\rV}[3]{
  \textcolor{velColor}{
    \vc{w}^{#1}_{#2#3}
  }}

%%% #1 Reference frame #2 Body frame
\newcommand{\bodyVel}[2]{
  \textcolor{bodyVelColor}{
    \vc{V}^{\bodyFrame}_{#1#2}
  }}
%%% #1 Reference frame #2 Body frame
\newcommand{\bodyTVDef}[2]{
  \textcolor{bodyVelColor}{
    \rotationInv{#1}{#2}\translationd{#1}{#2}
  }}
%%% #1 Reference frame #2 Body frame
\newcommand{\bodyRVDef}[2]{
  \textcolor{bodyVelColor}{
      (\rotationInv{#1}{#2}\rotationd{#1}{#2})^{\vee}
  }}


\newcommand{\bodyRVDefTwo}[2]{
  \textcolor{bodyVelColor}{
    \rotationInv{#1}{#2}\rotationd{#1}{#2}
  }}



%%% #1 Reference frame #2 Body frame
\newcommand{\bodyVelDef}[2]{
  \textcolor{bodyVelColor}{
    \vectorTwo{
      \bodyTVDef{#1}{#2}
    }{
      \bodyRVDef{#1}{#2}
    }
  }}

%%% #1:Reference frame of the screw, #2:Body frame attached to the screw motion, #3: angle
\newcommand{\screwBodyAcc}[3]{
  \textcolor{bodyVelColor}{
    \Vee{\twistTransform{#1}{#2(0)}  \twistC{#1}{#2}}\ddot{#3}
  }}
%%% #1:Reference frame of the screw, #2:Body frame attached to the screw motion, #3: angle
\newcommand{\screwBodyVel}[3]{
  \textcolor{bodyVelColor}{
    \Vee{\twistTransform{#1}{#2(0)}  \twistC{#1}{#2}}\dot{#3}
  }}
%%% #1:Reference frame of the screw, #2:Body frame attached to the screw motion, #3: angle
\newcommand{\screwSpatialAcc}[3]{
  \textcolor{spatialVelColor}{
    \twistC{#1}{#2}\ddot{#3}
  }}

%%% #1:Reference frame of the screw, #2:Body frame attached to the screw motion, #3: angle
\newcommand{\screwSpatialVel}[3]{
  \textcolor{spatialVelColor}{
    \twistC{#1}{#2}\dot{#3}
  }}



\newcommand{\bodyVelTwist}[2]{
  \textcolor{bodyVelColor}{
    \skewMatrix{\vc{V}}^{\bodyFrame}_{#1#2}
  }}
\newcommand{\bodyVelTwistDef}[2]{
  \textcolor{bodyVelColor}{
    \matrixTwo
    {
      \bodyRVTwist{#1}{#2}
    }{
      \bodyTV{#1}{#2}
    }{
      \zeroVector
    }{
      0
    }
  }}
\newcommand{\bodyVelTwistDefDef}[2]{
  \textcolor{bodyVelColor}{
    \matrixTwo
    {
      \rotationInv{#1}{#2}\rotationd{#1}{#2}
    }{
      \bodyTVDef{#1}{#2}
    }{
      \zeroVector
    }{
      0
    }
}}
\newcommand{\bodyTV}[2]{
  \textcolor{bodyVelColor}{
    \vc{v}^{\bodyFrame}_{#1#2}
  }}
\newcommand{\bodyTA}[2]{
  \textcolor{bodyVelColor}{
    \dot{\vc{v}}^{\bodyFrame}_{#1#2}
  }}
\newcommand{\bodyTADef}[2]{
  \textcolor{bodyVelColor}{
    % \rotationInvD{#1}{#2}\translationd{#1}{#2} + \rotationInv{#1}{#2}\translationdd{#1}{#2}
    \rotationInv{#1}{#2}\translationdd{#1}{#2} - \bodyRVTwist{#1}{#2}\bodyTV{#1}{#2}
  }}
\newcommand{\bodyTADefDef}[2]{
  \textcolor{bodyVelColor}{
    \rotationInvD{#1}{#2}\translationd{#1}{#2} + \rotationInv{#1}{#2}\translationdd{#1}{#2}
  }}


\newcommand{\bodyRA}[2]{
  \textcolor{bodyVelColor}{
    \dot{\angularVel}^{\bodyFrame}_{#1#2}
  }}
\newcommand{\bodyTVd}[2]{
  \textcolor{bodyVelColor}{
    \dot{\vc{v}}^{\bodyFrame}_{#1#2}
}}
\newcommand{\bodyTVTwist}[2]{
  \textcolor{bodyVelColor}{
    \skewMatrix{\vc{v}}^{\bodyFrame}_{#1#2}
}}

\newcommand{\bodyRV}[2]{
  \textcolor{bodyVelColor}{
    \vc{w}^{\bodyFrame}_{#1#2}
}}
\newcommand{\bodyRVd}[2]{
  \textcolor{geometricColorVariable}{
    \dot{\vc{w}}^{\bodyFrame}_{#1#2}
}}

\newcommand{\bodyForce}{
\force^{\bodyFrame}
}
\newcommand{\bodyTorque}{
\torque^{\bodyFrame}
}

\newcommand{\bodyRVTwist}[2]{
  \textcolor{bodyVelColor}{
    \skewMatrix{\vc{w}}^{\bodyFrame}_{#1#2}
}}



\newcommand{\twistCTransform}[2]{
  \textcolor{geometricColorVariable}{
    \twistTransform{#1}{#2}
  }
}
\newcommand{\twistCTransformDef}[2]{
  \textcolor{geometricColorVariable}{
    \twistTransformDef{#1}{#2}
  }
}
% transform a point from coordinate #1 to coordinate #2
\newcommand{\pointTransform}[2]{
  \textcolor{geometricColorVariable}{
    \transformInv{#1}{#2}
  }
}
\newcommand{\pointTransformDef}[2]{
  \textcolor{geometricColorVariable}{
    \transformInvDef{#1}{#2}
  }
}
\newcommand{\pointTransformTwo}[2]{
  \textcolor{geometricColorVariable}{
    \transform{#2}{#1}
  }
}
\newcommand{\pointTransformTwoDef}[2]{
  \textcolor{geometricColorVariable}{
    \transformDef{#2}{#1}
  }
}
% transform from coordinate #1 to coordinate #2
\newcommand{\twistTransformTwo}[2]{
  \textcolor{geometricColorVariable}{
    \adg{#2}{#1}
  }
}

\newcommand{\twistTransformdTwo}[2]{
  \textcolor{geometricColorVariable}{
    \derivative{\adg{#2}{#1}}
  }
}
\newcommand{\twistTransformdTwoDef}[2]{
  \textcolor{geometricColorVariable}{
    -\adg{#2}{#1}(\spatialVel{#1}{#2}\times)
  }
}

\newcommand{\twistTransformTwoDef}[2]{
  \textcolor{geometricColorVariable}{
    \adgDef{#2}{#1}
  }
}

\newcommand{\twistTransform}[2]{
  \textcolor{geometricColorVariable}{
    \adgInv{#1}{#2}
  }
}


% The direct transform: re-write the coordinates.
% #1 current frame, #2 target frame.
\newcommand{\dTransform}[2]{
  \textcolor{geometricColorVariable}{
    \matrixTwo{\tdTransform{#1}{#2}}{0}{0}{\tdTransform{#1}{#2}}
  }
}

% #1 current frame, #2 target frame.
\newcommand{\tdTransform}[2]{
  \textcolor{geometricColorVariable}{
    \rotationInv{#1}{#2}
  }
}


\newcommand{\twistTransformDef}[2]{
  \adgInvDef{#1}{#2}
}
\newcommand{\twistTransformd}[2]{
  \textcolor{geometricColorVariable}{
    \derivative{\twistTransform{#1}{#2}}
    }
}
\newcommand{\twistTransformdDef}[2]{
  \textcolor{geometricColorVariable}{
    -\twistTransform{#1}{#2}(\spatialVel{#1}{#2}\times)
    }
  }

\newcommand{\twistTransformdDefTwo}[2]{
  \textcolor{geometricColorVariable}{
    \twistTransform{#1}{#2}(\bodyVel{#2}{#1}\times)
    }
}





% #1: reference frame , #2 link(body) frame, and #3: New reference frame.
\newcommand{\spatialVelTransform}[3]{
  \textcolor{spatialVelColor}{
    % \spatialVel{#3}{#1} + \spatialVelStaticTransform{#3}{#1}\spatialVel{#1}{#2}
    \spatialVel{#3}{#1} + \twistTransform{#1}{#3}\spatialVel{#1}{#2}
  }
}
% #1: reference frame , #2 link(body) frame, and #3: New reference frame.
\newcommand{\spatialAccTransform}[3]{
  \textcolor{spatialVelColor}{
    % \spatialVel{#3}{#1} + \spatialVelStaticTransform{#3}{#1}\spatialVel{#1}{#2}
    \spatialAcc{#3}{#1}
    +\twistTransform{#1}{#3} (\spatialVelTwist{#1}{#2}\spatialVelTwist{#1}{#3} - \spatialVelTwist{#1}{#3}\spatialVelTwist{#1}{#2})^\vee
    % +\twistTransform{#1}{#3} (\bodyVelTwist{#3}{#1}\bodyVelTwist{#2}{#1})^\vee
    + \twistTransform{#1}{#3}\spatialAcc{#1}{#2}
  }
}
% #1: reference frame , #2 link(body) frame, and #3: New reference frame.
\newcommand{\spatialAccTransformTwo}[3]{
  \textcolor{spatialVelColor}{
    \spatialAcc{#3}{#1}
    +\twistTransform{#1}{#3}(\cross{\spatialVel{#1}{#2}}{\spatialVel{#1}{#3}})
    + \twistTransform{#1}{#3}\spatialAcc{#1}{#2}
  }
}

% #1: reference frame , #2 link(body) frame, and #3: New reference frame.
\newcommand{\bodyAccTransform}[3]{
  \textcolor{bodyVelColor}{
    \bodyAcc{#1}{#2}
    + \twistTransform{#1}{#2}(
    \bodyVelTwist{#2}{#1}\bodyVelTwist{#3}{#1}
    -\bodyVelTwist{#3}{#1}\bodyVelTwist{#2}{#1}
    )^{\vee}
    +\twistTransform{#1}{#2}\bodyAcc{#3}{#1}
  }
}
% #1: reference frame , #2 link(body) frame, and #3: New reference frame.
\newcommand{\bodyAccTransformTwo}[3]{
  \textcolor{bodyVelColor}{
    \bodyAcc{#1}{#2}
    % \twistTransformdDef{#1}{#2}
    +\twistTransform{#1}{#2}
    (\cross{\bodyVel{#2}{#1}}
    {\bodyVel{#3}{#1}})
    +\twistTransform{#1}{#2}\bodyAcc{#3}{#1}
  }
}

\newcommand{\spatialVelTransformTwo}[3]{
  \textcolor{spatialVelColor}{
    % \spatialVel{#3}{#1} + \spatialVelStaticTransform{#3}{#1}\spatialVel{#1}{#2}
    \spatialVel{#3}{#1} +
    \twistTransformTwo{#1}{#3}\spatialVel{#1}{#2}
  }
}

% Expressed in #2
% #1: reference frame , #2 link(body) frame, and #3: New reference frame.
\newcommand{\bodyVelTransform}[3]{
  \textcolor{bodyVelColor}{
    \twistTransform{#1}{#2}\bodyVel{#3}{#1} + \bodyVel{#1}{#2}
  }
}
\newcommand{\bodyVelTransformTwo}[3]{
  \textcolor{bodyVelColor}{
    \twistTransformTwo{#1}{#2}\bodyVel{#3}{#1} + \bodyVel{#1}{#2}
  }
}


% Transform from #1 to #2
\newcommand{\bodyWrench}{\bs{W}^{\bodyFrame}}



%%%%%%%%%%%%%%%%%%%%%%%%% Inertia and momentum:
%   #1 reference frame name, #2 body frame name
%  locally, we can specify: #1 body frame: F_i  #2 the local centroidal frame: F_com_i

%%%%%%%%%%%%%%%%%%%%%%%%%%%%%%%%%%%%%%%%%%%%%%%%%%%%%%%%%%%%%%%%%%%%%%%%%%%%%%%%
%% 10. SPATIAL VECTOR ALGEBRA (Featherstone)
%%%%%%%%%%%%%%%%%%%%%%%%%%%%%%%%%%%%%%%%%%%%%%%%%%%%%%%%%%%%%%%%%%%%%%%%%%%%%%%%

\newcommand{\svaMT}[2]{
  \svaFT{#1}{#2}
}


\newcommand{\svaTransform}[2]{
  \textcolor{svaColorVariable}{
    {^{#2}T_{#1}}
  }
}
% Transform from #1 to #2
\newcommand{\svaTransformDef}[2]{
  \textcolor{svaColorVariable}{
  \matrixTwo{
    \rotation{#2}{#1}
  }{
    -\rotation{#2}{#1}\translation{#1}{#2}
  }{
    \zeroVector
  }{
    1
  }
  }
}
\newcommand{\svaTransformDefTwo}[2]{
  \textcolor{svaColorVariable}{
  \matrixTwo{
    \rotationInv{#1}{#2}
  }{
    -\rotationInv{#1}{#2}\translation{#1}{#2}
  }{
    \zeroVector
  }{
    1
  }
  }
}

\newcommand{\svaT}{
  \textcolor{svaColorVariable}{
    \bs{X}
  }
}

% Transform from #1 to #2 (#2 is represented in #1)
\newcommand{\svaVT}[2]{
  \textcolor{svaColorVariable}{
  {^{#2}\svaT_{#1}}
  % \svaT{#1}{#2}
  }
}

% Transform Vel from #1 to #2
\newcommand{\svaVTDef}[2]{
  % eq.~2.24 of FeatherStone
  % \matrixTwo{
  %   \rotationInv{#1}{#2}
  % }{\zeroMatrix}{
  %   -\rotationInv{#1}{#2}\translationSkew{#1}{#2}
  % }
  % {
  %   \rotationInv{#1}{#2}
  % }
  \textcolor{svaColorVariable}{
  \matrixTwo{
    \rotation{#2}{#1}
  }{\zeroMatrix}{
    -\rotation{#2}{#1}\translationSkew{#1}{#2}
  }
  {
    \rotation{#2}{#1}
  }
  }
}

% Transform Vel from #1 to #2
\newcommand{\svaVTDefTwo}[2]{
  % eq.~2.26 of FeatherStone
  \textcolor{svaColorVariable}{
  \matrixTwo{
    \rotationInv{#1}{#2}
  }{\zeroMatrix
  }{
    \translationSkew{#2}{#1}\rotationInv{#1}{#2}
  }
  {
    \rotationInv{#1}{#2}
  }
  }
}

% Transform from #1 to #2
\newcommand{\svaFT}[2]{
  \textcolor{svaColorVariable}{
    {^{#2}\svaT^{*}_{#1}}
  }
}

% Transform from #1 to #2
\newcommand{\svaFTDef}[2]{
  % Following eq.2.27 of Featherstone
  % \matrixTwo
  % {\rotationInv{#2}{#1}}
  % {
  %   \translationSkew{#1}{#2}\rotationInv{#2}{#1}
  % }{
  %   \zeroMatrix
  % }
  %   {\rotationInv{#2}{#1}}
  \textcolor{svaColorVariable}{
  \matrixTwo
  {\rotationInv{#1}{#2}}
  {
    \translationSkew{#2}{#1}\rotationInv{#1}{#2}
  }{
    \zeroMatrix
  }
  {\rotationInv{#1}{#2}}
  }
}

% Transform from #1 to #2
\newcommand{\svaFTDefTwo}[2]{
  % Following eq.2.25 of Featherstone
  % \matrixTwo{
  %   \rotationInv{#1}{#2}
  % }{
  %   -\rotationInv{#1}{#2}\translationSkew{#1}{#2}
  % }{
  %   \zeroMatrix
  % }
  %   {\rotationInv{#1}{#2}}
  \textcolor{svaColorVariable}{
  \matrixTwo{
    \rotation{#2}{#1}
  }{
    -\rotation{#2}{#1}\translationSkew{#1}{#2}
  }{
    \zeroMatrix
  }
  {\rotation{#2}{#1}}
  }
}

% Reference frame: #1, link frame: #2
\newcommand{\svaBodyGInertiaDef}[2]{
  \textcolor{svaColorVariable}{
    \matrixTwo{\bodyMetaInertia{#1}{#2} }{\zeroMatrix}{\zeroMatrix}{\mass_{#2} \identityMatrix}
  }
}


% #1 Reference frame #2 body frame
\newcommand{\svaMomentaTransform}[2]{
  \transpose{\svaVT{#2}{#1}}
}
\newcommand{\svaMomentaTransformDef}[2]{
  \transpose{\svaVTDef{#2}{#1}}
}


% #1 and #2 defines the spatial generalized inertia to be transformed.
% #1 Reference frame #2 body frame
% #3 defines the new reference frame
\newcommand{\svaSpatialGInertiaTransform}[3]{
  \transpose{\svaVT{#3}{#1}}\spatialGInertia{#1}{#2}\svaVT{#3}{#1}
}

\newcommand{\svaSpatialCrbInertiaTransform}[3]{
  \transpose{\svaVT{#3}{#1}}\spatialCrbInertia{#1}{#2}\svaVT{#3}{#1}
}

%%% Generalized inertia transform
% #1:new reference frame
% #2:current reference frame
% #3:current Generalized inertia
\newcommand{\svaGInertiaTransform}[3]{
  \transpose{\svaVT{#1}{#2}}
  #3
  \svaVT{#1}{#2}
}


% #1 Reference frame #2 body frame
\newcommand{\svaSpatialGInertia}[2]{
  \textcolor{svaColorVariable}{
    \transpose{\svaVT{#1}{#2}} \bodyGInertia{#1}{#2}\svaVT{#1}{#2}
  }
}

% #1 Reference frame #2 body frame
\newcommand{\svaSpatialGInertiaDef}[2]{
  \textcolor{svaColorVariable}{
    \transpose{\svaVTDef{#1}{#2}} \svaBodyGInertiaDef{#1}{#2}\svaVTDef{#1}{#2}
  }
}



%%%%%%%%%%%%%%%%%%%%%%%%%%%%%%%%%%%%%%%%%%%%%%%%%%%%%%%%%%%%%%%%%%%%%%%%%%%%%%%%
%% 11. JOINT KINEMATICS
%%%%%%%%%%%%%%%%%%%%%%%%%%%%%%%%%%%%%%%%%%%%%%%%%%%%%%%%%%%%%%%%%%%%%%%%%%%%%%%%

\newcommand{\svaJointJac}[1]{\bf{\Phi}_{#1}}


\newcommand{\applyJangles}[1]{%
    \IfStrEqCase{#1}{%
    {actuated}{
    \bs{\theta}
    }%
    {fb}{
    \bs{q}_{\text{fb}}
    }%
    {generalized}{
    \bs{\mathcal{X}}
    }%
    }[\bs{q}] % "default" string
}


\newcommand{\applyJanglesPostscript}[1]{%
    \IfStrEqCase{#1}{%
    {current}{
    {t_k}
    }%
    {next}{
    {t_{k+1}}
    }%
    {previous}{
    {t_{k-1}}
    }%
    }[] % "default" string
}


\NewDocumentCommand{\jangles}{O{}O{}O{}O{}}{%
{    \ensuremath{%
        % Check the first optional argument for dot or double dot
        \applyOverheadSymbols{#1}{\applyJangles{#3}} % Third argument for actuated, fb joints
        % 'ref', 'mea', 'des'
        ^{\applySuperscript{#2}}%
        _{\applyJanglesPostscript{#4}}
    }
    }
}

\newcommand{\jvelocities}{\mathbf{\dot{q}}}
\newcommand{\jaccelerations}{\mathbf{\ddot{q}}}
\newcommand{\jtorques}{\boldsymbol{\tau}}
\newcommand{\omegaVector}{\boldsymbol{\omega}}
\newcommand{\tsInertia}{\boldsymbol{\Lambda}}




\newcommand{\janglesJump}{\jump \jangles}
\newcommand{\jvelocitiesJump}{\jump \jvelocities}
\newcommand{\jangle}{q}

%%%%%%%%%%%%%%%%%%%%%%%%%%%%%%%%%%%%%%%%%%%%%%%%%%%%%%%%%%%%%%%%% Bound
\newcommand{\jvelocitiesUpperBound}{\mathbf{\dot{\bar{\jangle}}}}
\newcommand{\jvelocitiesLowerBound}{\mathbf{\dot{\ubar{\jangle}}}}
\newcommand{\jvelocitiesJumpUpperBound}{\jump \jvelocitiesUpperBound}
\newcommand{\jvelocitiesJumpLowerBound}{\jump \jvelocitiesLowerBound}

%%%%%%%%%%%%%%%%%%%%%%%%%%%%%%%%%%%%%%%%%%%%%%%%%%%%%%%%%%%%%%%%%%%



%%%%%%%%%%%%%%%%%%%%%%%%%% Def equal
\newcommand{\jointVel}[2]{\twistVel^{#1}_{J#2}}%
\newcommand{\jointVelDef}[2]{\spatialVel{#1}{#2}}%
\newcommand{\jointAcc}[2]{\dot{\twistVel}^{#1}_{J#2}}%
\newcommand{\jointAccDef}[2]{\spatialAcc{#1}{#2}}%


% In Body coordinates: #1 Joint index
% *predecessor idx: #1 -1  *successor #1
\newcommand{\jointVelBody}[1]{
  \textcolor{bodyVelColor}{
    \twistVel_{J#1}
  }
}
\newcommand{\jointVelBodynew}[2]{
  \textcolor{bodyVelColor}{
    \twistVel^{#1}_{J#2}
  }
}
% Defines the Joint frame number. The joint locates on the link #1-1 and does not depend on the #1th joint position
% #1 Joint index
\newcommand{\jointFrameDef}[1]{
  (#1 -1, #1)
}
% In Body coordinates: #1 Joint index
% *predecessor idx: #1 -1  *successor #1
\newcommand{\jointVelBodyDef}[1]{
  \bodyVel{\jointFrameDef{#1}}{#1}
  % \bodyVel{(#1-1, #1)}{#1}
}
% In Body coordinates: #1 Reference Frame  #2 Joint index
% *predecessor idx: #2 -1  *successor #2
% Due to the rigid link assumption, this formula assumes \bodyVel{(#2-1)}{(#2-1, #2)} = 0.
\newcommand{\jointVelBodyDefDef}[2]{
  \bodyVel{#1}{#2} - \twistTransform{(#2-1)}{#2}\bodyVel{#1}{(#2-1)}
}

% In Body coordinates: #1 Joint index
% *predecessor idx: #1 -1  *successor #1
\newcommand{\jointAccBody}[1]{
  \textcolor{bodyVelColor}{
    \dot{\twistVel}_{J#1}
  }
}
% In Body coordinates: #1 Joint index
% *predecessor idx: #1 -1  *successor #1
\newcommand{\jointAccBodyDef}[2]{
  \bodyAcc{\jointFrameDef{#1}}{#1}
}


% #1 Spatial frame or body frame, #2 joint index.
\newcommand{\linkVel}[2]{\twistVel^{#1}_{L#2}}%
\newcommand{\linkVelBody}[2]{
  \textcolor{bodyVelColor}{
    \twistVel^{#1}_{L#2}
  }
}%
\newcommand{\linkAcc}[2]{\dot{\twistVel}^{#1}_{L#2}}%

% The force transmits across a joint.
% #1 Reference frame, #2 joint index.
\newcommand{\body}[1]{#1^{b}}
\newcommand{\spatial}[1]{#1^{s}}

% Joint positions are unified via \jangles[overhead][superscript][type][postscript]:
%   \jangles             -> \bs{q}              (default, joint position)
%   \jangles[dot]        -> \dot{\bs{q}}
%   \jangles[ddot]       -> \ddot{\bs{q}}
%   \jangles[][][actuated]    -> \bs{\theta}    (actuated joint angles)
%   \jangles[][][generalized] -> \bs{\mathcal{X}} (floating-base generalized coordinates)
%   \jangles[][][fb]     -> \bs{q}_{\text{fb}}  (floating-base joint position)
% Simple symbol aliases for use with subscripts in equations:
\newcommand{\jointPosition}{\bs{\theta}}
\newcommand{\jointPositiond}{\dot{\jointPosition}}
\newcommand{\jointPositiondd}{\ddot{\jointPosition}}

% ################## Coordinate frames

% The frame defined by Christian Ott:
\newcommand{\jointAngle}{\theta}
\newcommand{\jointAngled}{\dot{\jointAngle}}
% \jangles[][][actuated] is now unified as \jangles[][][actuated]

%%%%%%%%%%%%%%%%%%%%%%%%%%%%%%%%%%%%%%%%%%%%%%%%%%%%%%%%%%%%%%%%%%%%%%%%%%%%%%%%
%% 12. JACOBIANS
%%%%%%%%%%%%%%%%%%%%%%%%%%%%%%%%%%%%%%%%%%%%%%%%%%%%%%%%%%%%%%%%%%%%%%%%%%%%%%%%

\newcommand{\applyJacobianSubscript}[1]{%
    \IfStrEqCase{#1}{%
    {contact}{
    \contactPoint
    }%
    }[#1] % "default" string
}

\newcommand{\jacobianSymbol}{\text{J}}
\NewDocumentCommand{\jacobian}{O{}O{}}{%
{    \ensuremath{%
        % Check the first optional argument for dot or double dot
        \applyOverheadSymbols{#1}{\jacobianSymbol}
        % 'contact',
        _{\applyJacobianSubscript{#2}}%
    }
    }
}



\newcommand{\jacobianVector}{\bs{\jacobian}}

\newcommand{\jac}[1]{\jacobian_{#1}} % Jacobian of a special quantity
\newcommand{\jacd}[1]{\dot{\jacobian}_{#1}} % Jacobian of a special quantity
\newcommand{\jacdd}[1]{\ddot{\jacobian}_{#1}} % Jacobian of a special quantity

\newcommand{\jacobianMotion}{\jacobian_{\text{m}}}%


%%%%%%%%%%%%%%%% The articulated-body inertia
% Matrix spans a subspace
\newcommand{\rVelJac}[2]{{
  \textcolor{bodyVelColor}{
    \jacobian^{'\bodyFrame}_{#1#2}
  }
}}

% The relative velocity Jacobian: vel = vel(induced by cmm vel) + relative vel
% #1 inertial frame #2 body frame
\newcommand{\rVel}[2]{{
  \textcolor{bodyVelColor}{
    \vc{V}^{'\bodyFrame}_{#1#2}
  }
}}
% #1 inertial frame #2 body frame
\newcommand{\rVelJacDef}[2]{{
    (\bodyJacobian{#1}{#2} - \twistTransform{\cmmFrame}{#2}\bodyJacobian{#1}{\cmmFrame})
}}



\newcommand{\appBodyJacobian}[2]{
  \textcolor{bodyVelColor}{
    {\tilde{{\jacobian}}^{\bodyFrame\text{t}}_{#1 #2}}
  }
}
\newcommand{\tBodyJacobian}[2]{
  \textcolor{bodyVelColor}{
    {\jacobian^{\bodyFrame\text{t}}_{#1 #2}}
  }
}
\newcommand{\rBodyJacobian}[2]{
  \textcolor{bodyVelColor}{
    {\jacobian^{\bodyFrame\text{R}}_{#1 #2}}
  }
}

%  Define a macro with an optional argument for dot, or more.
\newcommand{\bodyJacobian}[3][]{
  \textcolor{bodyVelColor}{
      #1{\jacobian}^{\bodyFrame}_{#2 #3}
  }
}

\newcommand{\bodyJacobiand}[2]{
  \textcolor{bodyVelColor}{
    \dot{\jacobian}^{\bodyFrame}_{#1 #2}
  }
}


%%%%%%%%%%%%%%%%%%%%%%%%%% QP Control law


%%%%%%%%%%%%%%%%%%%%%%%%%%%%%%%%%%%%%%%%%%%%%%%%%%%%%%%%%%%%%%%%%%%%%%%%%%%%%%%%
%% 13. INERTIA, MASS, AND MOMENTUM
%%%%%%%%%%%%%%%%%%%%%%%%%%%%%%%%%%%%%%%%%%%%%%%%%%%%%%%%%%%%%%%%%%%%%%%%%%%%%%%%

\newcommand{\smm}{
  \textcolor{bodyVelColor}{
    \bs{A}
  }
}

\newcommand{\systemMomenta}{
  \textcolor{bodyVelColor}{
    \bs{h}
  }
}
\newcommand{\momenta}{
  \bs{h}
}

% The momentum transform looks the same as the wrench transform
%%% #1:current frame, #2:target frame
\newcommand{\geometricMT}[2]{
  \geometricFT{#1}{#2}
}
\newcommand{\geometricMTTwo}[2]{
  \geometricFTTwo{#1}{#2}
}
\newcommand{\systemVel}{
  \textcolor{bodyVelColor}{
    \bs{V}
  }
}
\newcommand{\systemJacobian}{
  \textcolor{bodyVelColor}{
   \jacobian
  }
}
\newcommand{\systemJacobiand}{
  \textcolor{bodyVelColor}{
   \jacobian[dot]
  }
}

\newcommand{\momentum}{\bs{h}}
\newcommand{\momentumd}{\dot{\bs{h}}}

%%% #1 reference frame, #2 link index
\newcommand{\momentumFrame}[2]{\bs{h}_{#1#2}}

%%% The canonical momenta:
\newcommand{\canoMomenta}{{\bs{h}_{J}}}

%%% The system momentum per link
\newcommand{\sMomentum}{{\bs{h}_{s}}}

%%% The contact momentum per contact
\newcommand{\cMomentum}{{\bs{h}_{c}}}

% #1:reference frame, #2:contact index
\newcommand{\cMomentumDef}[2]{
  \equivalentMass_{#2} \spatialVel{#1}{#2}
}

\newcommand{\cMomentumDefDef}[2]{
  {\equivalentMassDef{#1}{#2}} \spatialVel{#1}{#2}
}


\newcommand{\cmmMomenta}{
  \textcolor{bodyVelColor}{
  \bs{h}_{\cmmFrame}
 }
}
\newcommand{\cmmInertia}{
  \textcolor{bodyVelColor}{
    \bs{I}_{\cmmFrame}
  }
}

\newcommand{\cmmInertiaDef}{
  \agg{i}{\dof}{\spatialGInertia{\cmmFrame}{i}}
}
\newcommand{\cmmRelativeInertia}{
  \textcolor{bodyVelColor}{
    \bs{I}'_{\cmmFrame}
  }
}
\newcommand{\cmmRelativeInertiaDef}{
  \textcolor{bodyVelColor}{
    \quadraticObj{\rVelJac{}{}}{\systemInertia}
  }
}

\newcommand{\cmmInertiaDefTwo}{
  \agg{i}{\dof}{\spatialGInertia{\cmmFrame}{i}}
}

\newcommand{\cmmInertiaDetail}{
  \textcolor{bodyVelColor}{
    \bs{I}_{\cmmFrame}
  }
}

\newcommand{\linearMomentumSymbol}{P}
\NewDocumentCommand{\linearMomentum}{O{}O{}O{}}{%
{    \ensuremath{%
        % Check the first optional argument for dot or double dot
        \applyOverheadSymbols{#1}{\linearMomentumSymbol}
        % Check the second optional argument for 'ref'
        ^{\applySuperscript{#2}}
        \applyPostscript{#3}
    }
    }
}
%  The angular Momentum
\newcommand{\angularMomentumSymbol}{\textcolor{textColor}{\mathcal{L}}}
\NewDocumentCommand{\angularMomentum}{O{}O{}O{}O{}}{%
{    \ensuremath{%
        % Check the first optional argument for dot or double dot
        \applyOverheadSymbols{#1}{\angularMomentumSymbol}
        % Check the second optional argument for 'ref'
        _{#2}
        ^{\applySuperscript{#3}}
        \applyPostscript{#4}
    }
    }
}



%  The linear Momentum
\newcommand{\aMomentum}{\mathcal{L}}

\newcommand{\cmmSymbol}{
\textcolor{bodyVelColor}{
A
}
}




\newcommand{\applyCmmSuperscript}[1]{%
\textcolor{bodyVelColor}{
    \IfStrEqCase{#1}{%
    {l}{
    \text{t}
    }%
    {r}{
    \text{r}
    }%
    }[] % "default" string
    }
}


% %%%%%%%%%%%%%%%%%%%%%%%%% EQUATION VARIABLES:
\NewDocumentCommand{\cmm}{O{}O{default}O{default}O{default}}{%
{    \ensuremath{%
        % Check the first optional argument for dot or double dot
        \applyOverheadSymbols{#1}
        {\cmmSymbol}_{\textcolor{bodyVelColor}{\cmmFrame}}
        % Check the second optional argument for linear or angular
        ^{\applyCmmSuperscript{#2}}%
         \applyPostscript{#3}
    }
    }
}


\newcommand{\cmmt}{
  \textcolor{bodyVelColor}{
    A^{\text{t}}_{\cmmFrame}(\jangles)
  }
}


\newcommand{\cmmr}{
  \textcolor{bodyVelColor}{
    A^{\text{r}}_{\cmmFrame}(\jangles)
  }
}

\newcommand{\cmmLinear}{A_{\bs{v}\cmmFrame}}

\newcommand{\cmmDot}{
  \textcolor{bodyVelColor}{
  \dot{A}_{\cmmFrame}(\jangles[][][actuated])
}
}

\newcommand{\cmmPlane}{
  \textcolor{bodyVelColor}{
    A_{\cmmFrame_{x,y}}
  }
}

\newcommand{\aii}{\bs{\Phi}}
\newcommand{\aiiDef}[2]{#1 \inverse{(\transpose{#1} #2 #1)}\transpose{#1}}
% The joint acc of a rigid body constrianed by a joint.
% #1 the joint jacobian, #2 the inertia
\newcommand{\aInertia}{{I^A}}%

%%%%%%%%%%%%%%%% Motion constraint:

% The implicit constraint: \Phi(\dot{q}) = 0
% #1 Joint position
\newcommand{\fbVel}{
  \bs{V}_{\fbFrame}
}
\newcommand{\fbAcc}{
  \dot{\bs{V}}_{\fbFrame}
}

\newcommand{\cmmVel}{
  \textcolor{bodyVelColor}{
    \bs{V}_{\cmmFrame}
  }
}

\newcommand{\cmmVelDef}{
  \bodyVel{\inertialFrame}{\cmmFrame}
}
\newcommand{\cmmVelDefDef}{
  \inverse{\cmmInertia} \cmmMomenta
}
%%%%%%%%%%%% Exponential Map given unit skew matrix and a real number
%%% #1:input matrix
\newcommand{\crbVel}{
  {^{crb}\bs{V}}
}

%%% Spatial CRB inertia
% #1 inertial frame #2 body frame
\newcommand{\spatialCrbInertia}[2]{
  {^{crb}I^s_{#1#2}}
}

%%% CRB inertia
\newcommand{\crbGInertia}{
  {^{crb}I}
}

%%% System inertia
\newcommand{\systemInertia}{
  \textcolor{bodyVelColor}{
    I
  }
}


%%%%%%%%%%%

\newcommand{\averageVel}{
  \bar{\systemVel}
}
\newcommand{\relativeVel}{{\systemVel'}}
% The relative velocity Jacobian: vel = vel(induced by cmm vel) + relative vel
% #1 inertial frame #2 body frame
\newcommand{\equivalentMass}{{M_{eq}}}
% The 6-by-6 equivalent inertensor
\newcommand{\eInertiaTensor}{{I_{eq}}}

\newcommand{\osInertia}{{\Lambda}}
% The operational space bias force
\newcommand{\osBias}{{\bs{p}}}

% #1:reference frame #2:contact index
\newcommand{\equivalentMassDef}[2]{
  \inverse{\left(
      \spatialJacobian{#1}{#2} \inverse{\inertiaMatrix}_{#2}
     \transpose{{\spatialJacobian{#1}{#2}}}
   \right)}
}

\newcommand{\geometricFTDef}[2]{
  \textcolor{geometricColorVariable}{
    \adgInvTransDef{#2}{#1}
  }
}

% From frame #1 to frame #2
\newcommand{\geometricFTTwo}[2]{
  \textcolor{geometricColorVariable}{
    \adgTrans{#1}{#2}
  }
}
\newcommand{\geometricFTTwoDef}[2]{
  \textcolor{geometricColorVariable}{
    \adgTransDef{#1}{#2}
  }
}

% From frame #1 to frame #2
\newcommand{\geometricFT}[2]{
  \textcolor{geometricColorVariable}{
    \adgInvTrans{#2}{#1}
  }
}

% From frame #1 to frame #2
\newcommand{\geometricFTd}[2]{
  \textcolor{geometricColorVariable}{
    \derivative{\adgInvTrans{#2}{#1}}
  }
}

% From frame #1 to frame #2
\newcommand{\geometricFTdDef}[2]{
  \textcolor{geometricColorVariable}{
    -\transpose{\cross{\spatialVel{#2}{#1}}{}}\geometricFT{#1}{#2}
  }
}
\newcommand{\geometricFTdDefTwo}[2]{
  \textcolor{geometricColorVariable}{
    \dualCross{\spatialVel{#2}{#1}}{}\geometricFT{#1}{#2}
  }
}
\newcommand{\geometricFTdDefThree}[2]{
  \textcolor{geometricColorVariable}{
        \transpose{(\bodyVel{#1}{#2}\times)} \transpose{\twistTransform{#2}{#1}}
  }
}

\newcommand{\bodyGInertia}[2]{
\textcolor{bodyVelColor}{
  \body{\gInertia}_{#1 #2}
  }
}


\newcommand{\mechanicalConnection}{
  \bs{A}
}




\newcommand{\bodyGInertiaDef}[2]{
  \textcolor{bodyVelColor}{
    \matrixTwo{\mass_{#2} \identityMatrix }{\zeroMatrix}{\zeroMatrix}{\bodyMetaInertia{#1}{#2}}
  }
}

\newcommand{\spatialGInertia}[2]{
  \textcolor{spatialVelColor}{
    \spatial{\gInertia}_{#1#2}
  }
}

\newcommand{\spatialGInertiad}[2]{
  \textcolor{spatialVelColor}{
    \spatial{\dot{\gInertia}}_{#1#2}
  }
}

\newcommand{\spatialGInertiadDef}[2]{
  \textcolor{spatialVelColor}{
    % -\transpose{(\cross{\spatialVel{#1}{#2}}{})} \spatialGInertia{#1}{#2} - \spatialGInertia{#1}{#2}\cross{\spatialVel{#1}{#2}}{}
    \dualCrossDef{\spatialVel{#1}{#2}}{\spatialGInertia{#1}{#2}}
    - \spatialGInertia{#1}{#2}\cross{\spatialVel{#1}{#2}}{}
  }
}

\newcommand{\spatialGInertiadDefTwo}[2]{
  \textcolor{spatialVelColor}{
    \dualCross{\spatialVel{#1}{#2}}{\spatialGInertia{#1}{#2}}
    - \spatialGInertia{#1}{#2}\cross{\spatialVel{#1}{#2}}{}
  }
}

% #1 Current reference frame #2 New Reference frame
\newcommand{\momentaTransform}[2]{
  \transpose{\twistTransform{#2}{#1}}
}
\newcommand{\momentaTransformDef}[2]{
  \transpose{\twistTransformDef{#2}{#1}}
}

% #1 Current reference frame #2 New Reference frame
\newcommand{\spatialGInertiaTransform}[3]{
  \transpose{\twistTransform{#3}{#1}}\spatialGInertia{#1}{#2}\twistTransform{#3}{#1}
}

% #1 and #2 defines the spatial generalized inertia to be transformed.
% #1 Reference frame #2 body frame
% #3 defines the new reference frame
\newcommand{\spatialCrbInertiaTransform}[3]{
  \transpose{\twistTransform{#3}{#1}}\spatialCrbInertia{#1}{#2}\twistTransform{#3}{#1}
}

\newcommand{\gInertiaTransform}[3]{
  \transpose{\twistTransform{#1}{#2}}
  #3
  \twistTransform{#1}{#2}
}

\newcommand{\spatialGInertiaDef}[2]{
  \transpose{\twistTransform{#1}{#2}}
  % \bodyGInertia{#1}{#2}
  \spatialGInertia{#2}{\com_{#2}}
  \twistTransform{#1}{#2}
}

\newcommand{\spatialGInertiaDefTwo}[2]{
  \transpose{\twistTransformTwo{#1}{#2}}
  % \bodyGInertia{#1}{#2}
  \spatialGInertia{#2}{\com_{#2}}
  \twistTransformTwo{#1}{#2}
}
% #1 Reference frame #2 body frame
\newcommand{\spatialGInertiaDefDef}[2]{
  % \transpose{{\adgInvDef{#1}{#2}}} \bodyGInertiaDef{#2}\adgInvDef{#1}{#2}
  \transpose{{\twistTransformDef{#1}{#2}}}
  \spatialGInertia{#2}{\com_{#2}}
  % \bodyGInertiaDef{#1}{#2}
  \twistTransformDef{#1}{#2}
}
\newcommand{\spatialGInertiaDefDefTwo}[2]{
  % \transpose{{\adgInvDef{#1}{#2}}} \bodyGInertiaDef{#2}\adgInvDef{#1}{#2}
  \transpose{{\twistTransformTwoDef{#1}{#2}}}
  % \bodyGInertiaDef{#1}{#2}
  \spatialGInertia{#2}{\com_{#2}}
  \twistTransformTwoDef{#1}{#2}
}

% #1 Reference frame #2 body frame
\newcommand{\spatialGInertiaDetail}[2]{
  \textcolor{spatialVelColor}{
\matrixTwo{
  \mass_{#2} \identityMatrix
}{
  \mass_{#2}\transpose{\translationSkew{#1}{#2}}
}{
  \mass\translationSkew{#1}{#2}
}{
  \mass \translationSkew{#1}{#2} \transpose{\translationSkew{#1}{#2}} + \rotation{#1}{#2}\bodyMetaInertia{#1}{#2}\rotationInv{#1}{#2}
}
}
}
% #1 Reference frame #2 body frame
\newcommand{\spatialGInertiaDiagonalDetail}[2]{
  \textcolor{spatialVelColor}{
\matrixTwo{
  \mass_{#2} \identityMatrix
}{
  \zeroMatrix
}{
  \zeroMatrix
}{
  \rotation{#1}{#2}\bodyMetaInertia{#1}{#2}\rotationInv{#1}{#2}
}
}
}

% #1 Reference frame #2 body frame
\newcommand{\metaInertia}{\mathcal{I}}



% #1: the reference frame and #2 the link frame
\newcommand{\spatialMetaInertia}[2]
{
\textcolor{spatialVelColor}{
  \spatial{\metaInertia}_{{#1}{#2}}
}
}
% #1: the reference frame and #2 the link frame
\newcommand{\bodyMetaInertia}[2]
{
  \textcolor{bodyVelColor}{
      \body{\metaInertia}_{{#1 #2}}
    }
}
\newcommand{\spatialInertia}{\mathcal{I}^{\prime}}

\newcommand{\gInertia}{\mathcal{M}}

%%%%%%%%%%%%%%%%%%%%%%%%%%%%%%%%%%%%%%%%%%%%%%%%%%%%%%%%%%%%%%%%%%%%%%%%%%%%%%%%
%% 14. EQUATIONS OF MOTION
%%%%%%%%%%%%%%%%%%%%%%%%%%%%%%%%%%%%%%%%%%%%%%%%%%%%%%%%%%%%%%%%%%%%%%%%%%%%%%%%

\newcommand{\numArticulatedJoints}{n}
\newcommand{\nAJ}{n}
\newcommand{\ajDim}{\text{n}}
\newcommand{\actuationMatrix}{\text{B}}

\newcommand{\numLinks}{{n_{L}}}
\newcommand{\fbDim}{6}
\newcommand{\jsDim}{(\numArticulatedJoints+\fbDim)}
\newcommand{\selectionMatrix}{
  \textcolor{textColor}{
  \text{S}
  }
}

\newcommand{\subspace}{S}

% The apparent inverse inertia. #1 subspace #2 inertia
\newcommand{\jaccConstrained}[2]{
  \inverse{(\transpose{#1} #2 #1)}\transpose{#1}(\jaWrench{}{} - #2\dot{#1}\jvelocities - \biasForce )
}

% #1 Spatial frame, #2 joint index.
\newcommand{\cwBasis}{T}%
% The orthonormal basis of the Actuation wrench. T_a \in \RRm{6}{1}
\newcommand{\awBasis}{{T_a}}%
% The motion subspace:
\newcommand{\mBasis}{{S}}%


% The projection matrix that describes the connections between different links in a multi-body system
% The connectivity matrix
\newcommand{\pMatrix}{
  \textcolor{spatialVelColor}{
    P
  }
}%

\newcommand{\pMatrixBody}{
  \textcolor{bodyVelColor}{
    P
  }
}%


% The articulated-body inertia
\newcommand{\ic}[1]{\Phi(#1)}


% The constraint kernel: K \dot{q} = 0
\newcommand{\ick}{K}
\newcommand{\ickd}{\dot{\ick}}
%%% #1 Joint position
\newcommand{\ickDef}[1]{\pd{\ic{#1}}{#1}}

% The explicit constraint:  q = \gamma(x)
% #1 End-effector position
\newcommand{\ec}[1]{\gamma(#1)}

% The explicit constraint kernel:  \dot{q} = \eck\dot{x}
% #1 End-effector position
\newcommand{\eck}{G}
\newcommand{\eckDef}[1]{\pd{\ec{#1}}{#1}}

\newcommand{\eckd}{\dot{\eck}}


\newcommand{\cForce}{\bs{\lambda}}

\newcommand{\cTorque}{\jtorques_c}

% Dimension of the constriants
\newcommand{\cDim}{m}
\newcommand{\cRank}{m_{ic}}


%%%%%%%%%%%%%%%%


% Contact Vel:
\newcommand{\gravitySymbol}{\mathbf{g}}
\newcommand{\gravity}[1][default]{
    \IfStrEqCase{#1}{%
        {value}{9.81}%
        {vec}{\vectorThree{0}{0}{-1g}}%
    }[\gravitySymbol] % "default" string
}

% Operational space graviational forces
\newcommand{\sDim}{\text{n}}
% input dimension
\newcommand{\iDim}{\text{m}}


% First-order partial derivative #1 quantity, #2 variable
\newcommand{\gForce}{{\text{\bf g} }}
\newcommand{\gForceDef}{\vectorThree{0}{0}{-1}g}
\newcommand{\gValue}{\text{9.81}}
\newcommand{\gSymbol}{\text{g}}




\newcommand{\dof}{n}
\newcommand{\jsim}{\inertiaMatrix}
\newcommand{\inertiaMatrix}{\text{M}}
\newcommand{\inertiaMatrixFull}{\text{M}(\jangles)}
\newcommand{\gravityandcoriolis}{\mathbf{N}}
\newcommand{\gravityandcoriolisFull}{\gravityandcoriolis(\jangles, \jangles[dot])}
\newcommand{\coriolis}{\mathbf{C}}

\newcommand{\inertiaMatrixd}{\dot{\inertiaMatrix}}

%%% #1:row index, #2:column index
\newcommand{\inertiaMatrixComponentDef}[2]{
\sum^\dof_{l = max(#1,#2)}{
  \transpose{
    \twistC{S}{#1}(t_0)
  }\transpose{
    \twistTransform{#1}{l}
  }
  \spatialGInertia{S}{\com_l}(t_0)
  \twistTransform{#2}{l}\twistC{S}{#2}(t_0)
}
}

\newcommand{\inertiaMatrixComponentDefTwo}[2]{
\sum^\dof_{l = max(#1,#2)}{
  \transpose{
    \twistC{S}{#1}(t_0)
  }\transpose{
    \adg{l}{#1}
  }
  \spatialGInertia{S}{\com_l}(t_0)
  \adg{l}{#2}
  \twistC{S}{#2}(t_0)
}
}

%%%%%%%%%%%%%%%%%% inertiaMatrix partial derivative:
%%% #1:row index, #2:column index, #3 theta index
\newcommand{\inertiaMatrixdComponentDef}[3]{
  \sum^\dof_{l = max(#1,#2)}{\left(
      \transpose{
        \liebra{\adg{#3}{#1}\twistC{S}{#1}}{\twistC{S}{#3}}
      }
      \transpose{
        \adg{l}{#3}
      }
      \spatialGInertia{S}{\com_l}(t_0)
      \adg{l}{#2}
      \twistC{S}{#2}(t_0)
      +
      \transpose{
        \twistC{S}{#1}(t_0)
      }\transpose{
        \adg{l}{#1}
      }
      \spatialGInertia{S}{\com_l}(t_0)
      \adg{l}{#3}\liebra{\adg{#3}{#2}\twistC{S}{#2}}{\twistC{S}{#3}}.
      % \adg{l}{#2}
      % \twistC{S}{#2}(t_0)
    \right)
  }
}



\newcommand{\cMatrix}{C(\jointPosition, \jointPositiond)}
\newcommand{\nVector}{N}
\newcommand{\rpTransform}{\mathcal{T}}
\newcommand{\rpTransformd}[1]{\mathcal{#1}}

\newcommand{\coordinateTransformed}[2]{
\transpose{{\inverse{#2}}} #1 \inverse{#2}
}



\newcommand{\cristoffelSymbol}[3]{\Gamma_{#1#2#3}}

%%% #1 Reference frame: #2 notation
\newcommand{\cMatrixComponent}[3]{
  \Gamma_{#1#2#3}
}
%%%%%%%%%%% #1:row index, #2:column index, #3 index (between 1 and dof)
\newcommand{\cMatrixComponentDef}[3]{
  \frac{1}{2}\left(
    \pd{\inertiaMatrix_{#1#2}}{\jointPosition_{#3}}
    +\pd{\inertiaMatrix_{#1#3}}{\jointPosition_{#2}}
    -\pd{\inertiaMatrix_{#3#2}}{\jointPosition_{#1}}
  \right)
}



% #1:reference frame, #2:notation of the twist
\newcommand{\biasForce}{\bs{p}}

\newcommand{\kineticEnergy}{T(\jointPosition, \jointPositiond)}
\newcommand{\potentialEnergy}{V(\jointPosition)}
\newcommand{\totalEnergy}{H}

\newcommand{\lagR}{\mathcal{L}(\jointPosition, \jointPositiond)}

\newcommand{\lagRDef}{
  \quadratic{\jointPositiond}{\inertiaMatrix} + \potentialEnergy
}

%%% #1: dof number
\newcommand{\lagRDefDef}[1]{
  \frac{1}{2}\sum^{#1}_{i,j=1}\inertiaMatrix_{ij}(\jointPosition)\jointAngled_i\jointAngled_j - \potentialEnergy
}

% The Lagrange multipliers
\newcommand{\lm}{\bs{\lambda}}

%%% #1: link index
\newcommand{\lagREquations}[1]{
  \derivative{\pd{\lagR}{\jointAngled_{#1}}} - \pd{\lagR}{\jointPosition_{#1}} = \torque_{#1}
}

% The Lagrange's equation of a link
%%% #1: link index
\newcommand{\lagREquationsDef}[1]{
  \agg{j}{\dof}{
  (\inertiaMatrixd_{#1j}\jointPositiond_j  + \inertiaMatrix_{#1j}\jointPositiondd_j)
}
- \frac{1}{2}\agg{#1}{\dof}
{
  \pd{\inertiaMatrix_{#1j}}{\jointPosition_#1} \jointPositiond_#1\jointPositiond_j
  - \pd{\potentialEnergy}{\jointPosition_#1}
}
  = \torque_{#1}
}


\newcommand{\lagREquationsDefDef}[1]{
  \agg{j}{\dof}{
    \inertiaMatrix_{#1j}\jointPositiondd_j
  }
  + \frac{1}{2}\agg{j,k}{\dof}{(
    \pd{\inertiaMatrix_{#1j}}{\jointPosition_k}
    +\pd{\inertiaMatrix_{#1k}}{\jointPosition_j}
    - \pd{\inertiaMatrix_{kj}}{\jointPosition_#1}
    % \christoffelSymbolDef{#1}{j}{k}
  )\jointPositiond_k\jointPositiond_j}
+ \pd{\potentialEnergy}{\jointPosition_#1}
  = \torque_{#1}
}
%%%% #1 link index #2,#3 iteration index
\newcommand{\christoffelSymbol}[3]{
  \Gamma_{#1#2#3}
}
\newcommand{\christoffelSymbolDef}[3]{
  \frac{1}{2}(
    \pd{\inertiaMatrix_{#1#2}}{\jointPosition_#3}
    +\pd{\inertiaMatrix_{#1#3}}{\jointPosition_#2}
    - \pd{\inertiaMatrix_{#3#2}}{\jointPosition_#1})
}


% The argument are: force, torque

%%%%%%%%%%%%%%%%%%%%%%%%%%%%%%%%%%%%%%%%%%%%%%%%%%%%%%%%%%%%%%%%%%%%%%%%%%%%%%%%
%% 15. WRENCHES, FORCES, AND TORQUES
%%%%%%%%%%%%%%%%%%%%%%%%%%%%%%%%%%%%%%%%%%%%%%%%%%%%%%%%%%%%%%%%%%%%%%%%%%%%%%%%

\newcommand{\jointForce}[2]{\force^{#1}_{J#2}}%
\newcommand{\jointWrench}[2]{\wrench^{#1}_{J#2}}%

% #1 Reference frame, #2 joint index.
\newcommand{\linkWrench}[2]{\wrench^{#1}_{L#2}}%

% The actuation force of a joint
% #1 Reference frame, #2 joint index.
\newcommand{\jaForce}[2]{\force^{#1}_{a#2}}%
\newcommand{\jaWrench}[2]{\wrench^{#1}_{a#2}}%
% #1 joint index.
\newcommand{\jaWrenchDef}[1]{\wrench^\jointFrameDef{#1}_{a\jointFrameDef{#1}}}%


% The constraint force implements the motion constriant.
% #1 Reference frame, #2 joint index.
\newcommand{\jcForce}[2]{\force^{#1}_{c#2}}%
\newcommand{\jcWrench}[2]{\wrench^{#1}_{c#2}}%
% #1 joint index.
\newcommand{\jcWrenchDef}[1]{\wrench^{\jointFrameDef{#1}}_{c\jointFrameDef{#1}}}%

% The orthonormal basis of the constraint wrench. For example, T \in \RRm{6}{5} for a joint. and T \in \RRm{6}{1} for a point contact.
\newcommand{\inertialWrench}{{\wrench_{\inertialFrame}}}
\newcommand{\inertialForce}{{\force_{\inertialFrame}}}
\newcommand{\inertialTorque}{{\torque_{\inertialFrame}}}

\newcommand{\comForce}{{\force_{\com}}}
\newcommand{\comTorque}{{\torque_{\com}}}


% Z direction of the inertial frame
\newcommand{\cwc}{\bs{C}}

\newcommand{\wrenchC}[2]{
\begin{bmatrix}
#1\\
#2
\end{bmatrix}
}

% The argument are: force, torque
\newcommand{\wrenchR}[2]{
[#1^\top, #2^\top]^\top
}


%%%%%%%%%%%%%%%%%%%%%%%%%%%%%%%%%%%%%%%%%%%%%%%%%%%%%%%%%%%%%%%%%%%%%%%%%%%%%%%%
%% 16. CONTACT, FRICTION, AND IMPACT DYNAMICS
%%%%%%%%%%%%%%%%%%%%%%%%%%%%%%%%%%%%%%%%%%%%%%%%%%%%%%%%%%%%%%%%%%%%%%%%%%%%%%%%

\newcommand{\impactDuration}{\delta t}



\newcommand{\nextState}[1]{{#1}_{t_{k+1}}}
\newcommand{\nextStateB}[2]{{#1}_{#2,t_{k+1}}}
\newcommand{\nextStateTwo}[1]{{#1(t_{k+1})}}

\newcommand{\jump}{\Delta}
\newcommand{\impactduration}{\delta t}
\newcommand{\pointer}[2]{\overrightarrow{{#1}{#2}}}
\newcommand{\pointerDef}[2]{{\bs{#2}  - \bs{#1}}}

\newcommand{\pointerPerp}[2]{\overrightarrow{#1#2}_{\perp}}


\newcommand{\work}[1]{w_{#1}} % work done by #1
\newcommand{\impulse}{\bs{\iota}}
\newcommand{\impulseScalar}{\iota}

\newcommand{\quantityJump}{\jump \quantity}
\newcommand{\jtorquesJump}{\jump \jtorques}
\newcommand{\eevelocity}{\bs{\dot{x}}}
\newcommand{\eevelocityJump}{\jump \eevelocity}
\newcommand{\forceJump}{\jump \force}

\newcommand{\initial}[1]{{#1}_{0}}
\newcommand{\viable}[1]{{#1}^{\viableSymbol}}

\newcommand{\current}[1]{{#1}_{t_{k}}}
\newcommand{\currentTwo}[1]{#1(t_{k})}
\newcommand{\nextTwo}[1]{{#1(t_{k+1})}}
\newcommand{\previous}[1]{{#1}_{t_{k-1}}}

\newcommand{\impactNormal}{\bs{n}}

\newcommand{\surfaceNormal}{\bs{n}}
\newcommand{\projection}{P}

\newcommand{\projectionDef}[1]{{#1}{#1}^{\top}}

%%%%%%%%%%%%%%%%%%%%%%%%%%% C++ code

\newcommand{\lcpMass}{\text{M}_{\text{lcp}}}
\newcommand{\lcpd}{\text{d}_{\text{lcp}}}
% separation velocity:
\newcommand{\seVelocity}{\dot{\bs{\xi}}}
\newcommand{\seDistance}{\bs{\xi}}

%%%%%%%%%%%%%%%%%%%%%%%%%% Impact dynamics:

%%%%% Impact duration and structural dynamics

% The critical impact velocity that leads to visible plastic indentation.
\newcommand{\icVel}{\contactVel_{cr}}

% The yielding impact velocity that leads to plastic deformations that are not visible.
\newcommand{\iyVel}{\contactVel_{Y}}

%%%%%

\newcommand{\contactVel}{\bs{v}}
\newcommand{\nContactVel}{v_n}
\newcommand{\nContactVelR}{\contactVel_{nr}}
\newcommand{\nContactVelS}{\contactVel_{ns}}
\newcommand{\nContactVelL}{\contactVel_{nl}}
\newcommand{\nContactVelF}{\contactVel_{n,f}}
% The point where the tangential contactvel and its derivative are collinear.
\newcommand{\tVelL}{\contactVel_{gl}}
\newcommand{\tVelZero}{{\tp{\contactVel}}_{0}}

\newcommand{\cln}{\bs{s}}

% notation of the impacting body
\newcommand{\ib}{1}

% notation of the object
\newcommand{\ob}{2}

% virtual spring length
\newcommand{\vSpringL}{x}
\newcommand{\vSpringLd}{\dot{x}}
% spring coefficient
\newcommand{\vSpringC}{k}
\newcommand{\vDamperC}{c}

% The potential energy stored in the virtual spring
\newcommand{\pEnergy}{{{E}_p}}
\newcommand{\dEnergy}{{{E}_d}}
\newcommand{\dEnergyd}{{\dot{{E}}_d}}

% The kinetic energy
\newcommand{\kEnergy}{{E_k}}
\newcommand{\kEnergyd}{{\dot{E}_k}}

% #1 denotes the normal impulse
\newcommand{\pEnergyIntOne}[1]{\Phi_1(#1)}
% #1 denotes the normal impulse
\newcommand{\pEnergyIntTwo}[1]{\Phi_2(#1)}

\newcommand{\pEnergyd}{\dot{E}}

% Impulse at the end of the impact process
\newcommand{\impulseEnd}{\impulse_r}

% Derivative w.r.t. the normal impulse
\newcommand{\dni}[1]{\frac{d #1}{d \nImpulse} }
\newcommand{\dzi}[1]{\frac{d #1}{d \zImpulse} }
% Inverse inertia matrix
\newcommand{\iim}{W}
% The tangential plane iim
\newcommand{\iimtp}{B}
\newcommand{\iimtpDef}{\vectorTwo{\transpose{\basisVec{x}}}{\transpose{\basisVec{y}}} \iim \columnTwo{\basisVec{x}}{\basisVec{y}}}
% The z-axis iim
\newcommand{\iimcross}{\bs{d}}
\newcommand{\iimcrossDef}{\vectorTwo{\transpose{\basisVec{x}}}{\transpose{\basisVec{y}}} \iim \basisVec{z}}

\NewDocumentCommand{\coef}{m}{%
\IfStrEqCase{#1}
{
{f}{\mu}% friction
{s}{\bar{\mu}}% Coefficient of Stick
{r}{c_{\text{r}}}% Restitution coefficient
{er}{e_{\text{r}}}% Energetic restitution coefficient
}[] % "default" string}
}

\newcommand{\postImpact}[1]{{#1}^{+}}
\newcommand{\preImpact}[1]{#1^{-}}

% The impulse when the restitution phase ends
\newcommand{\impulseR}{\impulse_{r}}
\newcommand{\contactVelR}{\contactVel_{r}}
\newcommand{\contactVelC}{\contactVel_{c}}
% The impulse when the collinear condition applies
\newcommand{\impulseL}{\impulse_{l}}

% tangential direction unit vector
\newcommand{\tangentialU}{\hat{\bs{t}}}

% normal direction unit vector
\newcommand{\normalU}{\hat{\bs{n}}}

% The  unit vector of a coordinate system:
\newcommand{\basisVec}[1]{\widehat{\bs{\bs{#1}}}}


% The tangential impulse
\newcommand{\tImpulse}{\impulseScalar_{t}}
\newcommand{\tpImpulse}{\tp{\impulse}}
\newcommand{\tImpulsed}{\dot{\impulseScalar}_{t}}
% The tangential force
\newcommand{\tForce}{\forceScalar_{t}}

% Denote the tangential plane:
\newcommand{\tp}[1]{{#1}_{\perp}}

\newcommand{\tpForce}{\tp{\force}}

% The tangential linear Vel
\newcommand{\tVel}{\linearVelScalar_{t}}
\newcommand{\tVels}{\bs{\linearVelScalar}_{t}}
\newcommand{\hodog}{\dni{\tVels}}

% Curvature of the hodograph
\newcommand{\hodogC}{\kappa}
\newcommand{\hodogStepSize}{h}

% The normal linear Vel
\newcommand{\nVel}{\linearVelScalar_{n}}

% The normal impulse
\newcommand{\nImpulse}{{\impulseScalar_{n}}}
\newcommand{\nImpulseIni}{{\impulseScalar_{n,0}}}
\newcommand{\nImpulseF}{{\impulseScalar_{n,f}}}
\newcommand{\zImpulse}{{\impulseScalar_{z}}}
\newcommand{\nImpulsed}{\dot{\impulseScalar}_{n}}
\newcommand{\zImpulsed}{\dot{\impulseScalar}_{z}}
\newcommand{\impulseCone}{{^{\impulse}K}}
\newcommand{\fConeNum}{N_{\coef{f}}}

\newcommand{\applyConeSubscript}[1]{%
    %% \IfStrEqCase{#1}{%
    %% {}{
    %% \coef[f]
    %% }%
    %% {contact}{
    %% \text{sc}
    %% }%
    %% }[] % "default" string
    #1
}
\newcommand{\applyConeSuperscript}[1]{%
    \IfStrEqCase{#1}{%
    {f}{ % Friction cone
    \coef{f}
    }%
    {cwc}{ % Contact wrench cone
    \text{cwc}
    }%
    }[] % "default" string
}

\newcommand{\coneSymbol}{\text{K}}
\NewDocumentCommand{\cone}{O{}O{}}{
{
\coneSymbol^{\applyConeSuperscript{#1}}_{\applyConeSubscript{#2}}}
}

\newcommand{\generator}[1]{{\bs{\lambda}_{#1}}}
\newcommand{\generatorScalar}[1]{{\lambda_{#1}}}

\newcommand{\frictionCone}{{^{\coef{f}}K}}
\newcommand{\normalCone}{{^{\surfaceNormal}K}}
\newcommand{\impulseGenerator}{\bs{\lambda}}


% The normal force
\newcommand{\nForce}{{\forceScalar_{n}}}


% The normal impulse when the compression phase ends
\newcommand{\impulseC}{\impulse_{c}}
\newcommand{\nImpulseC}{\impulseScalar_{nc}}
\newcommand{\zImpulseC}{\impulseScalar_{zc}}
\newcommand{\zImpulseLCE}{\zeta_{zc}}
\newcommand{\czImpulseR}{\zeta_{r}} % candidate of the \zImpulseR
\newcommand{\nImpulseCEini}{\tilde{\impulseScalar}_{nc}}

% The normal impulse when the restitution phase ends
\newcommand{\nImpulseR}{\impulseScalar_{nr}}
\newcommand{\zImpulseR}{\impulseScalar_{zr}}

\newcommand{\zImpulseL}{\impulseScalar_{zl}}

% The smallest value of normal impulse at which v_t = 0.
\newcommand{\nImpulseS}{\impulseScalar_{ns}}
\newcommand{\nImpulseL}{\impulseScalar_{nl}}
\newcommand{\zImpulseS}{\impulseScalar_{zs}}
\newcommand{\impulseS}{\impulse_{s}}

% The derivative of impulse w.r.t. normal impulse before I_ns
\newcommand{\iDniB}{\bs{\delta}}
% The derivative of impulse w.r.t. normal impulse after the collinear condition applied.
\newcommand{\iDniL}{\bs{\delta}}

% The derivative of impulse w.r.t. normal impulse after I_ns
\newcommand{\iDniA}{\bs{\sigma}}

% The potential energy when the collinear condition applies
\newcommand{\pEnergyL}{\pEnergy_{l}}

% The potential energy when compression phase ends
\newcommand{\pEnergyC}{\pEnergy_{c}}

% The potential energy when v_t = 0 first holds
\newcommand{\pEnergyT}{\pEnergy_{s}}


%%%%%%%%%%%%%%%%%%%%%%%%%% state space models

%%%%%%%%%%%%%%%%%%%%%%%%%%%%%%%%%%%%%%%%%%%%%%%%%%%%%%%%%%%%%%%%%%%%%%%%%%%%%%%%
%% 17. LOCOMOTION AND BALANCE
%%%%%%%%%%%%%%%%%%%%%%%%%%%%%%%%%%%%%%%%%%%%%%%%%%%%%%%%%%%%%%%%%%%%%%%%%%%%%%%%

\newcommand{\ankleTorque}{{^a\bs{\tau}}}

% The arguments are: origin(or a point), surface normal
\newcommand{\plane}[2]{\Pi(#1, #2)}

% It rotates from frame #1 to #2.

%%%%%%%%%%%%%%%%%%%%%%%%% Polynomials
%% #1-#3 Coeffecients #4 variable name
\newcommand{\cop}[1][default]{
        \IfStrEq{#1}{text}{\text{CoP}}{
        {\text{cop}}
        }
}


\newcommand{\cmp}{{\textcolor{textColor}{\text{\bf c}}_{\text{mp}}}} % The centroidal moment pivot


\newcommand{\height}[1]{{d_{#1}}}

% The divergent component of motion (DCM) or the instantaneous capture point (ICP)
\newcommand{\dcmDef}{{\com + \frac{\com[dot]}{\pendulumC}}}
\newcommand{\dcmxDef}{{\com_x + \frac{\com[dot]_x}{\pendulumC}}}

% The center of mass (COM)
\newcommand{\comText}{{\text{CoM}}\xspace}


% stance foot location
\newcommand{\stanceFL}{{\bs{p}_s}}
% capture foot location
\newcommand{\captureFL}{{\bs{p}_c}}

% The arguments are: (1) starting point (2) the ending point. For instance: \vec{a} - \vec{b} = \pointer{b}{a}.
\newcommand{\areaSymbol}{\mathcal{S}}

\newcommand{\applyAreaSubscript}[1]{%
    \IfStrEqCase{#1}{%
    {zmp}{
     \zmp
    }%
    {com}{
     \com
    }%
    {comd}{
    \com[dot]
    }%
    {dcm}{
     \dcm
    }%
    }[] % "default" string
}

\NewDocumentCommand{\area}{O{}}{%
{    \ensuremath{%
        {\areaSymbol}_{\applyAreaSubscript{#1}}
    }
    }
}

\newcommand{\ssRegion}{{\set{\text{ssr}}}}

\newcommand{\controlArea}{{\mathcal{S}_{\com\zmp}}}
\newcommand{\contactArea}{{\mathcal{S}_{\contactPoint}}}
\newcommand{\footTask}{p}

% #1 joint number
\newcommand{\stepIndicator}{{\bf \text{O}}}
\newcommand{\stepIndicatorSum}{\stepIndicator_{\mpcHorizon}}

% #1 foot index
\newcommand{\footLocation}[1]{\contactPoint^{#1}}

\newcommand{\fL}{p}

\newcommand{\timeStep}{N}
% footStep
\newcommand{\footStep}{M}


% footStep
\newcommand{\step}[2]{
  {#1^{#2}}
}

%%%%%%%%%%%%%%%%%%%%%%%%%% Ellipsoid:

%%%%%%%%%%%%%%%%%%%%%%%%%%%%%%%%%%%%%%%%%%%%%%%%%%%%%%%%%%%%%%%%%%%%%%%%%%%%%%%%
%% 18. CONTROL
%%%%%%%%%%%%%%%%%%%%%%%%%%%%%%%%%%%%%%%%%%%%%%%%%%%%%%%%%%%%%%%%%%%%%%%%%%%%%%%%

\newcommand{\costFunctional}{J}
\newcommand{\stateWeight}{Q}

\newcommand{\inputWeight}{R}

\newcommand{\discountF}{\gamma}

% Set
\newcommand{\est}[1]{{\hat{#1}}}

% The damping ratio
\newcommand{\dRatio}{\xi}

% The critical frequency

\newcommand{\gainConstant}{\text{k}}
\newcommand{\stiffness}{\text{k}_{\text{s}}}

\newcommand{\pGain}{\gainConstant_{\text{p}}}
\newcommand{\dGain}{\gainConstant_{\text{d}}}
\newcommand{\iGain}{\gainConstant_{\text{i}}}

% #1:index of the pole
\newcommand{\pole}[1]{p_{#1}}

% The system momentum matrix
\newcommand{\samplingPeriod}{\Delta t}
\newcommand{\sampletime}{\Delta t}
\newcommand{\Kinv}{\bs{K}^{-1}}
\newcommand{\JEstimate}{\bs{K}_{\jvelocitiesJump}}
\newcommand{\JpinvEstimate}{\Kinv_{\jvelocitiesJump}}
\newcommand{\LambdaINVEstimate}{\bs{K}_{\impulse}}
\newcommand{\LambdaEstimate}{\Kinv_{\impulse}}

\newcommand{\weight}{w}

\newcommand{\cFreq}{{\omega_0}}
% The damped natural frequency
\newcommand{\dcFreq}{{\omega_d}}

\newcommand{\costate}{\lambda}%
\newcommand{\hamiltonian}{H}%

\newcommand{\initialTime}{t_0}%
\newcommand{\finalTime}{t_f}%

% #1 denotes the input weight: R; #2 denotes the state
\newcommand{\lqrControl}[2]{
  - \inverse{(#1 + \quadraticObj{B}{P_{n+1}})}\transpose{B}P_{n+1}A#2
}%

\newcommand{\manifold}{\mathcal{S}}%

% %%%%%%%%%%%%%%%%%%%%%%%%% LCP  VARIABLES:
\newcommand{\ssDim}{{n_s}}

\newcommand{\ssA}{A}
\newcommand{\ssB}{B}
\newcommand{\ssC}{C}
\newcommand{\ssD}{D}

\newcommand{\ssAd}{{A_d}}
\newcommand{\ssBd}{{B_d}}
\newcommand{\ssCd}{{C_d}}
\newcommand{\ssDd}{{D_d}}


\newcommand{\stateSymbol}{\text{\bf x}}
% #1 derivatives, #2 ref/min/, #3 x/y/xy, #4 current/next/previous
\NewDocumentCommand{\state}{O{}O{}O{}O{}}{%
  {    \ensuremath{%
      % Check the first optional argument for dot or double dot
      %% \applyOverheadSymbols{#1}\text{\bf z}
      \applyOverheadSymbols{#1}{\stateSymbol}
      % Check the second optional argument for 'ref'
      ^{\applySuperscript{#2}}%
      _{\applySubscript{#3}}
      \applyPostscript{#4}
    }
  }
}

% #1 current/next/previous
\NewDocumentCommand{\randomNoise}{O{}}{%
  {    \ensuremath{%
      \epsilon^{\applyPostscript{#1}}
    }
  }
}


\newcommand{\ssInput}{\text{\bf u}}

\newcommand{\x}[1]{{#1}_{x}}
\newcommand{\y}[1]{{#1}_{y}}
\newcommand{\z}[1]{{#1}_{z}}
\newcommand{\xy}[1]{{#1}_{x,y}}

\newcommand{\fwAngle}{\theta} % Fly wheel angle


%%%%%%%%%%%%%%%%%%%%%%%%%% QP commands:
\newcommand{\eqA}[1]{{\text{A}_{#1}}}
\newcommand{\eqb}[1]{{\text{b}_{#1}}}
\newcommand{\ieqC}[1]{{\text{C}_{#1}}}
\newcommand{\ieqd}[1]{{\text{d}_{#1}}}
\newcommand{\cpFrame}{\text{c}}

% Task-space elastic wrench:
\newcommand{\eWrench}{\wrenchFrame{\ee}{\Delta}}


% Object-level impedance control:
\newcommand{\ipdMass}{\text{K}_{\text{M}}}
\newcommand{\ipdDamper}{\text{K}_{\text{D}}}


% Operational space inertia matrix
\newcommand{\osim}[2][default]{
  \IfStrEq{#1}{def}
  {
    % Definition
    \inverse{(\quadraticObj{\jacobian_{#2}}
      {\inverse{\jsim}})}
  }{
    % Default notation
    \Lambda_{{#2}}
  }
}


% The desired frame
\newcommand{\oscc}[2][default]{%
  \IfStrEq{#1}{def}{
    % Gives the definition
    \transpose[inv]{\jacobian} \coriolis \inverse{\jacobian} - \osim{#2} \jacobian[dot]\inverse{\jacobian}
  }{
    \IfStrEq{#1}{dot}{
      % Defines the derivative
    }
    {
      % The default behavior
      \Gamma_{#2}}
  }
}


\newcommand{\osg}[2][default]{
  \IfStrEq{#1}{def}
  {
    % Definition
    \transpose{\jacobian_{#2}}
    \gravity
  }{
    \eta_{{#2}}
  }
}


%%%%%%%%%%%%%%%%%%%%%%%%%% MPC commands:
\newcommand{\olGain}{{\text{k}}}
\newcommand{\clGain}{{\text{K}}}



\newcommand{\totalCost}{{\text{J}_{0}}}


% Argument of the Minimum as a function of the perturbations around the i-th (x, u) pair
\newcommand{\argMinF}{{\text{Q}}}

% The argument denotes the step number
\newcommand{\costToGo}[1]{{\text{J}_{#1}}}
\newcommand{\costToGoDef}[1]{{\text{J}_{#1}}(\state, \controlSequence_{#1})}

% Value: the cost-to-go given a minimizing control sequence (not necessarily the optimal one), i.e., an instantiation of the cost-to-go.
\newcommand{\ctgIns}[1]{{\text{V}}(\state, #1)}

\newcommand{\ctgInsSimple}[1]{{\text{V}}^{#1}}

% #2 specifies the state
\newcommand{\ctgInsTwo}[2]{{\text{V}}(#2, #1)}

\newcommand{\ctgInsDef}[1]{\min_{\controlSequence_{#1}}\costToGoDef{#1}}

% #2 specifies the state
\newcommand{\ctgInsDefRecursion}[2]{\min_{\controlInput{#1}}
  (\runningCostDef{#1} + \ctgInsTwo{#1 + 1}{#2})
}



\newcommand{\totalCostDef}{\totalCost(\state, \controlSequence)}

\newcommand{\runningCost}{l}
\newcommand{\runningCostDef}[1]{\runningCost(\controlInput{#1}, \state_{#1})}

\newcommand{\finalCost}{{\runningCost_{\text{f}}}}
\newcommand{\finalCostDef}{\runningCost_{\text{f}}(\state_{\timeHL})}

\newcommand{\controlSequence}{{\bf{\text{U}}}}

% time horizon length
\newcommand{\timeHL}{N}

% time horizon
\newcommand{\timeH}{\text{T}}

% state dimension
\newcommand{\controlSequenceDef}[1]{\{\controlInput{#1}, \controlInput{#1 + 1},\ldots, \controlInput{\timeHL-1}\}}
\newcommand{\controlSequenceDefTwo}{\{\controlInput{0}, \controlInput{1},\ldots, \controlInput{\timeHL-1}\}}

% System Dynamics
\newcommand{\system}{{\text{f}}}

% Perturbed state
\newcommand{\perturb}[1]{{\delta #1}}
\newcommand{\perturbedState}[1]{{#1 + \delta #1}}

\newcommand{\objF}{J}

\newcommand{\mpcHorizon}{N}

\newcommand{\mpcfy}[1]{\bar{\bf #1}}



\newcommand{\mpcInput}{\mpcfy{\bf \controlInput{}}}
\newcommand{\mpcState}{\mpcfy{\state}}
\newcommand{\mpcStateWeight}{\mpcfy{Q}}
\newcommand{\mpcInputWeight}{\mpcfy{R}}

\newcommand{\mpcA}{\mpcfy{A}}
\newcommand{\mpcB}{\mpcfy{B}}
\newcommand{\mpcC}{\mpcfy{C}}
\newcommand{\mpcD}{\mpcfy{D}}

% Foot Step Location:
\newcommand{\decisionVariable}{\boldsymbol{\nu}}
\newcommand{\decisionVar}{\boldsymbol{\nu}}

% #1 velocity #2 mass
\newcommand{\desired}[1]{#1^{\text{des}}}
\newcommand{\optimized}[1]{#1^{*}}
\newcommand{\optimal}[1]{#1^{*}}
\newcommand{\reff}[1]{#1^{\text{ref}}}
\newcommand{\measured}[1]{#1^{\circ}}
\newcommand{\disturb}[1]{\delta{#1}}


% integration error
\newcommand{\iError}[1]{\int{\bs{e}_{#1}}}
\newcommand{\iErrorDef}[1]{\int{\errorDef{#1}}dt}

% Defines the function to be minimized in a QP
\newcommand{\function}[1]{\text{f}_{#1}}

\newcommand{\error}[1]{\bs{e}_{#1}}
\newcommand{\errord}[1]{\dot{\bs{e}}_{#1}}
\newcommand{\errordd}[1]{\ddot{\bs{e}}_{#1}}

\newcommand{\errorDef}[1]{
  #1 - \reff{#1}
}

% The constant term of the error dynamics
\newcommand{\cTerm}[1]{
  \text{C}_{#1}
}

% The quadratic term of the error dynamics
\newcommand{\qTerm}[1]{
  \text{Q}_{#1}
}

% The penalty term of the error dynamics
\newcommand{\pTerm}[1]{
  \text{P}_{#1}
}


\newcommand{\errorDefFull}[1]{
 \selectionMatrix #1 - \selectionMatrix \reff{#1}
}

% #1:index of the input
\newcommand{\controlInput}[1]{{\text{\bf{u}}_{#1}}}
\newcommand{\control}[1]{{\text{\bf{u}}_{#1}}}

%% #1 real part, #2 imaginery part

%%%%%%%%%%%%%%%%%%%%%%%%%%%%%%%%%%%%%%%%%%%%%%%%%%%%%%%%%%%%%%%%%%%%%%%%%%%%%%%%
%% 19. STATE ESTIMATION AND OBSERVERS
%%%%%%%%%%%%%%%%%%%%%%%%%%%%%%%%%%%%%%%%%%%%%%%%%%%%%%%%%%%%%%%%%%%%%%%%%%%%%%%%

\newcommand{\oGain}{K_O} % The diagonal observer gain matrix.

\newcommand{\reminder}{\bs{r}_{\text{M}}}
\newcommand{\reminderd}{\dot{\bs{r}}_{\text{M}}}
\newcommand{\reminderdDef}{\oGain(\momentum - \est{\momentumd})}

\newcommand{\ext}[1]{{#1}_{\text{ext}}}

%%%%%%%%%%%%%%%%%%%%%%%%%% special terms
\newcommand{\bNoise}[1]{\text{b}_{#1}}

% Gaussian noise
\newcommand{\gNoise}[1]{w_{#1}}


% covariance matrix
\newcommand{\cov}[1]{\text{Q}_{#1}}

% Leg kinematics
% Distance from the CoM to a contact point i: p_i - com
\newcommand{\disDef}[1]{\pointer{\com}{\contactPoint_{#1}}}
\newcommand{\dis}[1]{\text{s}_{#1}}



% Forward kinematics
\newcommand{\fk}[1]{\text{lkin}_{#1}(\jangles)}

\newcommand{\imu}{\text{m}}



% subequations

%%%%%%%%%%%%%%%%%%%%%%%%%% Impedance control commands:

% The compliant frame

%%%%%%%%%%%%%%%%%%%%%%%%%%%%%%%%%%%%%%%%%%%%%%%%%%%%%%%%%%%%%%%%%%%%%%%%%%%%%%%%
%% 20. COLLISION AVOIDANCE
%%%%%%%%%%%%%%%%%%%%%%%%%%%%%%%%%%%%%%%%%%%%%%%%%%%%%%%%%%%%%%%%%%%%%%%%%%%%%%%%

\newcommand{\link}[1]{\text{l}_{#1}}


%  Define a macro with an optional argument for dotting
\newcommand{\cDist}[3][]{#1{\text{dis}}_{#2,#3}}
% The secure distance that two links should never reach
\newcommand{\sDist}[1]{\sigma_{\text{s}, #1}}

% The threshold distance below which, the collision avoidance constraint will be activated
\newcommand{\tDist}[1]{\sigma_{\text{t}, #1}}

\newcommand{\dConstant}[1]{\text{c}_{#1}}

%%%%%%%%%%%%%%%%%%%%%%%%%% State estimation commands:

% bias (drift) noise

%%%%%%%%%%%%%%%%%%%%%%%%%%%%%%%%%%%%%%%%%%%%%%%%%%%%%%%%%%%%%%%%%%%%%%%%%%%%%%%%
%% 21. MISCELLANEOUS
%%%%%%%%%%%%%%%%%%%%%%%%%%%%%%%%%%%%%%%%%%%%%%%%%%%%%%%%%%%%%%%%%%%%%%%%%%%%%%%%

\newcommand{\uncertain}[1]{\widetilde{#1}}

\newcommand{\eePos}{\bs{x}} % End-effector position
\newcommand{\eeVel}{\dot{\eePos}}
\newcommand{\eeAcc}{\ddot{\eePos}}
\newcommand{\optVars}{\bs{u}}

% The center of pressure (COP) point
\newcommand{\specialJac}{\mathcal{J}}
\newcommand{\specialJacTwo}{\mathcal{C}}

%%%%%%%%%%%%%  Different support areas
\newcommand{\stateVar}{\bs{x}}

% #1 variable with different steps, #2 step number
\newcommand{\semiAxis}{
a
}

\newcommand{\density}{
\rho
}



%%%%%%%%%%%%%%%%%%%%%%%%%% Eigenvalue:


\newcommand{\eigV}{
\sigma
}

\newcommand{\eigVec}{
\bs{e}
}



%%%%%%%%%%%%%%%%%%%%%%%%%% Rigid body motion


\newcommand{\guassD}{\mathcal{D}}

\newcommand{\quantity}{\bs{d}}
\newcommand{\quantityd}{\dot{\quantity}}
\newcommand{\quantitydd}{\ddot{\quantity}}

%%%%%%%%%%%% Axis:
\newcommand{\pre}[1]{
  \setSymbol^p(#1)
}


%%% all the links between [root and predecessors] while(link.hasPredecessor()){set+=predecessor(link); link = link.predecessor}):
\newcommand{\recursivePre}[1]{
  \setSymbol^{rp}(#1)
}

%%% successor:
\newcommand{\suc}[1]{
  \setSymbol^s(#1)
}

%%% all the links between [root and predecessors] while(link.hasSuccessor()){set+=successor(link); link = link.successor}):
\newcommand{\recursiveSuc}[1]{
  \setSymbol^{rs}(#1)
}


%%% CRB vel


\newtheorem{theorem}{\oBlock{Theorem}}
\newtheorem{Lemma}{\oBlock{Lemma}}
\newtheorem{remark}{\oBlock{Remark}}[section]
\newtheorem{assumption}{assumption}
\newtheorem{hypothesis}{hypothesis}
\newtheorem{proposition}[theorem]{proposition}
\newtheorem{property}[section]{\oBlock{property}}
\newtheorem{corollary}[theorem]{corollary}
\newtheorem{Example}[theorem]{Example}
\newtheorem{problem}{Problem}
\newtheorem{Solution}{Solution}
\newtheorem{Definition}{Definition}


\newtheorem{thm}{Theorem}[section]
\newtheorem{cor}[thm]{Corollary}
\newtheorem{lem}[thm]{Lemma}
\newtheorem{prop}[thm]{Proposition}

\newtheorem{defn}[thm]{Definition}

\newtheorem{rem}[thm]{Remark}

\usepackage{../commands/yuquanTitle}

%%%%%%%%%%%%%%%%%%%%%%%%%%% -------------------------------------------

% Report-local notation, kept identical to rolling-contact-qp.tex so that a
% reader can move between the two documents without re-learning symbols.
\newcommand{\generalizedConfiguration}{\jangles[][][generalized]}
\newcommand{\generalizedVelocity}{%
  \applyOverheadSymbols{dot}{\applyJangles{generalized}}}
\newcommand{\generalizedAcceleration}{\jangles[ddot][][generalized]}
\newcommand{\rollingDirection}{\bs{t}}
\newcommand{\lateralDirection}{\bs{s}}
\newcommand{\wheelAngle}{\phi}
\newcommand{\wheelRadius}{r}
\newcommand{\wheelSelector}{\bs{S}^{\wheelAngle}}
\newcommand{\steeringSelector}{\bs{S}^{\delta}}
\newcommand{\tangentProjector}{\text{P}}
\newcommand{\slipVelocity}[1]{\contactVel_{\text{slip},#1}}
\newcommand{\rollingMatrix}{\text{A}}
\newcommand{\rollingTarget}{\bs{b}}
\newcommand{\rollingResidual}{\bs{e}^{\text{roll}}}
\newcommand{\slackVariable}{\bs{\sigma}}

% Test-suite notation.
\newcommand{\testTol}{\varepsilon}
\newcommand{\numGen}{n_{\text{g}}}
\newcommand{\parentDoc}{\texttt{rolling-contact-qp.tex}\xspace}
\newcommand{\harness}{\texttt{rolling\_contact\_qp\_checks.py}\xspace}

% A test card. #1 identifier, #2 title, #3 kind, #4 priority.
\newenvironment{testcard}[4]{%
  \par\medskip\noindent
  \tBlock{#1}\quad\textbf{#2}%
  \hfill{\small\textit{#3}\,/\,#4}\par\nopagebreak\vspace{2pt}
  \begin{description}[leftmargin=7.5em,style=sameline,itemsep=1pt,topsep=1pt]
}{%
  \end{description}\par\medskip
}

\project{\tinytf Tests for the rolling-contact QPs}
\author{
  Yuquan WANG \\
  Tencent Robotics X \\
  Shenzhen, P.R. China\\
  \vspace{5pt}
  \texttt{hasbragwang@tencent.com}\vspace{30pt} \\
}
\summary{
  This companion report specifies a unit- and smoke-test suite for an
  implementation of the two explicit planar quadratic programs of the
  rolling-contact report: the differential-drive QP and the four-steering-wheel
  QP. Each test states the invariant asserted, the exact regime in which that
  invariant holds, a tolerance tied to a physical scale, and the diagnosis a
  failure implies. Three invariants that are commonly assumed but that are only
  conditionally true---or, in one case, false---are corrected explicitly before
  the suite begins.
}

\rhead{Robotics-X \colorbox{teal}{\textcolor{white}{\quad \bf Rolling-contact QP tests}}}
%%%%%%%%%%%%%%%%%%%%%%%%%%% -------------------------------------------

\begin{document}

\maketitle

\section{Scope and how to read this suite}
\label{sec:test-scope}

This report specifies tests for an implementation of
\eqref{eq:differential-drive-qp} and \eqref{eq:four-steering-wheel-qp} of the
rolling-contact report \parentDoc. Throughout, \tBlock{T1} denotes the
differential-drive QP and \tBlock{T2} the four-steering-wheel QP. Equation
labels without further qualification refer to that report.

The suite tests the \emph{corrected} formulations. Readers working from an
earlier revision should note that the planar basis is now chassis aligned, the
reduced matrices are a congruence rather than a row selection, the normal load
enters through an exogenous row, the force generators are regularized, and the
four lateral rows of T2 carry slacks. Tests that would pass under the earlier
revision but fail under the corrected one, or the converse, are marked.

\subsection{Taxonomy}

\begin{description}[leftmargin=6.5em,style=sameline,itemsep=2pt]
  \item[unit] Isolates one assembly function: a frame, a
    matrix block, a single constraint row. No solver call.
  \item[smoke] A fast, decisive, end-to-end run of the
    assembled controller. Asserts trajectory-level outcomes.
  \item[property] Holds for every draw from a randomized
    fixture family; states the draw count and the seed policy.
  \item[negative] Asserts that something is \emph{not}
    done, or that a malformed input is rejected rather than silently accepted.
  \item[regression] Pins a value or a structure that has
    been wrong before.
\end{description}

Priorities are \oBlock{P0}, which must pass before the QP is trusted at all;
\oBlock{P1}, required for release; and \oBlock{P2}, valuable but deferrable.

\subsection{Tolerance policy}

A tolerance is only meaningful next to the scale it measures. Every entry
therefore states its tolerance relative to a physical quantity rather than as a
bare number. Three conventions are used throughout:
\begin{enumerate}[label=(T\arabic*)]
  \item \label{tol:exact} \tBlock{Exact algebra}: quantities that must agree to
    working precision, tolerance $\testTol_{\text{alg}}=10^{-9}$ relative to the
    larger operand. Used for identities that hold symbolically.
  \item \label{tol:fd} \tBlock{Finite-difference oracles}: tolerance
    $\testTol_{\text{fd}}=10^{-6}$ absolute with a step $h=10^{-6}$, together
    with an assertion that halving $h$ reduces the error quadratically. The
    convergence-rate assertion is the real test; a bare threshold can be passed
    by a wrong derivative with a small coefficient.
  \item \label{tol:phys} \tBlock{Physical residuals}: tolerance stated as a
    fraction of a named scale, for example $10^{-4}$ of $\mass\gravity$ for a
    force residual or $10^{-3}$\unitVelTS for a velocity residual.
\end{enumerate}

\subsection{Relation to the numerical harness}

A companion Python harness, \harness, already implements twelve of these
checks. Entries that it covers are marked \emph{implemented}; the remainder are
specified here for implementation against the production controller. The
harness is solver independent by design: it uses finite differences, matrix
ranks, and explicit nullspace constructions, so a failure points at the
mathematics rather than at a quadratic-programming library.

\section{Three invariants that are commonly assumed, and their corrections}
\label{sec:test-corrections}

Before the suite proper, three plausible invariants must be restated. Two hold
only in a named regime and one is false. Encoding any of them naively produces
a test that either cannot fail or that fails a correct implementation.

\subsection{Contact forces and gravity}
\label{sec:correction-gravity}

The natural statement, that the contact forces sum to the robot's weight, is a
\emph{static-equilibrium specialization} of Newton's law, not an invariant of
the controller. The general law for the centroidal linear momentum is
\quickEq{eq:test-newton-general}{
  \agg{i}{\numContacts}{\force_i} + \mass\gravity
  = \mass\com[ddot],
  \qquad\text{equivalently}\qquad
  \agg{i}{\numContacts}{\force_i}
  = \mass\left(\com[ddot]-\gravity\right).
}
Only when $\com[ddot]=\zeroVector$ and no other external wrench acts does this
reduce to $\agg{i}{\numContacts}{\force_i}=-\mass\gravity$. On a plane of slope
$\theta$ the useful specialization splits into
\quickEq{eq:test-ramp-split}{
  \agg{i}{\numContacts}{\force_{n,i}} = \mass\gravity[value]\cos\theta,
  \qquad
  \norm{\agg{i}{\numContacts}{\tangentProjector_i\force_i}}
  = \mass\gravity[value]\sin\theta,
}
which is feasible only while $\coef{f}\geq\tan\theta$.

A structural caveat matters more than either formula. Both planar QPs carry
only the three base coordinates $(v_x,v_y,\omega_{\Pi})$, and by
\eqref{eq:reduced-normal-annihilated} the reduced contact Jacobians annihilate
$\surfaceNormal_{\Pi}$. Neither QP therefore contains a surface-normal force
balance at all. The normal split is supplied exogenously by
\eqref{eq:planar-normal-load-row}. Asserting
\eqref{eq:test-ramp-split} against a planar QP solution therefore tests the
\emph{normal-load model}, not the QP; asserting it against the generic
whole-body QP \eqref{eq:ideal-rolling-whole-body-qp} does test the solver.
Tests \texttt{DYN-G1}--\texttt{DYN-G3} separate these two cases.

\subsection{Contact forces and the friction cone}
\label{sec:correction-cone}

The statement that every contact force lies in its friction cone is true, but
it is true \hlBlock{by construction}: with $\force_i=\cone[f][i]\generator{\force_i}$ and
$\generator{\force_i}\geq\zeroVector$ the reconstructed force cannot leave the
polyhedral cone. A test that recomputes
$\norm{(\force_{t,i},\force_{s,i})}\leq\coef{f}_i\force_{n,i}$ from the
solution exercises the reconstruction code and nothing else; it cannot fail on
a wrong model.

The tests with teeth are three different assertions:
\begin{enumerate}[label=(\roman*)]
  \item that $\cone[f][i]$ is rebuilt in the \emph{local} rolling frame
    $[\rollingDirection_i,\lateralDirection_i,\surfaceNormal_i]$ whenever the
    wheel steers or the terrain normal changes---a cone cached in a world frame
    is the classic defect, and it still passes the naive membership test;
  \item that the polygon is an \emph{inner} approximation, so its usable
    tangential fraction is only $\cos(\pi/\numGen)$ of the true cone, which
    bounds how hard the controller may push before the model, not the physics,
    refuses;
  \item that the cone actually \emph{binds}. Without
    \eqref{eq:planar-normal-load-row} it does not: since
    $\transpose{\cone[f][i]}$ maps a symmetric generator set to
    $\cone[f][i]\bs{1}\propto\surfaceNormal_i$, and the reduced rows annihilate
    $\surfaceNormal_{\Pi}$, the substitution
    $\generator{\force_i}\mapsto\generator{\force_i}+\text{s}\bs{1}$ with
    $\text{s}\geq0$ satisfies every constraint and changes no objective term.
    The cone is vacuous and the optimal set is an unbounded ray.
\end{enumerate}
Tests \texttt{FRI-01}--\texttt{FRI-06} implement these.

\subsection{Steering torques and yaw momentum}
\label{sec:correction-yaw}

The statement that the steering torques sum to the yaw-momentum derivative is
\hlBlock{false}, and it is false structurally rather than approximately.
Actuator torques are \emph{internal}: in
\eqref{eq:rolling-equations-of-motion} the actuation matrix
$\actuationMatrix$ has identically zero rows for the unactuated floating-base
coordinates, and the congruence $\text{B}_{\text{red}}
=\transpose{\text{P}}\actuationMatrix$ of
\eqref{eq:reduced-dynamics-congruence} preserves that. No steering torque can
appear in the chassis-yaw row at all.

What does hold is the centroidal angular-momentum balance. The plane-normal
component of the rate of centroidal angular momentum equals the net moment of
the \emph{contact} forces about the center of mass, gravity contributing none:
\quickEq{eq:test-yaw-balance}{
  \innerP{\surfaceNormal_{\Pi}}{\angularMomentum[dot][\comSymbol]}
  = \agg{i}{\numContacts}{
      \innerP{\surfaceNormal_{\Pi}}
        {\left(\cross{(\point[][][\inertialFrame,p_i]-\com)}{\force_i}\right)}
    }.
}
Separately, each steering torque satisfies its own joint row of the reduced
dynamics, balancing the steering-column inertia, the transmitted scrub and
self-aligning moment, and the coupling terms in $\bs{h}_{\text{red}}$. These
are different equations with different content.

Both facts are worth testing. Test \texttt{DYN-Y1} asserts
\eqref{eq:test-yaw-balance}. Test \texttt{DYN-Y2} is a \emph{negative} test
asserting $\transpose{\bs{e}_{\omega}}\text{B}_{\text{red}}=\transpose{\zeroVector}$,
which fails loudly for an implementation that wires steering torque into the
chassis yaw equation---a plausible error that
\eqref{eq:test-yaw-balance} alone would not catch.

\section{Shared fixtures}
\label{sec:test-fixtures}

Every test draws from one of the fixtures below so that a tolerance calibrated
in one entry stays meaningful in the others. Values are those of a
mobile-manipulator-scale wheeled base and match the defaults of \harness.

\begin{table}[h]
  \centering
  \caption{Shared numeric fixtures}
  \label{tab:test-fixtures}
  \begin{tabular}{|C{0.20\columnwidth}|C{0.30\columnwidth}|C{0.38\columnwidth}|}
    \hline
    Fixture & Parameters & Used by \\
    \hline
    \texttt{F-COMMON} &
      $\mass=45$\unitMass, $\gravity[value]=9.81$,
      $\wheelRadius_i=0.15$\unitMeter, $\coef{f}=0.7$,
      $\samplingPeriod=2\unitMs$, $\pGain=50$, $\numGen=6$ &
      every test; the base parameter set \\
    \hline
    \texttt{F-DIFF} &
      $b=0.50$\unitMeter, $\bs{\rho}_L=(0,\,b/2)$, $\bs{\rho}_R=(0,-b/2)$,
      yaw inertia $3.2$ &
      all T1 entries \\
    \hline
    \texttt{F-4S} &
      $\bs{\rho}_i=(\pm0.40,\,\pm0.30)$\unitMeter, four steering joints,
      $\weight_{\sigma}$ common across wheels &
      all T2 entries \\
    \hline
    \texttt{F-RAMP} &
      \texttt{F-COMMON} with slope $\theta\in\{0,10,20,30,40\}$\unitDegree
      about the chassis lateral axis &
      ramp and friction-margin entries \\
    \hline
    \texttt{F-ICR} &
      \texttt{F-4S} with steering angles solved for a prescribed instantaneous
      center of rotation at $(1.10,\,0.80)$\unitMeter &
      T2 feasibility and slack entries \\
    \hline
    \texttt{F-RANDOM} &
      randomized $\bs{\rho}_i$, $\delta_i$, $\wheelRadius_i$, $\coef{f}$,
      slope and references; $500$ draws, seed fixed and recorded &
      all property entries \\
    \hline
  \end{tabular}
\end{table}

Slope $40$\unitDegree in \texttt{F-RAMP} is deliberately beyond the friction
limit, since $\tan 40^{\circ}=0.839>\coef{f}$: it is the fixture for the
tests that must predict failure rather than report success.

\section{Layer A. Geometry and frame data}
\label{sec:test-layer-a}

These tests exercise the code that builds QP data before any row is written.
They require no solver, run in milliseconds, and catch the largest class of
defects.

\begin{testcard}{GEO-01}{Convention pin: planar geometry is chassis-heading invariant}{regression}{P0}
  \item[Targets] \eqref{eq:chassis-aligned-tangent-basis},
    \eqref{eq:wheel-center-planar-offset}, \eqref{eq:planar-carrier-map}
  \item[Regime] Flat plane or constant-slope ramp; rigid carrier mount.
  \item[Invariant] Rotate the whole scene about $\surfaceNormal_{\Pi}$ by an
    arbitrary heading $\psi$. Every planar quantity must be unchanged:
    $\bs{\rho}_i(\psi)=\bs{\rho}_i(0)$, $\text{H}_i(\psi)=\text{H}_i(0)$, and
    $\rollingDirection_i^{\Pi}(\psi)=\rollingDirection_i^{\Pi}(0)$.
  \item[Pass] Sweep $\psi$ over $[-\pi,\pi]$ in $64$ steps; maximum deviation
    below $\testTol_{\text{alg}}$, convention~\ref{tol:exact}.
  \item[Diagnosis] Any drift means the implementation resolved a planar
    quantity in a world-fixed basis. Every downstream row is then short of the
    Coriolis terms listed in Remark~\ref{re:planar-basis-choice}. This is the
    single highest-value test in the suite; run it first.
\end{testcard}

\begin{testcard}{GEO-02}{Chassis-aligned basis is orthonormal, right-handed, and body-constant}{unit}{P0}
  \item[Targets] \eqref{eq:chassis-aligned-tangent-basis},
    \eqref{eq:terrain-basis-body-frame}
  \item[Regime] Chassis forward axis not parallel to $\surfaceNormal_{\Pi}$.
  \item[Invariant] $\transpose{\text{E}_{\Pi}}\text{E}_{\Pi}=\bs{I}_2$;
    $\transpose{\text{E}_{\Pi}}\surfaceNormal_{\Pi}=\zeroVector$;
    $\det[\bs{e}_{\Pi,x}\ \bs{e}_{\Pi,y}\ \surfaceNormal_{\Pi}]=+1$; and
    ${}^{\fbFrame}\text{E}_{\Pi}$ is constant under a pure heading change.
  \item[Pass] Convention~\ref{tol:exact}; the determinant must be $+1$, not
    $\pm1$.
  \item[Diagnosis] A determinant of $-1$ flips the sign of every yaw moment
    while leaving all norms correct, so it survives every magnitude-only check.
\end{testcard}

\begin{testcard}{GEO-03}{Degenerate forward axis aborts rather than normalizing garbage}{negative}{P1}
  \item[Targets] \eqref{eq:chassis-aligned-tangent-basis}
  \item[Regime] Chassis pitched so its forward axis approaches
    $\surfaceNormal_{\Pi}$.
  \item[Invariant] As
    $\norm{\tangentProjector_{\Pi}\rotation{\inertialFrame}{\fbFrame}\basisVec{x}}
    \to0$ the assembly must raise a diagnosable error, not divide by a small
    number.
  \item[Pass] Sweep the pitch to within $10^{-8}$\unitDegree of the
    degeneracy; the implementation must reject below a declared threshold and
    must never emit a non-finite basis.
  \item[Diagnosis] A silent normalization produces a basis that is orthonormal
    to machine precision but whose direction is noise, so GEO-01 passes while
    the controller steers randomly.
\end{testcard}

\begin{testcard}{GEO-04}{Local contact triad and the tangent projector}{unit}{P0}
  \item[Targets] \eqref{eq:rolling-local-frame}, \eqref{eq:tangent-projector}
  \item[Regime] Any wheel pose; per-wheel normals allowed to differ from
    $\surfaceNormal_{\Pi}$ only if assumption~\ref{as:one-plane} is relaxed.
  \item[Invariant] $\lateralDirection_i=\cross{\surfaceNormal_i}{\rollingDirection_i}$;
    the triad is orthonormal and right handed;
    $\tangentProjector_i^2=\tangentProjector_i=\transpose{\tangentProjector_i}$;
    $\tangentProjector_i\surfaceNormal_i=\zeroVector$;
    $\text{tr}\,\tangentProjector_i=2$; and the claimed identity
    $\tangentProjector_i
     =\rollingDirection_i\transpose{\rollingDirection_i}
     +\lateralDirection_i\transpose{\lateralDirection_i}$ holds.
  \item[Pass] Convention~\ref{tol:exact}.
  \item[Diagnosis] The trace and idempotence checks localize the fault: a trace
    of $3$ means the normal was never subtracted, a trace of $1$ means it was
    subtracted twice.
\end{testcard}

\begin{testcard}{GEO-05}{$\text{H}_i$ is the planar form of $\contactVel+\cross{\bs{\omega}}{\bs{\rho}}$}{unit}{P0}
  \item[Targets] \eqref{eq:planar-carrier-map}
  \item[Regime] Rigid carrier mount, planar motion.
  \item[Invariant] For random $\bs{\xi}^{\Pi}_{\fbFrame}$ and $\bs{\rho}_i$,
    $\text{H}_i\bs{\xi}^{\Pi}_{\fbFrame}$ equals the tangential projection of
    the full three-dimensional rigid-body velocity
    $\contactVel_{\fbFrame}+\cross{\bs{\omega}}{\bs{d}_i}$ computed
    independently.
  \item[Pass] Convention~\ref{tol:exact} over $200$ random draws from
    \texttt{F-RANDOM}.
  \item[Diagnosis] A sign error in the $-y_i$ or $+x_i$ entry inverts the
    moment arm of one wheel, which no single-wheel test detects.
\end{testcard}

\begin{testcard}{GEO-06}{Contact point offset is invisible to the planar rows}{unit}{P1}
  \item[Targets] \eqref{eq:wheel-center-contact-point},
    \eqref{eq:reduced-normal-annihilated}
  \item[Regime] Planar reduction only.
  \item[Invariant] Replacing $\point[][][\inertialFrame,p_i]$ by
    $\point[][][\inertialFrame,c_i]$ changes no planar base row, because the
    two differ by $-\wheelRadius_i\surfaceNormal_i$ and the reduced rows
    annihilate $\surfaceNormal_{\Pi}$. It does change the wheel-spin column.
  \item[Pass] Base rows identical to $\testTol_{\text{alg}}$; wheel column
    differs by exactly $\wheelRadius_i$.
  \item[Diagnosis] This test tells the implementer where the carrier/contact
    distinction does and does not matter, preventing a spurious ``fix''.
\end{testcard}

\begin{testcard}{GEO-07}{Steering angle is chassis relative, and its rate is the joint rate}{negative}{P0}
  \item[Targets] \eqref{eq:steering-directions},
    \eqref{eq:steering-direction-derivatives},
    \eqref{eq:steering-wheel-angular-velocity}
  \item[Regime] T2, assumption~\ref{as:steering-axis}.
  \item[Invariant] With the steering modules mechanically locked and the
    chassis yawing at $\omega_{\Pi}$, every $\dot\delta_i$ read by the QP must
    be zero. Substituting the absolute heading rate
    $\dot\delta_i^{\text{abs}}=\dot\delta_i+\omega_{\Pi}$ must make this test
    fail.
  \item[Pass] $\abs{\dot\delta_i}<10^{-9}$\unitVelJS with the modules locked
    and $\omega_{\Pi}=1$\unitVelJS; the mutant must exceed
    $0.9$\unitVelJS.
  \item[Diagnosis] The absolute convention silently consumes the entire
    steering-rate budget during any turn, and biases
    $\dot{\rollingDirection}^{\Pi}_i$ by $\omega_{\Pi}\lateralDirection^{\Pi}_i$.
\end{testcard}

\begin{testcard}{GEO-08}{Reduced contact Jacobian: normal annihilation and cross-wheel sparsity}{unit}{P0}
  \item[Targets] \eqref{eq:reduced-normal-annihilated},
    \eqref{eq:reduced-dynamics-congruence}
  \item[Regime] Planar reduction, any steering configuration.
  \item[Invariant] $\transpose{\text{J}_{\text{red},i}}\surfaceNormal_{\Pi}
    =\zeroVector$ exactly; and $\text{J}_{\text{red},i}$ has bit-exact zeros on
    every joint column except wheel $i$'s own spin and steering columns.
  \item[Pass] Convention~\ref{tol:exact}; the sparsity assertion is exact
    equality to zero, not a threshold.
  \item[Diagnosis] A nonzero cross-wheel entry means one wheel's contact force
    is being fed into another wheel's torque row. A nonzero normal projection
    means the implementer restored an out-of-plane coordinate without saying so,
    which invalidates \eqref{eq:planar-normal-load-row}.
\end{testcard}

\begin{testcard}{GEO-09}{ICR concurrency is exactly the rank condition on the lateral block}{property}{P1}
  \item[Targets] \eqref{eq:steering-expanded-constraints},
    Remark~\ref{re:four-steering-lateral-slack}
  \item[Regime] T2, fixture \texttt{F-ICR} and \texttt{F-RANDOM}.
  \item[Invariant] Let $\text{A}$ stack
    $\transpose{(\lateralDirection^{\Pi}_i)}\text{H}_i$. Then
    $\text{rank}(\text{A})=2$ if and only if the four axle lines are
    concurrent; otherwise $\text{rank}(\text{A})=3$.
  \item[Pass] $500$ draws: concurrent constructions give
    $\sigma_3/\sigma_1<10^{-10}$ after nondimensionalization by a length scale;
    random ones give $\sigma_3/\sigma_1>10^{-3}$. Nondimensionalize before
    thresholding, since $\text{A}$ mixes dimensionless and length entries.
  \item[Diagnosis] Thresholding the raw singular value makes the verdict depend
    on whether offsets are in metres or millimetres.
\end{testcard}

\begin{testcard}{GEO-10}{Rows and contact Jacobian are built from one geometry}{unit}{P1}
  \item[Targets] \eqref{eq:planar-carrier-map},
    \eqref{eq:reduced-dynamics-congruence}
  \item[Regime] Both chassis.
  \item[Invariant] Perturb a single stored $\delta_i$ or $\bs{\rho}_i$ and
    require that the kinematic rows \emph{and} $\text{J}_{\text{red},i}$ both
    change consistently: they must be functions of the same cached geometry.
  \item[Pass] Both change; neither is stale. Compare against a full rebuild.
  \item[Diagnosis] Catches a partially updated cache, in which the kinematics
    steer but the force map does not.
\end{testcard}

\section{Layer B. Constraint-row assembly}
\label{sec:test-layer-b}

These tests assert that each printed row is the exact derivative of its
velocity-level constraint, and that the row set has the intended rank.

\begin{testcard}{ROW-01}{T1 rows are the exact derivative at $\pGain=0$}{unit}{P0}
  \item[Targets] \eqref{eq:differential-drive-qp},
    \eqref{eq:differential-geometric-constraints}
  \item[Regime] Assumptions~\ref{as:planar-basis} and~\ref{as:rigid-wheel};
    fixture \texttt{F-DIFF}.
  \item[Invariant] With $\pGain=0$, the three printed rows equal
    $\frac{d}{dt}$ of the velocity constraints along a trajectory that
    satisfies them, verified by central differences.
  \item[Pass] Convention~\ref{tol:fd}, including the quadratic convergence
    assertion.
  \item[Diagnosis] Implemented in \harness. A residual proportional to
    $\omega_{\Pi}v_x$ means the inertial basis crept back in; see ROW-05.
\end{testcard}

\begin{testcard}{ROW-02}{The $b/2$ signs match the wheel geometry}{unit}{P0}
  \item[Targets] \eqref{eq:differential-drive-qp},
    \eqref{eq:differential-wheel-geometry}
  \item[Regime] \texttt{F-DIFF}.
  \item[Invariant] Driving the right wheel faster than the left must yield
    $\omega_{\Pi}>0$, matching
    \eqref{eq:differential-angular-velocity}.
  \item[Pass] $\omega_{\Pi}=(\wheelRadius_R\dot{\wheelAngle}_R
    -\wheelRadius_L\dot{\wheelAngle}_L)/b$ to $\testTol_{\text{alg}}$, sign
    included.
  \item[Diagnosis] A swapped sign inverts every turn while leaving speeds
    correct, so it passes any test that checks only $v_x$.
\end{testcard}

\begin{testcard}{ROW-03}{The measured residual decays at the discrete rate set by $\pGain$}{unit}{P0}
  \item[Targets] \eqref{eq:differential-drive-qp},
    \eqref{eq:geometric-wheel-acceleration-task}
  \item[Regime] Nonzero measured lateral velocity, $0<\pGain\samplingPeriod<2$.
  \item[Invariant] One cycle multiplies the lateral residual by exactly
    $(1-\pGain\samplingPeriod)$. At $\pGain=0$ the residual is frozen, not
    removed.
  \item[Pass] Per-cycle ratio within $10^{-9}$ of $1-\pGain\samplingPeriod$;
    the $\pGain=0$ branch must hold the residual to $\testTol_{\text{alg}}$.
  \item[Diagnosis] Implemented in \harness. Asserting a continuous-time decay
    $e^{-\pGain\samplingPeriod}$ instead is a false-failure trap.
\end{testcard}

\begin{testcard}{ROW-04}{Basis-convention pin on the rows themselves}{negative}{P0}
  \item[Targets] \eqref{eq:differential-drive-qp},
    Remark~\ref{re:planar-basis-choice}
  \item[Regime] Turning motion with $\omega_{\Pi}v_x\neq0$.
  \item[Invariant] The chassis-aligned rows are exact. The inertial-convention
    rows differ from them by exactly $-\omega_{\Pi}v_x$ in the lateral row and
    $+\omega_{\Pi}v_y$ in each rolling row.
  \item[Pass] Both branches asserted: the implemented rows match the first, and
    a deliberately mutated inertial assembler reproduces the second offset to
    convention~\ref{tol:fd}.
  \item[Diagnosis] Implemented in \harness. The two-branch form guarantees the
    test is discriminating rather than vacuous.
\end{testcard}

\begin{testcard}{ROW-05}{T2 rows are the exact derivative, with no $\dot{\text{H}}_i$ term}{unit}{P0}
  \item[Targets] \eqref{eq:four-steering-wheel-qp},
    \eqref{eq:steering-expanded-constraints}
  \item[Regime] Assumptions~\ref{as:planar-basis}--\ref{as:steering-axis}.
  \item[Invariant] The printed rolling and lateral rows equal the central
    difference of $\transpose{\rollingDirection^{\Pi}_i}\text{H}_i
    \bs{\xi}^{\Pi}_{\fbFrame}$ and its lateral counterpart. No
    $\dot{\text{H}}_i$ term appears, since $\dot{\text{H}}_i=\zeroMatrix$.
  \item[Pass] Convention~\ref{tol:fd}. Additionally assert that the term
    $-\omega_{\Pi}^2\innerP{\rollingDirection^{\Pi}_i}{\bs{\rho}_i}$, required
    under the inertial convention, is \emph{absent}.
  \item[Diagnosis] Implemented in \harness.
\end{testcard}

\begin{testcard}{ROW-06}{Direction derivatives and their preserved invariants}{unit}{P0}
  \item[Targets] \eqref{eq:steering-direction-derivatives}
  \item[Regime] T2, chassis-relative $\dot\delta_i$.
  \item[Invariant] $\dot{\rollingDirection}^{\Pi}_i
    =\dot\delta_i\lateralDirection^{\Pi}_i$ and
    $\dot{\lateralDirection}^{\Pi}_i
    =-\dot\delta_i\rollingDirection^{\Pi}_i$, with
    $\innerP{\rollingDirection^{\Pi}_i}{\dot{\rollingDirection}^{\Pi}_i}=0$ and
    $\innerP{\rollingDirection^{\Pi}_i}{\lateralDirection^{\Pi}_i}=0$ preserved
    along the motion.
  \item[Pass] Convention~\ref{tol:fd} against central differences; the
    orthogonality invariants to $\testTol_{\text{alg}}$.
  \item[Diagnosis] Implemented in \harness. A missing minus sign in
    $\dot{\lateralDirection}^{\Pi}_i$ is invisible at
    $\dot\delta_i=0$, so the fixture must steer.
\end{testcard}

\begin{testcard}{ROW-07}{Rate-prediction block is an exact affine substitution}{unit}{P1}
  \item[Targets] \eqref{eq:four-steering-rate-prediction},
    \eqref{eq:four-steering-decision-vector}
  \item[Regime] T2.
  \item[Invariant] Eliminating $\bs{\nu}_{\mathrm{rot}}^{\mathrm{4s},+}$ by
    substitution leaves the solution, the active set, and every bound margin
    unchanged; only the problem dimension drops by eight.
  \item[Pass] Solutions agree to $\testTol_{\text{alg}}$ on the shared
    variables; the active sets are identical as sets.
  \item[Diagnosis] A discrepancy means the prediction was applied with the
    wrong $\samplingPeriod$ or to the measured rather than the optimized rate.
\end{testcard}

\begin{testcard}{ROW-08}{Lateral-row redundancy for T1, structural rank for both}{unit}{P0}
  \item[Targets] \eqref{eq:differential-geometric-constraints},
    \eqref{eq:differential-drive-qp}
  \item[Regime] \texttt{F-DIFF}, $b>0$, $\wheelRadius_i>0$.
  \item[Invariant] Both wheels yield the identical lateral row, so only one
    copy belongs in the QP; and the printed $3\times5$ kinematic block has full
    row rank $3$ for every $b>0$ and $\wheelRadius_i>0$.
  \item[Pass] Row equality to $\testTol_{\text{alg}}$; rank exactly $3$ over
    $500$ draws from \texttt{F-RANDOM}.
  \item[Diagnosis] Implemented in \harness. Two copies are algebraically
    harmless but degrade conditioning and can confuse an active-set solver.
\end{testcard}

\begin{testcard}{ROW-09}{The lateral slacks are exactly the range-space defect}{property}{P1}
  \item[Targets] \eqref{eq:four-steering-wheel-qp},
    Remark~\ref{re:four-steering-lateral-slack}
  \item[Regime] T2, common $\weight_{\sigma}$, $\weight_{\sigma}$ large.
  \item[Invariant] As $\weight_{\sigma}\to\infty$,
    $\min\norm{\slackVariable^{\text{lat}}}\to\norm{\transpose{\bs{Z}}\bs{c}}$
    where $\bs{Z}$ spans $\text{null}(\transpose{\text{A}})$ and $\bs{c}$ is
    the measured right-hand side.
  \item[Pass] Agreement within $1\%$ at
    $\weight_{\sigma}/\weight_{v_x}=10^{6}$, with the residual decreasing
    monotonically in $\weight_{\sigma}$ over four decades.
  \item[Diagnosis] Confirms the slack reports a genuine kinematic
    incompatibility rather than absorbing an assembly error.
\end{testcard}

\begin{testcard}{ROW-10}{Degenerate steering configurations}{unit}{P1}
  \item[Targets] \eqref{eq:steering-expanded-constraints}
  \item[Regime] T2, \texttt{F-4S}.
  \item[Invariant] With all wheels tangent to circles about the base origin,
    only rotation is admissible and the admissible twist direction is
    $(0,0,1)$. With all steering angles equal, yaw is forced to zero and only
    translation survives.
  \item[Pass] The computed nullspace direction matches the expected one to
    $10^{-8}$ after normalization.
  \item[Diagnosis] These are the two configurations where an implementation
    that silently pseudo-inverts produces motion that the mechanism cannot
    execute.
\end{testcard}

\begin{testcard}{ROW-11}{T2 rows carry no proportional stabilization}{negative}{P2}
  \item[Targets] \eqref{eq:four-steering-wheel-qp},
    \eqref{eq:differential-drive-qp}
  \item[Regime] Both chassis.
  \item[Invariant] T1's rows carry $-\pGain$ times the measured residual; T2's
    printed rows do not. An implementation that reuses one assembler for both
    must not silently apply $\pGain$ to T2.
  \item[Pass] Setting $\pGain\neq0$ changes T1's right-hand side and leaves
    T2's unchanged.
  \item[Diagnosis] Documents a deliberate asymmetry so it is not ``fixed'' by a
    later refactor. If stabilization is wanted for T2 it must be added
    explicitly, alongside the slacks.
\end{testcard}

\begin{testcard}{ROW-12}{Steering-angle wrap-around}{regression}{P2}
  \item[Targets] \eqref{eq:steering-directions}
  \item[Regime] T2, $\delta_i$ crossing $\pm\pi$.
  \item[Invariant] The rows are $2\pi$-periodic in $\delta_i$ and change sign
    under $\delta_i\mapsto\delta_i+\pi$; $\dot\delta_i$ must be computed from an
    unwrapped angle.
  \item[Pass] Row values agree across the wrap to $\testTol_{\text{alg}}$; the
    finite-differenced $\dot\delta_i$ shows no spike at the crossing.
  \item[Diagnosis] A wrap spike injects a large spurious
    $\dot{\rollingDirection}^{\Pi}_i$ for one cycle, which reads as a transient
    slip.
\end{testcard}

\section{Layer C. Dynamics, equilibrium and momentum}
\label{sec:test-layer-c}

This layer contains the corrected forms of the requester's claims (1) and (3),
together with the tests that protect the congruence reduction. It needs a
full-model dynamics oracle, for example a recursive Newton--Euler or composite
rigid-body implementation of \eqref{eq:rolling-equations-of-motion}.

\begin{testcard}{DYN-01}{Reduced inertia is a congruence, not a row slice}{unit}{P0}
  \item[Targets] \eqref{eq:reduced-dynamics-congruence}
  \item[Regime] Both chassis; full model available.
  \item[Invariant] $\text{M}_{\text{red}}=\transpose{\text{P}}\jsim\text{P}$.
    In particular $\text{M}_{\text{red}}$ is symmetric positive definite, which
    a row selection of $\jsim$ is not.
  \item[Pass] Blockwise agreement to $10^{-9}$ relative to
    $\norm{\jsim}$; symmetry to $\testTol_{\text{alg}}$; smallest eigenvalue
    strictly positive.
  \item[Diagnosis] A row-slicing implementation fails the symmetry assertion
    immediately, which is why symmetry is the primary check rather than a
    secondary one.
\end{testcard}

\begin{testcard}{DYN-02}{$\bs{h}_{\text{red}}$ carries the $\dot{\text{P}}\bs{u}$ term}{unit}{P0}
  \item[Targets] \eqref{eq:reduced-dynamics-congruence},
    Remark~\ref{re:planar-basis-choice}
  \item[Regime] Turning motion, so that $\dot{\text{P}}\bs{u}\neq\zeroVector$.
  \item[Invariant] $\bs{h}_{\text{red}}=\transpose{\text{P}}
    (\jsim\dot{\text{P}}\bs{u}+\gravityandcoriolis)$. The velocity-product part
    is not optional: it is where the Coriolis and centripetal terms of the
    rotating chassis-aligned basis live.
  \item[Pass] Compare against the full model at nonzero $\omega_{\Pi}$;
    agreement to $10^{-4}$ of $\mass\gravity[value]$,
    convention~\ref{tol:phys}. Separately assert that dropping
    $\dot{\text{P}}\bs{u}$ changes $\bs{h}_{\text{red}}$ by more than that
    tolerance, so the test is discriminating.
  \item[Diagnosis] Omitting the term makes the controller correct in a straight
    line and wrong in every turn, in proportion to $\omega_{\Pi}^2$.
\end{testcard}

\begin{testcard}{DYN-03}{Correction of claim (1): Newton's law, general and static}{property}{P0}
  \item[Targets] \eqref{eq:test-newton-general},
    \eqref{eq:rolling-equations-of-motion}
  \item[Regime] Stated per branch; see \secRef{sec:correction-gravity}.
  \item[Invariant] General:
    $\agg{i}{\numContacts}{\force_i}=\mass(\com[ddot]-\gravity)$.
    Static: $\agg{i}{\numContacts}{\force_i}=-\mass\gravity$. In the planar
    QPs only the two in-plane components are earned by the solver; the
    surface-normal component is supplied by
    \eqref{eq:planar-normal-load-row} and is therefore a test of the load
    model.
  \item[Pass] In-plane residual below $10^{-4}\mass\gravity[value]$; the
    normal component asserted separately against the load model, not against
    the QP.
  \item[Diagnosis] Implemented in \harness. Asserting the naive
    ``forces equal gravity'' during acceleration produces a false failure of
    magnitude $\mass\norm{\com[ddot]}$.
\end{testcard}

\begin{testcard}{DYN-04}{Correction of claim (3): yaw momentum equals the contact moment}{property}{P0}
  \item[Targets] \eqref{eq:test-yaw-balance},
    \eqref{eq:reduced-dynamics-congruence}
  \item[Regime] Both chassis, any motion.
  \item[Invariant] $\innerP{\surfaceNormal_{\Pi}}{\angularMomentum[dot][\comSymbol]}$
    equals the net plane-normal moment of the contact forces about the center
    of mass. Actuator torques contribute nothing.
  \item[Pass] Residual below $10^{-4}\mass\gravity[value]\ell$ with $\ell$ the
    largest carrier offset, convention~\ref{tol:phys}.
  \item[Diagnosis] Implemented in \harness. A residual proportional to
    $\sum_i\jtorques_{\delta_i}$ is the signature of the defect DYN-05 targets.
\end{testcard}

\begin{testcard}{DYN-05}{Steering torque cannot enter the chassis yaw row}{negative}{P0}
  \item[Targets] \eqref{eq:reduced-dynamics-congruence},
    \secRef{sec:correction-yaw}
  \item[Regime] Both chassis.
  \item[Invariant] $\text{B}_{\text{red}}$ has identically zero rows for the
    three floating-base coordinates; in particular the plane-normal base row is
    exactly zero.
  \item[Pass] Exact equality to zero, not a threshold. A mutant that adds
    $-\jtorques_{\delta_i}$ to the yaw row must fail both this test and DYN-04.
  \item[Diagnosis] Implemented in \harness. This is the structural form of the
    correction: it fails on inspection of the matrix, without any simulation,
    and it catches the error that DYN-04 alone might absorb into tolerance.
\end{testcard}

\begin{testcard}{DYN-06}{Contact power vanishes under ideal rolling}{property}{P1}
  \item[Targets] \eqref{eq:ideal-rolling-velocity},
    \eqref{eq:reduced-dynamics-congruence}
  \item[Regime] Ideal rolling, no slip, no sliding contacts.
  \item[Invariant] The material contact point has zero velocity, so each
    contact force does no work:
    $\transpose{\bs{u}}\transpose{\text{J}_{\text{red},i}}\force_i=0$.
    Consequently actuator power equals the rate of mechanical energy.
  \item[Pass] Contact power below $10^{-6}$ of the actuator power; the energy
    balance closes to $10^{-4}$ relative over a $5$\unitTime rollout.
  \item[Diagnosis] A strong independent oracle: it uses no constraint row
    directly, so it catches sign and frame errors that the row-level tests
    share a convention with.
\end{testcard}

\begin{testcard}{DYN-07}{Normal-load model closes the three discarded equilibrium rows}{unit}{P0}
  \item[Targets] \eqref{eq:planar-normal-load-row},
    \eqref{eq:reduced-dynamics-congruence}
  \item[Regime] Assumption~\ref{as:normal-load}; \texttt{F-DIFF} with a caster,
    \texttt{F-4S} with a load model.
  \item[Invariant] The supplied $\reff{\force}_{n,i}$ must satisfy the three
    out-of-plane equations discarded by the reduction: normal force balance and
    the two tilting-moment balances about the in-plane axes, including the
    inertial load transfer $\mass\com[ddot]$ under acceleration.
  \item[Pass] Residual below $10^{-3}\mass\gravity[value]$ statically and under
    a $2$\unitAccTS longitudinal acceleration.
  \item[Diagnosis] For T1 the system is over-determined, so a two-wheel model
    without a caster share cannot close and must be reported as such rather
    than least-squares fitted silently. For T2 it is indeterminate, so the test
    checks the chosen distribution rule, not uniqueness.
\end{testcard}

\begin{testcard}{DYN-08}{Sign conventions, asserted individually}{unit}{P0}
  \item[Targets] \secRef{sec:rolling-implementation},
    \eqref{eq:wheel-torque-row}
  \item[Regime] Both chassis.
  \item[Invariant] Positive $\dot{\wheelAngle}_i$ produces positive rim speed
    along $\rollingDirection_i$; positive $\jtorques_{\wheelAngle_i}$ produces
    positive $\ddot{\wheelAngle}_i$ in isolation; and the drive torque enters
    the generalized force with the sign of
    \eqref{eq:rolling-equations-of-motion}.
  \item[Pass] Each of the three asserted as a separate pass/fail so a failure
    localizes.
  \item[Diagnosis] A pair of compensating sign errors reproduces correct
    straight-line motion, so bundling these three into one assertion hides the
    common case.
\end{testcard}

\begin{testcard}{DYN-09}{Quasi-static torque bound must not duplicate the wheel row}{negative}{P1}
  \item[Targets] \eqref{eq:quasistatic-wheel-force},
    \eqref{eq:rolling-equations-of-motion}
  \item[Regime] Both chassis.
  \item[Invariant] Imposing
    $\abs{\innerP{\rollingDirection_i}{\force_i}}
     \leq\abs{\jtorques_{\wheelAngle_i}}/\wheelRadius_i$ alongside the complete
    wheel row must be detected as duplication: it removes feasible points that
    the full dynamics admits, because it neglects the wheel's rotational
    inertia.
  \item[Pass] Construct a state with nonzero $\ddot{\wheelAngle}_i$ for which
    the full row is satisfied and the quasi-static inequality is violated;
    assert the QP with both constraints is infeasible and the QP with the row
    alone is feasible.
  \item[Diagnosis] Confirms the warning after
    \eqref{eq:quasistatic-wheel-force} is respected.
\end{testcard}

\begin{testcard}{DYN-10}{Momentum-integral rollout}{regression}{P2}
  \item[Targets] \eqref{eq:test-newton-general}, \eqref{eq:test-yaw-balance}
  \item[Regime] Closed-loop, $5$\unitTime, both terrains.
  \item[Invariant] The discrete change in linear and plane-normal angular
    momentum matches the integrated contact wrench, up to the integrator's
    order.
  \item[Pass] Relative error below $10^{-3}$ and decreasing at the expected
    order as $\samplingPeriod$ is halved.
  \item[Diagnosis] Turns the two instantaneous balances into a trajectory-level
    statement, catching a defect that is small per cycle but accumulates.
\end{testcard}

\section{Layer D. Contact force and friction}
\label{sec:test-layer-d}

The corrected form of the requester's claim (2) lives here. The organizing
principle is that cone membership is a tautology, so the tests target the force
\emph{path}, the cone \emph{rebuild}, and whether the cone \emph{binds}.

\begin{testcard}{FRI-01}{Cone membership is a tautology: test the path, not the value}{unit}{P1}
  \item[Targets] \eqref{eq:rolling-force-generators},
    \secRef{sec:correction-cone}
  \item[Regime] Both chassis.
  \item[Invariant] Recomputing
    $\norm{(\force_{t,i},\force_{s,i})}\leq\coef{f}_i\force_{n,i}$ from
    $\force_i=\cone[f][i]\generator{\force_i}$ with
    $\generator{\force_i}\geq\zeroVector$ cannot fail. The test must therefore
    assert instead that the checker is \emph{sensitive}: injecting a force that
    is not of the form $\cone[f][i]\generator{\force_i}$ must be rejected.
  \item[Pass] The tautological branch passes trivially; the injected
    out-of-cone force must be flagged.
  \item[Diagnosis] Implemented in \harness. Without the sensitivity branch this
    entry is a test that cannot fail, which is worse than no test.
\end{testcard}

\begin{testcard}{FRI-02}{The cone is rebuilt when the wheel steers}{negative}{P0}
  \item[Targets] \eqref{eq:rolling-force-generators},
    \eqref{eq:steering-directions}
  \item[Regime] T2, nonzero $\delta_i$.
  \item[Invariant] $\cone[f][i]$ must be expressed in the current local frame
    $[\rollingDirection_i,\lateralDirection_i,\surfaceNormal_i]$. Because the
    cone \emph{set} is rotationally symmetric about $\surfaceNormal_i$, the
    stale-cone defect is invisible in the set; it is visible only in the
    \emph{map} $\generator{\force_i}\mapsto\force_i$.
  \item[Pass] Assert exact rotational equivariance:
    $\cone[f][i](\delta_i)=\text{R}(\delta_i)\cone[f][i](0)$ to
    $\testTol_{\text{alg}}$. A cone cached at $\delta_i=0$ must fail for
    $\delta_i=0.6$~rad of steering.
  \item[Diagnosis] The classic defect. Testing only cone membership will never
    catch it, which is why this entry asserts the map.
\end{testcard}

\begin{testcard}{FRI-03}{The cone is rebuilt when the terrain normal changes}{negative}{P1}
  \item[Targets] \eqref{eq:rolling-force-generators}, \texttt{F-RAMP}
  \item[Regime] Ramp, or a terrain-normal update.
  \item[Invariant] Each generator satisfies
    $\innerP{\surfaceNormal_i}{\bs{c}_k}=1$ in the current frame. A stale
    normal breaks this at first order in the tilt.
  \item[Pass] Deviation below $\testTol_{\text{alg}}$ for the fresh cone; the
    stale-normal mutant must deviate by at least $1-\cos\theta$ at
    $\theta=20$\unitDegree.
  \item[Diagnosis] Distinguishes a normal-update fault from a steering-update
    fault, which FRI-02 cannot.
\end{testcard}

\begin{testcard}{FRI-04}{Inner approximation and its quantified conservatism}{property}{P1}
  \item[Targets] \eqref{eq:rolling-friction-cone},
    \eqref{eq:rolling-force-generators}
  \item[Regime] Any $\numGen\geq3$.
  \item[Invariant] Every representable force satisfies
    $\norm{(\force_{t,i},\force_{s,i})}\leq\coef{f}_i\force_{n,i}$, and the
    attained inradius is exactly $\coef{f}_i\cos(\pi/\numGen)$: tangential
    demands between $\coef{f}_i\cos(\pi/\numGen)\force_{n,i}$ and
    $\coef{f}_i\force_{n,i}$ are rejected at unfavourable headings.
  \item[Pass] Both bounds attained to $10^{-9}$; the two-sided form is the
    assertion, since the upper bound alone is implied by construction.
  \item[Diagnosis] Implemented in \harness. Fixes the achievable traction
    budget, which the ramp tests depend on.
\end{testcard}

\begin{testcard}{FRI-05}{The normal-load row makes the cone bind}{smoke}{P0}
  \item[Targets] \eqref{eq:planar-normal-load-row},
    \eqref{eq:reduced-normal-annihilated}
  \item[Regime] Both chassis, corrected formulation.
  \item[Invariant] Without \eqref{eq:planar-normal-load-row} the feasible set
    contains the ray $\generator{\force_i}\mapsto
    \generator{\force_i}+\text{s}\bs{1}$, $\text{s}\geq0$, on which every
    constraint holds and the objective is unchanged; the cone never binds and
    the optimal set is unbounded. With the row, the generator set is a compact
    simplex and the ray is blocked.
  \item[Pass] Two branches. Removing the row must reproduce an unbounded or
    solver-dependent $\generator{\force_i}$; with the row, the feasible
    generator set must be bounded and the cone must be attained at a commanded
    traction limit.
  \item[Diagnosis] Implemented in \harness. This is the highest-severity defect
    the corrections addressed; keep both branches so a regression is caught.
\end{testcard}

\begin{testcard}{FRI-06}{Normal-load input validation and detachment}{negative}{P1}
  \item[Targets] \eqref{eq:planar-normal-load-row}
  \item[Regime] Both chassis.
  \item[Invariant] $\reff{\force}_{n,i}<0$ must be rejected at assembly.
    $\reff{\force}_{n,i}=0$ forces $\generator{\force_i}=\zeroVector$ and hence
    zero tangential force, which is detachment and must be handled by the
    contact-mode logic rather than by the QP.
  \item[Pass] Negative input raises; zero input yields exactly zero tangential
    force and a mode-manager notification.
  \item[Diagnosis] A silently clamped negative load produces a wheel that
    generates traction while airborne.
\end{testcard}

\begin{testcard}{FRI-07}{Tension is unrepresentable, so lift-off must be guarded outside the QP}{negative}{P1}
  \item[Targets] \eqref{eq:rolling-force-generators},
    \eqref{eq:reduced-normal-annihilated}
  \item[Regime] Both chassis.
  \item[Invariant] $\force_{n,i}\geq0$ holds by construction, and the planar
    reduction cannot express tension at all. Therefore a manoeuvre that would
    physically lift a wheel produces no QP infeasibility and no warning.
  \item[Pass] Command a load transfer that drives $\reff{\force}_{n,i}$ to
    zero; assert the QP still solves and that an external monitor raises the
    lift-off condition.
  \item[Diagnosis] Prevents the false belief that QP feasibility certifies
    contact maintenance in the planar model. It does not.
\end{testcard}

\begin{testcard}{FRI-08}{Ramp station-keeping and the heading-dependent failure slope}{smoke}{P1}
  \item[Targets] \eqref{eq:test-ramp-split}, \texttt{F-RAMP}
  \item[Regime] Static hold on slopes $0$ to $40$\unitDegree.
  \item[Invariant] The normal and tangential sums follow
    \eqref{eq:test-ramp-split}, and the controller fails between
    $\arctan(\coef{f}\cos(\pi/\numGen))$ and $\arctan\coef{f}$ depending on the
    heading relative to the fall line, because the polyhedral cone is
    anisotropic.
  \item[Pass] Force split within $10^{-3}\mass\gravity[value]$; the observed
    failure slope inside the predicted band for every heading tested.
  \item[Diagnosis] Asserting a single failure slope $\arctan\coef{f}$ is a
    false-failure trap at unfavourable headings; the band is the correct
    assertion.
\end{testcard}

\begin{testcard}{FRI-09}{Friction margin: log both margins and assert their ordering}{unit}{P2}
  \item[Targets] \secRef{sec:rolling-cases},
    \eqref{eq:rolling-friction-cone}
  \item[Regime] Both chassis.
  \item[Invariant] The circular margin
    $\coef{f}_i\force_{n,i}-\norm{(\force_{t,i},\force_{s,i})}$ and the
    polyhedral margin are both recorded, and the polyhedral margin is never
    larger.
  \item[Pass] Ordering holds on every cycle of a $5$\unitTime rollout; both are
    reported in the metrics record.
  \item[Diagnosis] Reporting only one margin makes a near-saturation event
    ambiguous between model conservatism and genuine physical limit.
\end{testcard}

\begin{testcard}{FRI-10}{Sliding-direction row is not imposed while sticking or rolling}{negative}{P1}
  \item[Targets] \eqref{eq:sliding-friction-direction}
  \item[Regime] Contact modes $\set{r}$ and $\set{f}$.
  \item[Invariant] The affine direction row is applied only for sustained
    sliding. During rolling or sticking, tangential force must remain an
    interior decision of the cone.
  \item[Pass] With all contacts rolling, assert the row is absent from the
    assembled problem and that the solution's tangential force is strictly
    interior for a non-saturating command.
  \item[Diagnosis] Imposing it during rolling pins the force to the cone
    boundary, producing a controller that always saturates friction. Note also
    that under a polyhedral inner approximation the row is not automatically
    cone compatible, so its $\varepsilon$ guard must be tested.
\end{testcard}

\section{Layer E. Bounds, actuation and prediction}
\label{sec:test-layer-e}

\begin{testcard}{BND-01}{Rate limits act on the one-cycle-ahead rate}{unit}{P0}
  \item[Targets] \eqref{eq:four-steering-rate-prediction},
    \eqref{eq:differential-drive-qp}, \eqref{eq:four-steering-wheel-qp}
  \item[Regime] Both chassis.
  \item[Invariant] The rate box is imposed on
    $\dot{\wheelAngle}_i+\samplingPeriod\ddot{\wheelAngle}_i$, not on the
    measured rate. The assembled row is therefore $[\,1\ \ \samplingPeriod\,]$
    against the measured rate, and the box scales with $\samplingPeriod$.
  \item[Pass] Halving $\samplingPeriod$ must halve the reachable acceleration
    window implied by the rate box, exactly.
  \item[Diagnosis] Bounding the measured rate makes the constraint a constant,
    so it never restricts the decision variable and silently disappears.
\end{testcard}

\begin{testcard}{BND-02}{A yawing chassis with locked steering consumes no steering-rate budget}{negative}{P0}
  \item[Targets] \eqref{eq:four-steering-wheel-qp},
    \eqref{eq:steering-direction-derivatives}
  \item[Regime] T2, $\omega_{\Pi}=1$~rad/s, steering modules locked.
  \item[Invariant] Because $\delta_i$ is chassis relative, the steering-rate
    bound is independent of $\omega_{\Pi}$: the assembled bound row must have a
    zero coefficient on $\dot\omega_{\Pi}$.
  \item[Pass] Exact zero coefficient. A mutant using the absolute heading rate
    must show a coefficient of $1$ and must consume the full rate budget at
    $\omega_{\Pi}=1$~rad/s.
  \item[Diagnosis] The paired mutant is what makes this test discriminating;
    without it the assertion is satisfied by an implementation that has no
    steering bound at all.
\end{testcard}

\begin{testcard}{BND-03}{Rate box and acceleration box become jointly empty at a known threshold}{unit}{P1}
  \item[Targets] \eqref{eq:four-steering-wheel-qp}
  \item[Regime] Measured rate already outside, or near, its bound.
  \item[Invariant] The two boxes are jointly infeasible exactly when
    $\dot{\wheelAngle}_i+\samplingPeriod\ddot{\wheelAngle}_{i,\text{max}}
    <\dot{\wheelAngle}_{i,\text{min}}$ or the mirrored condition holds. The
    implementation must report this rather than clip.
  \item[Pass] Sweep the measured rate across the analytic threshold; the
    transition to reported infeasibility must occur within one step of it.
  \item[Diagnosis] The report explicitly forbids clipping. A clipping
    implementation passes every steady-state test and fails only on recovery
    from a disturbance.
\end{testcard}

\begin{testcard}{BND-04}{Each limit binds at its own analytic threshold, in the right order}{smoke}{P1}
  \item[Targets] \eqref{eq:differential-drive-qp},
    \eqref{eq:wheel-torque-bound}, \eqref{eq:rolling-friction-cone}
  \item[Regime] \texttt{F-DIFF}, reference ramped monotonically.
  \item[Invariant] As the commanded acceleration is increased, the rate,
    acceleration, torque and friction limits activate at thresholds computable
    in closed form from the fixture, and in a predictable order.
  \item[Pass] Each activation within $2\%$ of its analytic threshold, and the
    order matches the prediction.
  \item[Diagnosis] Localizes a scaling error to a single limit. A uniformly
    shifted set of thresholds indicates a units error in the fixture rather
    than in the bounds.
\end{testcard}

\begin{testcard}{BND-05}{Torque bounds are selected at the right stride within each wheel block}{unit}{P1}
  \item[Targets] \eqref{eq:steering-wheel-decision-block},
    \eqref{eq:four-steering-torque-vector}
  \item[Regime] T2.
  \item[Invariant] The drive-torque limit must attach to
    $\jtorques_{\wheelAngle_i}$ and the steering-torque limit to
    $\jtorques_{\delta_i}$, at the correct offset inside the four-entry block.
  \item[Pass] Set each of the eight torque limits to a distinct sentinel value
    and read them back from the assembled bound vector by index.
  \item[Diagnosis] An off-by-one stride swaps drive and steering limits, which
    is invisible when the two happen to be similar and catastrophic when they
    are not.
\end{testcard}

\begin{testcard}{BND-06}{The lateral slacks carry no box}{negative}{P0}
  \item[Targets] \eqref{eq:four-steering-wheel-qp},
    Remark~\ref{re:four-steering-lateral-slack}
  \item[Regime] T2.
  \item[Invariant] $\slackVariable^{\text{lat}}$ must be unbounded. Bounding it
    re-imposes the generically infeasible hard rows whenever the required slack
    exceeds the bound.
  \item[Pass] Assert no finite bound is attached to the slack columns; then
    supply data requiring a large slack and confirm the QP still solves.
  \item[Diagnosis] A well-meaning ``sanity bound'' on the slack reintroduces
    exactly the defect the slack was added to remove.
\end{testcard}

\begin{testcard}{BND-07}{Normal-load row requires a strictly positive right-hand side}{unit}{P1}
  \item[Targets] \eqref{eq:planar-normal-load-row}
  \item[Regime] Both chassis.
  \item[Invariant] The row is an equality on a nonnegative combination, so
    $\reff{\force}_{n,i}>0$ is required for a nonempty generator set with
    interior; $\reff{\force}_{n,i}=0$ collapses it to the origin.
  \item[Pass] Positive values give a bounded nonempty simplex; zero gives
    exactly the origin; negative is rejected.
  \item[Diagnosis] Complements FRI-06 at the level of the assembled row rather
    than the model input.
\end{testcard}

\begin{testcard}{BND-08}{Bounds are recomputed every cycle and dimensions are constant}{regression}{P1}
  \item[Targets] \secRef{sec:rolling-implementation}
  \item[Regime] Closed loop across a contact-mode change.
  \item[Invariant] The right-hand sides of all bounds are refreshed from the
    current measurement each cycle, while the problem dimension and the index
    map are unchanged.
  \item[Pass] Dimensions and index maps bit-identical across $10^{3}$ cycles
    including a mode change; bound values demonstrably tracking the
    measurement.
  \item[Diagnosis] A stale bound produces a controller that is correct on the
    first cycle after a reset and progressively wrong afterwards.
\end{testcard}

\section{Layer F. QP well-posedness and solver contract}
\label{sec:test-layer-f}

Throughout this layer the objective is written
$q(\decisionVar)=\transpose{\decisionVar}\weightMatrix{H}\decisionVar
+\transpose{\bs{g}}\decisionVar$, with no factor of one half. Every numeric
assertion on $\weightMatrix{H}$ below assumes that convention; if the
implementation uses $\tfrac12\transpose{\decisionVar}\weightMatrix{H}\decisionVar$
then the expected Hessian entries double. Pin this once in the test fixture and
assert it, because two entries disagreeing about the factor is a defect that
looks like a tolerance failure.

\begin{testcard}{QP-01}{Hessian is positive definite under the corrected objective}{unit}{P0}
  \item[Targets] \eqref{eq:planar-force-regularization},
    \eqref{eq:planar-twist-weight}
  \item[Regime] $\varepsilon_{\lambda}>0$, $\weightMatrix{W}_{\xi}\succ0$, all
    motion weights positive.
  \item[Invariant] $\weightMatrix{H}\succ0$ on the whole decision space,
    including the generator block. Maintain a per-block ledger so a failure
    identifies which block is singular.
  \item[Pass] Smallest eigenvalue exceeds
    $\tfrac12\varepsilon_{\lambda}$; the ledger accounts for every column.
  \item[Diagnosis] Implemented in \harness.
\end{testcard}

\begin{testcard}{QP-02}{Setting $\varepsilon_{\lambda}=0$ reproduces the singular case}{regression}{P0}
  \item[Targets] \eqref{eq:planar-force-regularization},
    \eqref{eq:reduced-normal-annihilated}
  \item[Regime] Both chassis, $\numGen>3$.
  \item[Invariant] With $\varepsilon_{\lambda}=0$ and no load-distribution
    term, the generator block has nullity exactly $\numGen-3$ per contact and
    the minimizer is not unique.
  \item[Pass] Measured nullity equals $\numGen-3$; at $\numGen=3$ the
    degeneracy disappears, which is the boundary case worth pinning.
  \item[Diagnosis] Implemented in \harness. Keeping the singular branch as an
    explicit expectation prevents a later ``simplification'' from silently
    removing the regularization.
\end{testcard}

\begin{testcard}{QP-03}{Primal solution is unique across warm starts and solvers}{property}{P0}
  \item[Targets] \eqref{eq:planar-force-regularization}
  \item[Regime] Corrected objective, $\varepsilon_{\lambda}>0$.
  \item[Invariant] Solving from several distinct warm starts, and with two
    independent solvers, yields the same primal solution.
  \item[Pass] Agreement to $10^{-6}$ relative on
    $\dot{\bs{\xi}}^{\Pi}_{\fbFrame}$ and to
    $\sqrt{\testTol_{\text{alg}}/\varepsilon_{\lambda}}$ on
    $\generator{\force_i}$, since the generator block is only
    $\varepsilon_{\lambda}$-strongly convex. Compare duals only where strict
    complementarity holds.
  \item[Diagnosis] The $\varepsilon_{\lambda}$-aware tolerance is essential: a
    uniform tight tolerance on $\generator{\force_i}$ is a false-failure trap.
\end{testcard}

\begin{testcard}{QP-04}{Equality-row ledger and rank}{unit}{P0}
  \item[Targets] \eqref{eq:differential-drive-qp},
    \eqref{eq:four-steering-wheel-qp}
  \item[Regime] Both chassis.
  \item[Invariant] The equality rows number exactly as counted from the
    formulation, and the equality matrix has full row rank for every fixture
    with $b>0$, $\wheelRadius_i>0$ and, for T2, with the slacks present.
  \item[Pass] Exact counts; rank equal to the row count over $500$ draws from
    \texttt{F-RANDOM}, after nondimensionalization.
  \item[Diagnosis] A rank deficiency here is a duplicated row, most likely the
    second lateral copy in T1 or a reinstated $\text{G}^{\text{rot}}_i$
    identity block in T2.
\end{testcard}

\begin{testcard}{QP-05}{No lateral row can appear in a T2 infeasibility certificate}{smoke}{P0}
  \item[Targets] \eqref{eq:four-steering-wheel-qp},
    Remark~\ref{re:four-steering-lateral-slack}
  \item[Regime] T2, arbitrary measured steering angles and rates.
  \item[Invariant] Because each lateral row owns its own slack column, no
    lateral row can be part of a minimal infeasible subsystem. Any residual
    infeasibility must come from the bounds, the dynamics, or the cone.
  \item[Pass] Over $500$ adversarial draws including deliberately
    non-concurrent steering, the QP always solves; when infeasibility is forced
    by bounds, the returned certificate contains no lateral row index.
  \item[Diagnosis] Directly regression-tests the defect the slacks were added
    to remove.
\end{testcard}

\begin{testcard}{QP-06}{Unit rescaling leaves the physical solution invariant}{property}{P1}
  \item[Targets] \eqref{eq:planar-twist-weight},
    \secRef{sec:rolling-implementation}
  \item[Regime] Both chassis.
  \item[Invariant] Re-expressing lengths in millimetres and forces in
    kilonewtons, with weights transformed consistently, must leave the physical
    solution and the active set unchanged.
  \item[Pass] Physical quantities agree to $10^{-6}$ relative; the active set is
    identical as a set.
  \item[Diagnosis] A failure means a weight is a bare scalar where a matrix is
    required, which is exactly the defect
    \eqref{eq:planar-twist-weight} was introduced to fix.
\end{testcard}

\begin{testcard}{QP-07}{KKT residual as a solver-independent optimality oracle}{property}{P1}
  \item[Targets] both QPs
  \item[Regime] Any feasible fixture.
  \item[Invariant] Stationarity, primal and dual feasibility, and complementary
    slackness hold, computed from an \emph{independently} re-derived
    $\weightMatrix{H}$, $\bs{g}$ and constraint set rather than from the
    solver's own internal matrices.
  \item[Pass] Each residual below $10^{-8}$ after normalization by the
    corresponding row norm.
  \item[Diagnosis] Independent re-derivation is what makes this an oracle
    rather than a self-consistency check.
\end{testcard}

\begin{testcard}{QP-08}{Stability window of the proportional stabilization}{property}{P1}
  \item[Targets] \eqref{eq:differential-drive-qp},
    \eqref{eq:geometric-wheel-acceleration-task}
  \item[Regime] T1, discrete time.
  \item[Invariant] The residual recursion is stable exactly for
    $0\leq\pGain\samplingPeriod<2$, and monotone for
    $\pGain\samplingPeriod\leq1$. Above $2$ it diverges.
  \item[Pass] Sweep $\pGain\samplingPeriod$ over $[0,2.5]$; observed decay,
    oscillatory decay and divergence must match the predicted regions.
  \item[Diagnosis] Implemented in \harness. Fixes an upper limit on $\pGain$
    that is otherwise chosen by trial and error.
\end{testcard}

\begin{testcard}{QP-09}{Runtime contract}{smoke}{P1}
  \item[Targets] \secRef{sec:rolling-implementation}
  \item[Regime] Closed loop with mode changes.
  \item[Invariant] Problem dimensions and index maps are constant; the solver
    is deterministic given identical input; warm starting reproduces the cold
    solution; and the solve time stays within budget.
  \item[Pass] Bit-identical dimensions; cold and warm solutions agree to
    $10^{-8}$; the $99$th-percentile solve time under the declared budget.
  \item[Diagnosis] Determinism must be asserted before any golden-file test is
    meaningful.
\end{testcard}

\begin{testcard}{QP-10}{Rejection contract}{negative}{P0}
  \item[Targets] \secRef{sec:rolling-implementation}
  \item[Regime] Deliberately infeasible or malformed input.
  \item[Invariant] Certified infeasibility triggers the documented safe
    fallback; no non-finite value ever reaches an actuator; invalid model data
    is rejected at assembly rather than propagated.
  \item[Pass] Inject a non-finite bound, a negative $\reff{\force}_{n,i}$, a
    zero wheel radius and a non-positive-definite weight; each must be rejected
    with a distinguishable error, and no command must be emitted.
  \item[Diagnosis] Accepting an infinite bound is a known defect class in this
    codebase; it must fail validation, not pass through.
\end{testcard}

\section{Layer G. Independent oracles and property tests}
\label{sec:test-layer-g}

An oracle is only useful if it does not share a convention with the code under
test. Each entry below therefore states what makes it independent.

\begin{testcard}{ORC-01}{Richardson finite differences with a convergence-rate assertion}{unit}{P0}
  \item[Targets] \eqref{eq:steering-direction-derivatives},
    \eqref{eq:reduced-dynamics-congruence}
  \item[Regime] Every hand-derived time derivative in the formulation.
  \item[Invariant] The analytic derivative matches the central difference, and
    halving $h$ reduces the error by a factor of four.
  \item[Pass] Convention~\ref{tol:fd}; the measured convergence order must lie
    in $[1.8,\,2.2]$ over three refinements. Choose the fixture so the quantity
    differentiated is $O(1)$, otherwise the rate estimate is dominated by
    round-off.
  \item[Diagnosis] The rate assertion is the real content: a threshold alone is
    passed by a wrong derivative whose error happens to be small.
\end{testcard}

\begin{testcard}{ORC-02}{Reduced planar QP against its full-coordinate preimage}{smoke}{P0}
  \item[Targets] \eqref{eq:reduced-dynamics-congruence},
    \eqref{eq:ideal-rolling-whole-body-qp}
  \item[Regime] A planar fixture that the full whole-body QP can also express.
  \item[Invariant] Lifting the reduced solution through $\text{P}$ must yield a
    point that the full QP also returns, on the shared variables. Because the
    full problem has additional nullspace, ``agree'' means: the lifted reduced
    solution is feasible for the full problem, and its objective value matches
    the full optimum to within the regularization gap.
  \item[Pass] Shared variables agree to $10^{-6}$ relative; the feasibility
    residual of the lifted point is below $10^{-8}$.
  \item[Diagnosis] The single most valuable oracle in the suite: it validates
    the entire reduction in one assertion. Failure isolates to $\text{P}$,
    $\text{M}_{\text{red}}$ or $\bs{h}_{\text{red}}$ via DYN-01 and DYN-02.
\end{testcard}

\begin{testcard}{ORC-03}{Closed-form differential-drive oracle on the solved prediction}{smoke}{P1}
  \item[Targets] \eqref{eq:differential-linear-velocity},
    \eqref{eq:differential-angular-velocity}
  \item[Regime] T1, ideal rolling, $\pGain$ arbitrary.
  \item[Invariant] The solved one-cycle predicted twist satisfies the closed
    forms, with the wheel rates taken from the solution rather than from the
    command.
  \item[Pass] $10^{-3}$\unitVelTS on $v_x$ and the corresponding angular
    tolerance on $\omega_{\Pi}$, convention~\ref{tol:phys}.
  \item[Diagnosis] Using the commanded rather than the solved rates makes this
    a tautology; the distinction is the whole point.
\end{testcard}

\begin{testcard}{ORC-04}{Known-ICR analytic oracle for T2}{smoke}{P1}
  \item[Targets] \eqref{eq:steering-expanded-constraints}, \texttt{F-ICR}
  \item[Regime] T2, steering angles constructed from a prescribed ICR.
  \item[Invariant] The solved twist lies along the analytically known
    admissible direction, and $\norm{\slackVariable^{\text{lat}}}$ is at the
    soft-lateral minimum rather than at an arbitrary value.
  \item[Pass] Direction agreement to $10^{-6}$ after normalization; slack at
    its predicted minimum to $1\%$.
  \item[Diagnosis] Makes the soft-lateral limit the primary assertion, so the
    test remains meaningful for finite $\weight_{\sigma}$.
\end{testcard}

\begin{testcard}{ORC-05}{Metamorphic relations}{property}{P1}
  \item[Targets] both QPs
  \item[Regime] \texttt{F-RANDOM}.
  \item[Invariant] Three relations must hold for every draw. Rotating the whole
    scene about $\surfaceNormal_{\Pi}$ leaves the chassis-frame solution
    unchanged. Scaling every $\wheelRadius_i$ by $\alpha$ and every wheel rate
    by $1/\alpha$ leaves the chassis twist unchanged, exactly on the kinematic
    rows and up to the changed wheel inertia in the dynamics. Mirroring left
    and right mirrors the T1 solution.
  \item[Pass] $500$ draws, seeded and recorded; each relation to $10^{-6}$
    relative. The radius-scaling relation must be asserted only on the
    kinematic rows unless the wheel inertia is scaled consistently.
  \item[Diagnosis] Metamorphic relations need no reference solution, so they
    remain valid where no oracle exists.
\end{testcard}

\begin{testcard}{ORC-06}{Randomized property sweep with seeded, replayable streams}{property}{P1}
  \item[Targets] all layers
  \item[Regime] \texttt{F-RANDOM}.
  \item[Invariant] Every invariant marked \emph{property} in this suite holds
    for every draw.
  \item[Pass] $500$ draws from a single recorded seed sequence; every failure
    is shrunk to a minimal reproducer and appended to a replayed corpus that
    runs on every build.
  \item[Diagnosis] Without the recorded seed and the replay corpus, a
    randomized failure is unreproducible and will be dismissed as flaky.
\end{testcard}

\begin{testcard}{ORC-07}{Golden-file regression}{regression}{P2}
  \item[Targets] both QPs
  \item[Regime] Five canonical fixtures spanning both chassis and both
    terrains.
  \item[Invariant] Stored solutions reproduce within tolerance.
  \item[Pass] Gated on QP-03 uniqueness and on strict complementarity; where
    either fails, store and compare only the physically determined quantities,
    not $\generator{\force_i}$.
  \item[Diagnosis] An ungated golden file over a non-unique minimizer produces
    a test that fails on solver upgrades for no physical reason.
\end{testcard}

\begin{testcard}{ORC-08}{Closed-loop cross-check against an independent simulator}{smoke}{P2}
  \item[Targets] both QPs
  \item[Regime] Rigid-body simulator with the same contact model.
  \item[Invariant] Trajectories agree, with the discrepancy decreasing at the
    integrator's order as $\samplingPeriod$ shrinks.
  \item[Pass] The order-of-convergence assertion, not a fixed trajectory
    tolerance. Legitimate sources of disagreement---contact-model differences,
    solver tolerance, the polyhedral cone---must be enumerated and bounded in
    advance.
  \item[Diagnosis] A fixed tolerance here is unreliable; convergence order is
    the assertion that distinguishes a bug from a discretization difference.
\end{testcard}

\section{Layer H. Smoke tests}
\label{sec:test-layer-h}

\begin{testcard}{SMK-01}{One-cycle assembly contract}{smoke}{P0}
  \item[Targets] both QPs
  \item[Regime] Nominal fixture.
  \item[Invariant] One cycle assembles with the expected dimensions, no
    identically zero row, a positive-definite Hessian, and a solved status.
  \item[Pass] All four conditions; the zero-row check must exclude rows that
    are legitimately zero only for the current data.
  \item[Diagnosis] The cheapest test that can fail; run it first in the
    sequence.
\end{testcard}

\begin{testcard}{SMK-02}{Static hold on flat ground and on a ramp}{smoke}{P0}
  \item[Targets] \eqref{eq:test-ramp-split},
    \eqref{eq:planar-normal-load-row}
  \item[Regime] Zero command, \texttt{F-RAMP}.
  \item[Invariant] Nothing moves; the load model closes the three out-of-plane
    equations; the in-plane forces balance the gravity component.
  \item[Pass] Twist below $10^{-6}$\unitVelTS; force residuals per DYN-03 and
    DYN-07.
  \item[Diagnosis] Separates a controller fault from a load-model fault,
    because the two are asserted independently.
\end{testcard}

\begin{testcard}{SMK-03}{Straight-line drive}{smoke}{P0}
  \item[Targets] both QPs
  \item[Regime] Constant forward reference, flat ground.
  \item[Invariant] The chassis tracks the reference; for T1 the lateral
    residual decays at the rate of ROW-03, while for T2 it is only driven by
    the objective, since T2 carries no proportional term.
  \item[Pass] Steady-state tracking error below $1\%$; the two decay behaviours
    asserted separately.
  \item[Diagnosis] Expecting T2 to show T1's decay is a false-failure trap; the
    asymmetry is deliberate, see ROW-11.
\end{testcard}

\begin{testcard}{SMK-04}{In-place rotation}{smoke}{P0}
  \item[Targets] \eqref{eq:four-steering-wheel-qp},
    \eqref{eq:differential-drive-qp}
  \item[Regime] Zero linear reference, nonzero yaw reference.
  \item[Invariant] T1 produces equal and opposite wheel rates. T2 steers all
    wheels tangent to a circle about the chassis centre; read from inside the
    QP, every $\dot\delta_i$ must be zero once the configuration is reached,
    confirming the chassis-relative convention.
  \item[Pass] Yaw rate tracked to $1\%$; the T2 steering rates below
    $10^{-6}$~rad/s in the steady configuration.
  \item[Diagnosis] The steady-state steering rate is the cleanest end-to-end
    read-out of the GEO-07 convention.
\end{testcard}

\begin{testcard}{SMK-05}{Constant-radius arc and the centripetal term}{smoke}{P1}
  \item[Targets] \eqref{eq:reduced-dynamics-congruence},
    Remark~\ref{re:planar-basis-choice}
  \item[Regime] Steady arc, both chassis.
  \item[Invariant] All wheels share one ICR, and the lateral force required is
    the centripetal force $\mass v_x\omega_{\Pi}$, which must be supplied
    through $\bs{h}_{\text{red}}$ and the contact forces.
  \item[Pass] ICR consistency to $10^{-3}$ of the turn radius; lateral force
    within $10^{-3}\mass\gravity[value]$ of the analytic centripetal value.
  \item[Diagnosis] This is the end-to-end counterpart of DYN-02: an
    implementation missing $\dot{\text{P}}\bs{u}$ tracks a straight line and
    fails here.
\end{testcard}

\begin{testcard}{SMK-06}{Crab translation: the manoeuvre that separates T1 from T2}{smoke}{P0}
  \item[Targets] \eqref{eq:four-steering-wheel-qp},
    \eqref{eq:differential-drive-qp}
  \item[Regime] Pure lateral reference, all steering angles equal.
  \item[Invariant] T2 executes it. T1 cannot: its lateral row forbids lateral
    motion, so it must report the reference as untrackable rather than
    approximating it with a turn.
  \item[Pass] T2 tracks to $1\%$. T1 leaves a lateral tracking error equal to
    the full reference, and its solved lateral velocity remains below
    $10^{-6}$\unitVelTS.
  \item[Diagnosis] The single most discriminating scenario: it fails
    immediately if the two chassis share an assembler that ignores the
    differential-drive lateral constraint.
\end{testcard}

\begin{testcard}{SMK-07}{Incompatible four-wheel data surfaces as slack}{smoke}{P0}
  \item[Targets] \eqref{eq:four-steering-wheel-qp},
    Remark~\ref{re:four-steering-lateral-slack}
  \item[Regime] T2, deliberately non-concurrent steering angles.
  \item[Invariant] The QP solves, returns \emph{one} chassis twist, and reports
    the incompatibility through $\slackVariable^{\text{lat}}$ with magnitude
    bounded below per ROW-09. It must not return four independent wheel
    velocities.
  \item[Pass] Solved status; slack within $1\%$ of its predicted lower bound;
    exactly one twist in the solution.
  \item[Diagnosis] The end-to-end regression for the infeasibility defect.
\end{testcard}

\begin{testcard}{SMK-08}{Steering-rate-limited slalom}{smoke}{P1}
  \item[Targets] \eqref{eq:four-steering-wheel-qp} rate bounds
  \item[Regime] T2, reference that saturates the steering-rate bound.
  \item[Invariant] The predicted-rate bound activates and the tracking error
    grows in the predictable way, matching a rate-and-acceleration-limited
    reference model.
  \item[Pass] Rate never exceeds the bound; tracking error within $10\%$ of the
    limited-reference prediction.
  \item[Diagnosis] Confirms BND-01 and BND-02 end to end, including the
    chassis-relative convention under sustained yaw.
\end{testcard}

\begin{testcard}{SMK-09}{Slope sweep to friction and torque failure}{smoke}{P1}
  \item[Targets] \eqref{eq:test-ramp-split}, \texttt{F-RAMP}
  \item[Regime] Slopes $0$ to $40$~deg.
  \item[Invariant] Station keeping succeeds below the predicted band of FRI-08
    and fails above it; the friction margin decreases monotonically with slope.
  \item[Pass] The observed failure slope inside the band; monotone margin.
  \item[Diagnosis] The band must be computed from the load model, not assumed,
    or this test measures the load model and reports it as a controller fault.
\end{testcard}

\begin{testcard}{SMK-10}{The tangent basis rotates with the chassis}{smoke}{P0}
  \item[Targets] \eqref{eq:differential-pose-rate},
    \eqref{eq:chassis-aligned-tangent-basis}
  \item[Regime] Closed circular path returning to the start pose.
  \item[Invariant] Integrating \eqref{eq:differential-pose-rate} with
    $\text{E}_{\Pi}$ propagated as it rotates closes the circle. Freezing
    $\text{E}_{\Pi}$ at its initial value does not.
  \item[Pass] Closure error below $10^{-3}$ of the circle radius for the
    correct integration; the frozen-basis variant must fail by at least the
    radius itself.
  \item[Diagnosis] Directly targets the integration error flagged after
    \eqref{eq:differential-pose-rate}. The two-branch form guarantees the test
    discriminates.
\end{testcard}

\begin{testcard}{SMK-11}{Odometry cross-check, including where it must disagree}{smoke}{P1}
  \item[Targets] \eqref{eq:differential-pose-rate},
    \secRef{sec:rolling-implementation}
  \item[Regime] T1, flat no-slip ground; then a ramp.
  \item[Invariant] On flat ground, closed-form wheel odometry agrees with the
    integrated QP state. On a ramp, an odometry model that assumes a world
    horizontal plane must \emph{disagree}, and the test asserts the
    disagreement rather than tolerating it.
  \item[Pass] Flat: agreement to $10^{-3}$ of path length. Ramp: disagreement
    at least $\sin\theta$ times path length for the horizontal-plane model.
  \item[Diagnosis] Prevents an odometry model from being ``validated'' on flat
    ground and silently used on slopes.
\end{testcard}

\begin{testcard}{SMK-12}{Terrain transition}{smoke}{P1}
  \item[Targets] \eqref{eq:chassis-aligned-tangent-basis},
    \eqref{eq:rolling-force-generators}
  \item[Regime] Roll across a $0$ to $10$~deg slope discontinuity over three to
    five cycles.
  \item[Invariant] $\bs{\rho}_i$ and $\text{H}_i$ stay bounded and continuous
    in the chassis frame; the cone is rebuilt for the new normal; no solver
    failure occurs at the transition.
  \item[Pass] No discontinuity in the solved twist exceeding the commanded
    acceleration times $\samplingPeriod$; cone rebuild confirmed per FRI-03.
  \item[Diagnosis] Both terrains are otherwise exercised only as steady states;
    this is the only entry that tests the transition between them.
\end{testcard}

\begin{testcard}{SMK-13}{Long-horizon closed loop}{smoke}{P1}
  \item[Targets] both QPs
  \item[Regime] $60$~s of varied commands, both terrains.
  \item[Invariant] The rolling residual stays bounded with no monotone growth;
    T1 shows input-to-state stability from its proportional term, while T2's
    residual is bounded only by the objective and may drift more.
  \item[Pass] Residual bounded by a stated envelope; no monotone trend over the
    last $30$~s. Assert different envelopes for T1 and T2.
  \item[Diagnosis] Applying T1's envelope to T2 is a false-failure trap.
\end{testcard}

\section{Smoke sequence and continuous integration}
\label{sec:test-sequence}

The smoke tests are ordered cheapest and most decisive first, so that a broken
build fails in seconds rather than minutes:
\begin{enumerate}
  \item \texttt{SMK-01}, one-cycle assembly contract;
  \item \texttt{SMK-02}, static hold on flat ground and on a ramp;
  \item \texttt{SMK-03}, straight-line drive;
  \item \texttt{SMK-04}, in-place rotation;
  \item \texttt{SMK-06}, crab translation, which separates T1 from T2;
  \item \texttt{SMK-07}, incompatible four-wheel data surfacing as slack;
  \item \texttt{SMK-10}, closed-circle pose integration;
  \item \texttt{SMK-05}, constant-radius arc and the centripetal term;
  \item \texttt{SMK-09}, slope sweep to failure;
  \item \texttt{SMK-08}, steering-rate-limited slalom;
  \item \texttt{SMK-12}, terrain transition;
  \item \texttt{SMK-11}, odometry cross-check;
  \item \texttt{SMK-13}, long-horizon closed loop.
\end{enumerate}

Every smoke test records the metrics required in
\secRef{sec:rolling-cases}: $\norm{\rollingResidual_i}$,
$\norm{\slipVelocity{i}}$, $\norm{\slackVariable_i}$, normal force, friction
margin, wheel-torque margin, solver status, and solve time. Recording them
unconditionally, including on passing runs, is what makes a later regression
diagnosable.

Suggested gating: all \oBlock{P0} entries on every commit; \oBlock{P0} and
\oBlock{P1} on every merge; the full suite including \oBlock{P2} and the
randomized sweeps nightly.

\section{Traceability}
\label{sec:test-traceability}

\begin{table}[h]
  \centering
  \caption{Formulation element to test mapping}
  \label{tab:test-traceability}
  \begin{tabular}{|C{0.34\columnwidth}|C{0.54\columnwidth}|}
    \hline
    Element of the formulation & Entries \\
    \hline
    Chassis-aligned basis, \eqref{eq:chassis-aligned-tangent-basis} &
      GEO-01, GEO-02, GEO-03, ROW-04, SMK-10 \\
    \hline
    Carrier map, \eqref{eq:planar-carrier-map} &
      GEO-05, GEO-10, ROW-05 \\
    \hline
    Congruence reduction, \eqref{eq:reduced-dynamics-congruence} &
      DYN-01, DYN-02, GEO-08, ORC-02, SMK-05 \\
    \hline
    Normal annihilation, \eqref{eq:reduced-normal-annihilated} &
      GEO-06, GEO-08, FRI-05, FRI-07 \\
    \hline
    Normal-load row, \eqref{eq:planar-normal-load-row} &
      DYN-07, FRI-05, FRI-06, BND-07, SMK-02 \\
    \hline
    Force regularization, \eqref{eq:planar-force-regularization} &
      QP-01, QP-02, QP-03 \\
    \hline
    Twist weight, \eqref{eq:planar-twist-weight} &
      QP-06 \\
    \hline
    T1 kinematic rows &
      ROW-01, ROW-02, ROW-03, ROW-08, ORC-03 \\
    \hline
    T1 stabilization, $\pGain$ &
      ROW-03, ROW-11, QP-08, SMK-03 \\
    \hline
    T2 kinematic rows &
      ROW-05, ROW-06, ROW-10, ROW-12, ORC-04 \\
    \hline
    T2 lateral slacks &
      ROW-09, BND-06, QP-05, SMK-07 \\
    \hline
    Rate prediction, \eqref{eq:four-steering-rate-prediction} &
      ROW-07, BND-01, BND-03 \\
    \hline
    Steering convention &
      GEO-07, BND-02, SMK-04, ROW-06 \\
    \hline
    Friction cone, \eqref{eq:rolling-force-generators} &
      FRI-01, FRI-02, FRI-03, FRI-04, FRI-09, SMK-12 \\
    \hline
    Bounds and actuation &
      BND-01 to BND-08, SMK-08 \\
    \hline
    Assumptions \ref{as:planar-basis}--\ref{as:normal-load} &
      GEO-01, GEO-03, ROW-01, ROW-05, DYN-07, FRI-03 \\
    \hline
  \end{tabular}
\end{table}

\section{Residual gaps}
\label{sec:test-gaps}

The suite does not currently cover the following, and each is a deliberate
omission rather than an oversight:
\begin{itemize}
  \item \tBlock{Objective linear term.} The reference quantities
    $\bs{\xi}^{\Pi,\text{ref}}_{\fbFrame}$,
    $\reff{\ddot{\wheelAngle}}_i$, $\reff{\ddot\delta}_i$,
    $\reff{\dot{\wheelAngle}}_i$ and $\reff{\dot\delta}_i$ enter only through
    $\bs{g}$. No entry above asserts that $\bs{g}$ is assembled with the right
    sign and the right $\samplingPeriod$ factor. Add a unit test asserting
    $\bs{g}=-2\samplingPeriod\weightMatrix{W}_{\xi}
    (\bs{\xi}^{\Pi}_{\fbFrame}-\bs{\xi}^{\Pi,\text{ref}}_{\fbFrame})$ on the
    twist columns, under the factor convention pinned in
    \secRef{sec:test-layer-f}.
  \item \tBlock{Reference frame of $\bs{\xi}^{\Pi,\text{ref}}_{\fbFrame}$.}
    \eqref{eq:planar-twist-weight} is well posed only if the reference is
    expressed in the same chassis-aligned basis as the measurement. Nothing
    asserts this; a reference published in a world frame would rotate under
    GEO-01.
  \item \tBlock{Sliding contacts.} The suite assumes
    $\set{s}=\emptyset$, consistent with the ideal formulations. Testing
    \eqref{eq:sliding-friction-direction} beyond the negative test FRI-10
    requires the mode-aware QP \eqref{eq:rolling-whole-body-qp}.
  \item \tBlock{Wheel and steering-column inertia.} The reduced models carry
    them inside $\text{M}_{\text{red}}$, but no entry isolates them; a wrong
    gear ratio would show up only as a torque-magnitude error.
  \item \tBlock{Rolling resistance and self-aligning moment.} Outside the
    point-contact model of assumption~\ref{as:steering-axis}; these set the
    real steering torque and are not represented.
  \item \tBlock{Non-coplanar contacts.} Assumption~\ref{as:one-plane} is
    asserted, not relaxed. A leg-wheel machine on non-coplanar terrain needs
    the generic formulation, not these planar QPs.
  \item \tBlock{Timing under load.} QP-09 asserts a solve-time budget on an
    unloaded machine; contention and cache effects on the target are not
    modelled.
\end{itemize}

\bibliography{rolling-contact-qp}
\bibliographystyle{unsrtnat}

\end{document}
